\documentclass[11pt]{article}
\usepackage[T2A]{fontenc}
\usepackage[utf8]{inputenc}
\usepackage[russian]{babel}
\AtBeginDocument{\shorthandoff{"}}
\usepackage{lmodern}
\usepackage[margin=1in]{geometry}
\usepackage{microtype}

\emergencystretch=2em
\tolerance=1500
\hbadness=10000
\binoppenalty=700
\relpenalty=500

\usepackage{amsmath,amssymb,amsthm,mathtools,bm,mathrsfs}
\usepackage{longtable,booktabs,array,multirow}
\usepackage{textcomp}
\usepackage{graphicx}
\usepackage{enumitem}
\usepackage{xcolor}
\usepackage{framed}
\usepackage{fancyvrb}
\usepackage{ragged2e}
\usepackage{hyperref}
\usepackage{xurl}
\usepackage{tikz}
\usepackage{tikz-cd}
\usetikzlibrary{arrows.meta, positioning, calc, shapes.geometric, fit, decorations.pathreplacing}
\usepackage{caption}

% Theorem environments (Russian)
\theoremstyle{plain}
\newtheorem{theorem}{Теорема}[section]
\newtheorem{lemma}[theorem]{Лемма}
\newtheorem{proposition}[theorem]{Предложение}
\newtheorem{corollary}[theorem]{Следствие}

\theoremstyle{definition}
\newtheorem{definition}[theorem]{Определение}
\newtheorem{example}[theorem]{Пример}
\newtheorem{remark}[theorem]{Замечание}
\newtheorem{convention}[theorem]{Соглашение}

% Custom commands (same as English)
\newcommand{\RS}{\mathrm{R\text{-}S}}
\newcommand{\cS}{\mathcal{S}}
\newcommand{\cF}{\mathcal{F}}
\newcommand{\cM}{\mathcal{M}}
\newcommand{\cC}{\mathcal{C}}
\newcommand{\cD}{\mathcal{D}}
\newcommand{\cL}{\mathcal{L}}
\newcommand{\cP}{\mathcal{P}}
\newcommand{\fM}{\mathfrak{M}}
\newcommand{\Meta}{\mathfrak{Meta}}
\newcommand{\LFnd}{\mathcal{L}_{\mathrm{Fnd}}}
\newcommand{\LCls}{\mathcal{L}_{\mathrm{Cls}}}
\newcommand{\LClsMax}{\mathcal{L}_{\mathrm{Cls}}^{\top}}
\newcommand{\LAbs}{\mathcal{L}_{\mathrm{Abs}}}
\newcommand{\Sint}{\cS_{\mathrm{int}}}
\newcommand{\Ob}{\mathrm{Ob}}
\newcommand{\Mor}{\mathrm{Mor}}
\newcommand{\Hom}{\mathrm{Hom}}
\newcommand{\Fun}{\mathrm{Fun}}
\newcommand{\Cl}{\mathrm{Cl}}
\newcommand{\Lan}{\mathrm{Lan}}
\newcommand{\Ran}{\mathrm{Ran}}
\newcommand{\Set}{\mathrm{Set}}
\newcommand{\Cat}{\mathrm{Cat}}
\newcommand{\Syn}{\mathrm{Syn}}
\newcommand{\Mod}{\mathrm{Mod}}
\newcommand{\Con}{\mathrm{Con}}
\newcommand{\PA}{\mathsf{PA}}
\newcommand{\ZFC}{\mathsf{ZFC}}
\newcommand{\MLTT}{\mathsf{MLTT}}
\newcommand{\HoTT}{\mathsf{HoTT}}
\newcommand{\ETT}{\mathsf{ETT}}
\newcommand{\CIC}{\mathsf{CIC}}
\newcommand{\NBG}{\mathsf{NBG}}
\newcommand{\AFA}{\mathsf{AFA}}
\newcommand{\CZF}{\mathsf{CZF}}
\newcommand{\II}{\mathbf{I}}
\newcommand{\cU}{\mathcal{U}}
\newcommand{\tpi}{\widetilde{\pi}}
\newcommand{\cE}{\mathcal{E}}
\newcommand{\fME}{\fM_{\cE}}
\newcommand{\LAbsE}{\LAbs^{\cE}}
\newcommand{\LClsE}{\LCls^{\cE}}
\newcommand{\LClsMaxE}{(\LClsMax)^{\cE}}
\newcommand{\LFndE}{\LFnd^{\cE}}
\newcommand{\SSE}{\cS_{S}^{\cE}}
\newcommand{\Eff}{\ensuremath{\mathsf{Eff}}}
\newcommand{\NCG}{\ensuremath{\mathsf{NCG}}}
\DeclareMathOperator{\colim}{colim}
\DeclareMathOperator{\image}{image}
\DeclareMathOperator{\id}{id}
\DeclareMathOperator{\ev}{ev}
\DeclareMathOperator{\Trace}{Trace}
\DeclareMathOperator{\dep}{depth}

\newcommand{\gr}{\operatorname{gr}}
\newcommand{\OWL}{\ensuremath{\mathsf{OWL}}}
\newcommand{\SSEbase}{\cS_{S}^{\cE,\,\mathrm{base}}}
\newcommand{\SSEglobal}{\cS_{S}^{\cE,\,\mathrm{global}}}

\title{Пространство модулей формальных систем:\\
Классификация, стабилизация и no-go-теорема для абсолютных оснований}

\author{%
  \footnotesize Макс Середа \\[-3pt]
  \scriptsize\textit{Независимый исследователь} \\[-3pt]
  \scriptsize\texttt{\href{mailto:max@gst.st}{max@gst.st}}%
}

\date{\scriptsize\today}

\begin{document}
\maketitle

\begin{abstract}
Я изучаю $(\infty, n)$-категорное пространство модулей $\fM$ формальных систем (далее \textbf{MSFS}): классифицирующий $2$-стек Морита-классов эквивалентности Rich-оснований ($\LFnd$ — теории, удовлетворяющие условиям (R1)--(R5)). Я устанавливаю пять структурных результатов и одно граничное следствие.

\emph{(i) Плюрализм в $\LCls$.} $\infty$-космои (Риль--Верити), Унивалентные основания (Воеводский) и когезивные высшие топосы (Шрайбер) — попарно $2$-неэквивалентные частичные мета-фреймворки, каждый классифицирующий собственный строгий под-стек $\fM$; они лежат вне жёсткого максимального подкласса $\Meta_{\mathrm{Cls}}^{\top}$.

\emph{(ii) Условная мета-категоричность.} Любые два члена $\Meta_{\mathrm{Cls}}^{\top}$ над одной и той же Rich-метатеорией являются $(\infty, \infty)$-эквивалентными согласно Гротендик--Люри-выпрямке с совместной точностью экстенсионального и интенсионального классификационных функторов.

\emph{(iii) Срез-локальное интенсиональное уточнение.} Функтор отображённого $2$-класса $\II : \cF^{\mathrm{op}} \to \Sint$ срез-локален над $\fM$: интенсиональные различия ($\MLTT$ против $\ETT$) расслаиваются над одиночными точками $\fM$, инвариантно разделяясь внутренне-языковой вычислимостью в эффективном топосе $\mathrm{Eff}$ Хайланда.

\emph{(iv) Мета-стабилизация на теоретическом уровне.} Итерированная мета-классификация воспроизводит ту же $(\infty, \infty)$-теорию на каждом шаге (унитарность Барвика--Шоммер-Приса), хотя теоретико-множественная инстанциация поднимается по иерархии универсумов Гротендика $\kappa_1 < \kappa_2 < \cdots$.

\emph{(v) AC/OC-двойственность.} Сопряжение Ламбека--Скотта $(\Syn, \Mod)$ из (R5b), собранное равномерно по $\RS$ и над $n \in \mathbb{N} \cup \{\infty\}$, индуцирует $(\infty, \infty)$-эквивалентность $\fM \simeq \fM_\cE$ между теоретико-центричным $2$-стеком модулей $\fM$ и актом-центричным $2$-стеком модулей $\fM_\cE$ заострённых Rich-актов. Дуал AFN-T утверждает, что акт-абсолютный стратум также пуст.

\emph{Граничное следствие: No-go-теорема для абсолютных оснований (AFN-T).} По синтаксис-семантическому сопряжению, лежащему в основе (R5), $\fM$ не имеет максимальной точки: страта $\LAbs$ (формально определимое, нередуцируемое и максимально порождающее основание) пуста. AFN-T обобщает классическую no-go-серию Кантор~\cite{Cantor1891}, Рассел~\cite{Russell1903}, Гёдель~\cite{Godel1931}, Тарский~\cite{Tarski1936}, Ловер~\cite{Lawvere1969}, Эрнст~\cite{Ernst2015} равномерно по $\RS$ и всем $(\infty, n)$, $n \in \mathbb{N} \cup \{\infty\}$, и расширяется двойственно на акт-абсолютную страту $\LAbs^\cE$.

\smallskip
\noindent\textbf{Ключевые слова:} пространство модулей формальных систем, классифицирующие $2$-стеки, категорная логика, $(\infty, n)$-категории, теория высших топосов, интенсиональная теория типов, мета-классификация, эффективный топос, no-go-теоремы.

\smallskip
\noindent\textbf{MSC2020:} 03A05, 03B30, 03C45, 03F40, 18A05, 18N10, 18N50, 18N60, 18N65.

\end{abstract}

\newpage
\tableofcontents
\newpage

\section{Введение}

Эта статья изучает \emph{пространство модулей формальных систем} $\fM$ — классифицирующий $(\infty, 2)$-стек, точки которого — Морита-классы эквивалентности Rich-оснований (формальных теорий $F$, удовлетворяющих условиям (R1)--(R5) Определения~\ref{def:rs}), а $k$-морфизмы — точные интерпретации и доказуемые естественные эквивалентности. Главные результаты касаются структуры $\fM$:
\begin{itemize}
\item \textbf{Стратификация} (\S\ref{sec:hierarchy}): $\LFnd \xrightarrow{\mathrm{Cls}} \LCls \supsetneq \LClsMax \xrightarrow{\mathrm{Gen}} \LAbs = \emptyset$ --- страты высекаются разнотипными пакетами условий ((R1)--(R5) на формальных системах против (M1)--(M5) на классификаторах, Предложение~\ref{prop:level-structure}~(ii)); единственное включение --- $\LClsMax \subsetneq \LCls$.
\item \textbf{Структурный плюрализм в $\LCls$} (Теорема~\ref{thm:meta-mult}).
\item \textbf{Условная категоричность $\LClsMax$} (Теорема~\ref{thm:meta-cat}).
\item \textbf{Стабилизация на теоретическом уровне с восхождением по универсумам} (Теорема~\ref{thm:meta-stab}).
\item \textbf{Срез-локальность интенсионального уточнения} (Теорема~\ref{thm:slice-locality}).
\item \textbf{Интенсиональная градуировка слоёв среза} (Определение~\ref{def:display-grade}, Предложение~\ref{prop:grading}): каждый слой $\mathrm{Int}([F])$ несёт каноническую функториальную $\mathbb{N}$-фильтрацию по display-грейду; экстенсиональность — в точности слепота к грейду (Следствие~\ref{cor:grade-collapse}).
\item \textbf{Пятиосевая абсолютность} (\S\ref{sec:five-axis}).
\item \textbf{AC/OC Морита-двойственность} (Теорема~\ref{thm:ac-oc-duality}, Теорема~\ref{thm:dual-afnt}).
\item \textbf{Внешняя граница: $\LAbs = \emptyset$} --- No-go-теорема для абсолютных оснований (Теорема~\ref{thm:afnt-alpha}, Теорема~\ref{thm:afnt-beta}).
\end{itemize}

\paragraph{Машинная верификация.}  Каждая теорема, лемма и следствие этой статьи имеет сопутствующее машинно-проверенное формальное доказательство в \textbf{Verum MSFS Corpus} (\url{https://github.com/gst-st/msfs}), самодостаточном по модулю ZFC + 2 сильно недостижимых кардинала.  Точка входа в один абзац — Приложение~\ref{app:verification-ledger}; живая аудит-поверхность, карта «статья $\to$ корпус» и инструкции воспроизводимости — в репозитории корпуса.

\subsection{Обозначения и соглашения}

\begin{convention}\label{conv:zfc-inacc}
Я работаю всюду в $\ZFC$, дополненной двумя недостижимыми кардиналами $\kappa_1 < \kappa_2$, дающими универсумы Гротендика $\mathbf{U}_1 \subset \mathbf{U}_2$, достаточные для применяемой $2$-категорной техники. Категорный аппарат следует Келли~\cite{Kelly1982} для обогащённых и $2$-категорий, Люри~\cite{LurieHTT} для $(\infty, 1)$-категорий, Люри~\cite{LurieHA} для высшей алгебры и Риль--Верити~\cite{RiehlVerity2022} для $\infty$-космоев. Для теории доступных функторов следую Адамеку--Росицкому~\cite{AdamekRosicky1994}.
\end{convention}

\begin{convention}[Ключевые символы]\label{conv:notation}
\begin{itemize}
\item $\cF$ обозначает $2$-категорию \emph{Rich-оснований} ($\LFnd$): объекты — основополагающие теории в смысле, формализованном в \S\ref{sec:rs} ниже (условия (R1)--(R5)); $1$-морфизмы — точные интерпретации; $2$-морфизмы — доказуемые естественные эквивалентности между интерпретациями.
\item $\rho : \cF \to (\infty, n)\text{-}\Cat$ обозначает \emph{классификационный функтор}, сопоставляющий каждому $F \in \cF$ его категорию моделей (или, для теоретико-типовых $F$, синтаксическую $(\infty, n)$-категорию $\Syn(F)$).
\item \emph{Калибровочное преобразование} $\tau : F_1 \to F_2$ — это $1$-морфизм в $\cF$, сопровождаемый выделенной $2$-эквивалентностью $\rho(F_1) \simeq \rho(F_2)$ на модельно-категорном уровне. Два основания \emph{калибровочно эквивалентны}, если между ними существует обратимое калибровочное преобразование. \emph{Связь с Морита-эквивалентностью.} Калибровочная эквивалентность на уровне $\cF$ соответствует, через фактор $\fM := \cF/\text{калибровка}$, Морита-эквивалентности на уровне стека: $F_1, F_2 \in \cF$ калибровочно эквивалентны тогда и только тогда, когда $[F_1] = [F_2]$ в $\fM$, что по построению есть Морита-класс эквивалентности. Всюду в статье «калибровка» (уровень $1$-морфизмов) и «Морита» (уровень стека) именуют одно и то же лежащее в основе отношение эквивалентности на своих категорных стратах, связанные фактор-функтором $\cF \to \fM$; никакого третьего отношения не вводится.
\item \emph{Морита-редуцируемость на $(\infty, n)$-уровне}: кандидат $X$ в $\LAbs$ и любой $Y \in \cS_S$ рассматриваются как заострённые $(\infty, n)$-категорные структуры через каноническое вложение Йонеды в тотальную $(\infty, n)$-категорию $\mathbf{StrCat}_{S, n}$ $S$-структурированных $(\infty, n)$-категорий (Определение~\ref{def:strcat}, сформулированное сразу ниже). Конкретно: $X$, заданный формулой $\phi_X$, соответствует паре $(\Syn(S), [\phi_X])$, где $[\phi_X]$ — класс изоморфизма объекта, определённого $\phi_X$ в синтаксической категории; $Y \in \Ob(\cM) \subseteq \cS_S^{\mathrm{local}}$ соответствует $(\cM, Y)$. Я пишу $X \to_{M, n} Y$ и говорю, что $X$ \emph{Морита-редуцируем на $(\infty, n)$-уровне} к $Y$, если существует $(\infty, n)$-функтор $\Phi : \Syn(S) \to \cM$ в $\mathbf{StrCat}_{S, n}$, переводящий $[\phi_X]$ в класс изоморфизма $Y$, причём ограничение $\Phi |_{[\phi_X]}$ — $(\infty, n)$-эквивалентность на свой образ. Это расширяет классическую кольцевую Морита-эквивалентность (сравнивающую $\Mod(X)$ и $\Mod(Y)$ как $1$-категории), вбирая полную $(\infty, n)$-категорную структуру; оба понятия совпадают, когда $X, Y$ $1$-категорны и $n = 1$. Я использую сокращение «Морита-эквивалентны» для симметричного понятия и «Морита-редуцируем» для направленной версии.
\item $\fM$ обозначает \emph{классифицирующий $2$-стек} над $\cF$: фактор $\fM := \cF / \text{калибровка}$, рассматриваемый как $(\infty, 2)$-стек Морита-классов эквивалентности Rich-оснований. Формально $\fM$ — выпрямление картезианова расслоения $\int \rho \to \mathrm{R\text{-}S}$, слой которого над $S$ — $2$-группоид $\rho^{-1}(S)$ оснований над $S$ по модулю калибровки. Конструкция следует Люри HTT~\cite{LurieHTT} \S 3.2 (выпрямление--развыпрямление картезиановых расслоений).
\item Для $2$-категории $\mathbf{F}$ обозначение $\mathrm{Aut}_2(\mathbf{F})$ — $2$-группа $2$-автоэквивалентностей $\mathbf{F} \xrightarrow{\simeq_2} \mathbf{F}$ (моноидальная по композиции, с обратимыми $1$-морфизмами и когерентными $2$-клетками).
\item $\Sint$ обозначает $2$-категорию display-$2$-категорий (строгое определение отнесено в Приложение~\ref{app:categorical} ниже).
\item Все $2$-категории локально малы и $\kappa_1$-доступны, если иное не оговорено. (Для конкретных Rich-оснований $S$, чья сила непротиворечивости требует второго недостижимого, параметр доступности $\Syn(S)$ обобщается до $\kappa_S \in \{\kappa_1, \kappa_2\}$ согласно клаузе доступности (R5a) определения Rich-метатеории, формализуемого в \S\ref{sec:rs} ниже; последующие $\kappa_1$-утверждения следует читать равномерно как $\kappa_S$-утверждения, где $S$ индексирован конкретной рассматриваемой $\LFnd$-точкой. Все доказательства статьи устойчивы к этому равномерному чтению, поскольку аргументы доступности используют лишь регулярность кардинала, а не его конкретное значение.)
\item \emph{Усечение и поднятие по глубинам}: для $n \leq n_S$ я пишу $\Syn^{(\infty, n)}(S) := \tau_{\leq n}(\Syn(S))$ для $(\infty, n)$-усечения синтаксической категории. Для $n > n_S$ я пишу $\Syn^{(\infty, n)}(S)$ для канонического $(\infty, n)$-поднятия, получаемого левым Кан-расширением $\Syn(S)$ вдоль вложения $(\infty, n_S)\text{-}\Cat \hookrightarrow (\infty, n)\text{-}\Cat$ (Люри HTT~\cite{LurieHTT} \S 5.4). При этом соглашении сравнения объекта $X$ на уровне $n$ с $S$-определимым объектом $Y$ на родной глубине $n_S$ происходят на уровне $\max(n, n_S)$ через канонические поднятия; при $n = \infty$ объемлющим служит копредел $\colim_{k < \infty} \Syn^{(\infty, k)}(S)$.
\end{itemize}
\end{convention}

\begin{definition}[Тотальная $(\infty, n)$-категория $\mathbf{StrCat}_{S, n}$]\label{def:strcat}
Для $S \in \RS$ и $n \in \mathbb{N} \cup \{\infty\}$ \emph{тотальная $(\infty, n)$-категория $S$-структурированных $(\infty, n)$-категорий}, обозначаемая $\mathbf{StrCat}_{S, n}$, имеет:
\begin{itemize}
\item \textbf{Объекты}: пары $(\cC, \mathrm{str}_\cC)$, где $\cC$ — $(\infty, n)$-категория, либо существенно $\mathbf{U}_1$-малая, либо $\kappa_1$-доступная (два случая покрывают синтаксические категории и категории моделей соответственно; ср.\ (R5a)), а $\mathrm{str}_\cC$ — $S$-модельная структура на $\cC$ (то есть интерпретация $S$-сигнатуры в $\cC$, удовлетворяющая аксиомам $S$).
\item \textbf{$1$-морфизмы}: $S$-сохраняющие $(\infty, n)$-функторы $\Phi : (\cC_1, \mathrm{str}_1) \to (\cC_2, \mathrm{str}_2)$, то есть функторы, коммутирующие с соответствующими $S$-модельными структурами с точностью до когерентного $(\infty, n)$-изоморфизма.
\item \textbf{Высшие морфизмы}: когерентные $(\infty, n)$-естественные преобразования между $S$-сохраняющими функторами.
\end{itemize}
Объекты $\Syn(S)$ и любой модели $\cM \models S$ я вкладываю в $\mathbf{StrCat}_{S, n}$ конструкцией пойнтинга $X \mapsto (\Syn(S), X)$ и $Y \mapsto (\cM, Y)$ соответственно. Морита-редуцируемость, как указано в Соглашении~\ref{conv:notation}, означает существование $S$-сохраняющего $(\infty, n)$-функтора в $\mathbf{StrCat}_{S, n}$ между такими заострёнными структурами, являющегося эквивалентностью на свой образ.
\end{definition}

\subsection{Реестр допущений}\label{sec:assumption-register}

Каждая теорема ниже доказана \emph{из} конкретного пакета допущений, и заглавные утверждения статьи («$\LAbs = \emptyset$», «$\Meta_{\mathrm{Cls}}^{\top}$ категоричен») ровно настолько сильны, насколько силён этот пакет. Поскольку последующие работы строятся на MSFS как на базе --- в частности, любое расширение, позиционирующее себя как \emph{предельная} формальная теория, обязано либо дискаржить, либо явно унаследовать каждый пункт, --- я собираю допущения здесь в единый реестр, каждое со своим формальным статусом. В этом подразделе нет новой математики; это доверенная поверхность статьи, сделанная перечислимой.

\begin{description}
\item[\textbf{(A1) Основание.}] $\ZFC$ + два сильно недостижимых кардинала $\kappa_1 < \kappa_2$ (Соглашение~\ref{conv:zfc-inacc}). \emph{Статус: основной аксиомный пакет.} Все результаты — теоремы этой метатеории; по Гёделю~II пакет не может удостоверить собственную непротиворечивость, и статья нигде не утверждает обратного (ср.\ открытый вопрос (Q4), \S\ref{sec:consequences}). Утверждения о силе непротиворечивости ниже по течению могут быть в лучшем случае равносогласованностями с (A1).
\item[\textbf{(A2) Скоуп «метатеории».}] No-go квантифицирует по $\RS$, очерченному условиями (R1)--(R5) (\S\ref{sec:rs}): рекурсивно перечислимая, интерпретирующая арифметику Робинсона $Q$, обладающая моделью в $\mathbf{U}_2$, $\mathbf{StrCat}$-совместимая, с существенно $\mathbf{U}_1$-малой $\Syn(S)$ с $\kappa_S$-доступным $\mathrm{Ind}$-пополнением (R5a). \emph{Статус: определенческий скоуп.} Системы вне $\RS$ --- не-р.п.\ оракульные теории, инфинитарные логики без доступного синтаксиса --- не подпадают под вердикт. Теорема~\ref{thm:afnt-alpha} универсальна над $\RS$, а не над неформальной тотальностью «всех оснований».
\item[\textbf{(A3) Формальное содержание «абсолютного основания».}] Членство в $\LAbs$ \emph{определено} как конъюнкция $(F_S) \wedge (\Pi_{4, S, n}) \wedge (\Pi_{3\text{-max}, S, n})$ (Определение~\ref{def:level6}). \emph{Статус: допущение интерпретативной адекватности.} AFN-T закрывает ровно эту тройку; философское чтение «абсолютного основания не существует» верно ровно в той мере, в какой тройка захватывает подразумеваемое понятие. Пятиосевой анализ \S\ref{sec:five-axis} укрепляет это чтение ось за осью, но не устраняет его интерпретативного характера.
\item[\textbf{(A4) Отношение сравнения.}] Объекты сравниваются через Морита-редуцируемость: существование $S$-сохраняющего $(\infty, n)$-функтора в $\mathbf{StrCat}_{S, n}$, являющегося эквивалентностью на свой образ, с объектами, вложенными конструкцией пойнтинга (Определение~\ref{def:strcat}). \emph{Статус: структурная конвенция.} Более тонкая или более грубая топология сравнения перечертила бы границы стратов; \S\ref{sec:intensional-grading} измеряет в точности то, что забывает Морита-фактор.
\item[\textbf{(A5) Исчерпываемость осей.}] Пространством параметров побега объявлены пять осей $(S, n, \mu, \xi, \pi)$ из \S\ref{sec:five-axis}. \emph{Статус: определенческий скоуп, не теорема.} Исчерпываемость выполнена по определению того, что статья считает параметром побега; подлинно новая ось потребовала бы новой теоремы, а не противоречила бы существующей.
\item[\textbf{(A6) Аксиомы мета-классификатора.}] Классифицирующий аппарат удовлетворяет (M1)--(M5), а для максимального класса — (Max-1)--(Max-4), включая клаузу совместной консервативности (Max-4) для пары $(\Cl_{\mathbf{F}}, \mathbb{I}_{\mathbf{F}})$, питающую шаг точности Теоремы~\ref{thm:meta-cat}. \emph{Статус: аксиомы на классификатор, мотивированные, но не выведенные.}
\item[\textbf{(A7) Локализационная семантика эквивалентности.}] AC/OC-двойственность (\S\ref{sec:ac-oc-duality}) читает «эквивалентность» через $\mathrm{gauge}_\cE$ — $2$-локализацию $\cE$ по классу рефлексий $\mathcal{R}_\cE$ (исчисление дробей Пронк, (BF1)--(BF5)): единица рефлексии обращается \emph{локализацией}, а не в самой $\cE$. \emph{Статус: семантическая конвенция} того же рода, что кольцевая Морита-эквивалентность.
\item[\textbf{(A8) Минимальный display-класс.}] Интенсиональная структура измеряется относительно \emph{канонически минимального} display-класса $\mathcal{D}_{\mathbf{F}}$, индуктивно порождённого (D1)--(D4); display-грейд \S\ref{sec:intensional-grading} относителен к этому порождению. \emph{Статус: конвенция канонической минимальности;} больший display-класс огрубил бы градуировку, но не потревожил бы экстенсиональный фактор.
\end{description}

\begin{remark}[Интерфейс для предельных расширений]\label{rem:register-interface}
Реестр спроектирован как формальный интерфейс. Расширение MSFS, претендующее на любую форму \emph{предельного} статуса, может быть аудировано пункт за пунктом: (A1) можно лишь унаследовать (равносогласованность — потолок); (A2) и (A5) можно \emph{интернализовать} --- доказать эквивалентными структурным компромиссам внутри расширения, --- но не обойти; (A3), (A4), (A7), (A8) — конвенции, которые расширение обязано либо принять дословно, либо пере-обосновать; (A6) можно усилить. В частности, предельное расширение консистентно с AFN-T ровно тогда, когда его предельное притязание индексировано этим реестром --- максимальность \emph{внутри} стратов типа $\LClsMax$, --- а не как членство в $\LAbs$, которую Теорема~\ref{thm:afnt-alpha} оставляет пустой.
\end{remark}

\section{Стратифицированная иерархия пространства модулей}\label{sec:hierarchy}

Я стратифицирую $\fM$ на четыре формальные страты, различаемые двумя структурно различными мета-операциями $\mathrm{Cls}$ (горизонтальная, классифицирующая) и $\mathrm{Gen}$ (вертикальная, порождающая). Страты, определяемые своими категорными условиями, таковы:
\[
\LFnd \;\xrightarrow{\mathrm{Cls}}\; \LCls \;\supsetneq\; \LClsMax \;\xrightarrow{\mathrm{Gen}}\; \LAbs.
\]

\begin{figure}[ht]
\centering
\begin{tikzpicture}[
  level/.style={draw, rounded corners=2pt, minimum width=5cm, minimum height=0.8cm, align=center, font=\small},
  bluefill/.style={fill=blue!6},
  violetfill/.style={fill=violet!10},
  emptybox/.style={draw, dashed, rounded corners=2pt, minimum width=5cm, minimum height=0.8cm, align=center, font=\small, fill=red!5},
  novelty/.style={-{Stealth[length=2mm]}, thick},
  horizmeta/.style={-{Stealth[length=2mm]}, thick, blue!60!black, densely dashed},
  vertmeta/.style={-{Stealth[length=2mm]}, thick, red, dashed},
  node distance=0.45cm
]
\node[level, violetfill] (L5) {$\LFnd$ --- Rich-основания (R1)--(R5)};
\node[level, violetfill, above=of L5] (L5plus) {$\LCls$ --- классификаторы (M1)--(M5)};
\node[level, violetfill, above=of L5plus] (L5max) {$\LClsMax$ --- максимальные (Max-1)--(Max-4)};
\node[emptybox, above=of L5max] (L6) {$\LAbs$ --- \emph{пуста по AFN-T}};
\draw[horizmeta] (L5.north) -- (L5plus.south) node[midway, right, font=\scriptsize, blue!60!black] {$\,\mathrm{Cls}$};
\draw[novelty] (L5plus.north) -- (L5max.south) node[midway, right, font=\scriptsize] {$\,\mathrm{max}$};
\draw[vertmeta] (L5max.north) -- (L6.south) node[midway, right, font=\scriptsize, red] {$\,\mathrm{Gen}$};
\draw[red, very thick] ($(L6.north west)+(0.4,-0.2)$) -- ($(L6.south east)+(-0.4,0.2)$);
\draw[red, very thick] ($(L6.north east)+(-0.4,-0.2)$) -- ($(L6.south west)+(0.4,0.2)$);
\end{tikzpicture}
\caption{Формальная стратификация $\fM$. Синяя пунктирная стрелка $\mathrm{Cls}$ — \emph{горизонтальная мета-операция} (классификация без порождения; см.\ \S\ref{sec:indexing-scheme}), переводящая $\LFnd$ в $\LCls$. Ограничение максимальности выделяет подкласс $\LClsMax$. Красная пунктирная стрелка $\mathrm{Gen}$ — \emph{вертикальная мета-операция} (максимальная порождающая способность), образ которой $\LAbs$ пуст по AFN-T (Теорема~\ref{thm:afnt-alpha}). $\LFnd$ включает $\ZFC$, $\MLTT$, $\HoTT$, теорию $(\infty, 1)$-топосов и некоммутативную геометрию. $\LCls$ распадается на максимальный подкласс $\Meta_{\mathrm{Cls}}^{\top}$ (категоричный по Теореме~\ref{thm:meta-cat}) и плюральное дополнение \emph{частичных} классификаторов (Теорема~\ref{thm:meta-mult}; $\infty$-космои, Унивалентные основания, когезивные $(\infty, 1)$-топосы), каждый из которых классифицирует строгий под-стек $\fM$. По Теореме~\ref{thm:ac-oc-duality} диаграмма дословно применима к акт-центричному стеку модулей $\fM_\cE$ через $(\infty, \infty)$-эквивалентность $\fM \simeq \fM_\cE$: дуальные страты $\LFndE \to \LClsE \supsetneq \LClsMaxE \to \LAbsE$ зеркалят страты выше, и $\LAbsE$ также пуста (Теорема~\ref{thm:dual-afnt}).}
\label{fig:hierarchy}
\end{figure}

\subsection{Формальные страты}

\begin{definition}[Стратифицированная иерархия --- формально]\label{def:hierarchy}
Я определяю формальные \emph{страты} как классы математических объектов:

\begin{tabular}{@{}ll@{}}
\toprule
Класс & Критерий членства \\ \midrule
$\LFnd$ & Формальные системы, удовлетворяющие (R1)--(R5) (уточняются в \S\ref{sec:rs} ниже) \\
$\LCls$ & Классификаторы, удовлетворяющие (M1)--(M5) (уточняются в \S\ref{sec:meta} ниже) \\
$\LClsMax$ & Классификаторы, дополнительно удовлетворяющие (Max-1)--(Max-4) (\S\ref{sec:meta} ниже) \\
$\LAbs$ & Объекты, удовлетворяющие $(F_S) \wedge (\Pi_{4, S, n}) \wedge (\Pi_{3\text{-max}, S, n})$ для некоторого $(S, n)$ (\S\ref{sec:conditions} ниже) \\
\bottomrule
\end{tabular}
\end{definition}

\begin{proposition}[Структурные свойства иерархии]\label{prop:level-structure}
\begin{enumerate}
\item[\textbf{(i)}] \emph{Определимость.} Каждая из четырёх страт — (возможно, собственный) класс, определимый в $\ZFC + 2\text{-inacc}$, дополненной $2$-категорной техникой Соглашения~\ref{conv:zfc-inacc}. Для $\LFnd$ условия — (R1)--(R5) (\S\ref{sec:rs} ниже); для $\LCls$ — (M1)--(M5) (\S\ref{sec:meta} ниже); для $\LClsMax$ они усиливаются до (Max-1)--(Max-4) (\S\ref{sec:meta} ниже); $\LAbs$ определяется тройкой $(F_S) \wedge (\Pi_{4, S, n}) \wedge (\Pi_{3\text{-max}, S, n})$ для некоторого $(S, n)$ (\S\ref{sec:conditions} ниже; структурная самопротиворечивость этой тройки, зафиксированная в \S\ref{sec:afnt-alpha}, даёт Теорему~\ref{thm:afnt-alpha}).
\item[\textbf{(ii)}] \emph{Многостратное членство.} $\LFnd$ и $\LCls$ характеризуются логически разными наборами условий (первопорядковые $(R1)$--$(R5)$ на формальную систему против $2$-категорных $(M1)$--$(M5)$ на $2$-категорию). Математический объект, несущий обе структуры, удовлетворяет обоим наборам одновременно: например, Унивалентные основания принадлежат $\LCls$, а их подлежащая $\HoTT$ — $\LFnd$. Я пишу $\mathrm{Stratum}(X) \subseteq \{\LFnd, \LCls, \LClsMax, \LAbs\}$ для множества формальных стратов, содержащих $X$.
\item[\textbf{(iii)}] \emph{Строгое включение на максимальном слое.} $\LClsMax \subsetneq \LCls$. Теорема~\ref{thm:meta-cat} утверждает, что любые два члена $\LClsMax$, параметризованные одной Rich-метатеорией, $(\infty, \infty)$-эквивалентны. Теорема~\ref{thm:meta-mult} даёт $\geq 3$ попарно не-$2$-эквивалентных членов дополнения $\LCls \setminus \LClsMax$: $\infty$-космои Риля--Верити, программу Унивалентных оснований и когезивную рамку Шрайбера.
\item[\textbf{(iv)}] \emph{Пустота абсолютной страты.} $\LAbs = \emptyset$ по Теореме~\ref{thm:afnt-alpha}. Следовательно, $\LFnd \cap \LAbs = \LCls \cap \LAbs = \LClsMax \cap \LAbs = \emptyset$ тривиально.
\end{enumerate}
\end{proposition}

\begin{proof}
\textbf{(i)} Каждый критерий выразим первопорядково; указанные классы определимы.

\textbf{(ii)} Наборы условий касаются непересекающихся аспектов объекта: аксиомы с исчислением выводов против $2$-категорной структуры с классифицирующим функтором. Объект, несущий обе структуры, удовлетворяет обоим критериям одновременно.

\textbf{(iii)} Включение — из двух наборов условий, уточняемых в \S\ref{sec:meta} ниже: (Max-1)--(Max-4) усиливают (M1)--(M5). Строгость — из теоремы структурной множественности (Теорема~\ref{thm:meta-mult}, доказываемая в \S\ref{sec:meta} ниже): $\mathbf{F}_{\mathrm{cosmoi}}, \mathbf{F}_{\mathrm{univalent}}, \mathbf{F}_{\mathrm{cohesive}}$ лежат в $\LCls \setminus \LClsMax$, поскольку образы их классификаций — строгие под-стеки $\fM$, что нарушает (Max-1).

\textbf{(iv)} Немедленно из Теоремы~\ref{thm:afnt-alpha}.
\end{proof}

\subsection{Примеры}

\paragraph{Страта $\LFnd$.} Основания, с характерными чертами:

\begin{longtable}{@{}p{3.8cm}l p{4cm} p{4.6cm}@{}}
\toprule
Основание & Год & Автор & Характеристика \\ \midrule
Z (Цермело) & 1908 & Цермело & Без схемы подстановки \\
ZF/ZFC & 1922 & Цермело, Френкель & Классическая теория множеств \\
ZFC + недостижимые & 1963+ & Гротендик & С универсумами \\
NBG & 1925--40 & фон Нейман, Бернайс, Гёдель & Классы как объекты \\
MK & 1955 & Морс, Келли & Импредикативные классы \\
KP/KPU & 1965 & Крипке--Платек & Допустимые множества \\
ETCS & 1964 & Ловер & Элементарная теория категории множеств \\
ETCC & 1966 & Ловер & Элементарная теория категорий \\
NF; NFU & 1937; 1969 & Куайн; Йенсен & Стратифицированное свёртывание \\
CZF & 1978 & Акцель & Конструктивная ZF \\
IZF & 1970-е & Фридман & Интуиционистская ZF \\
PA & 1889 & Пеано & Первопорядковая арифметика \\
$\mathsf{Z}_2$ & XX в. & Гильберт--Бернайс & Второпорядковая арифметика \\
MLTT & 1984 & Мартин-Лёф & Зависимые типы \\
CIC & 1988--90 & Кокан--Юэ; Кокан--Полен & Индуктивные конструкции \\
HoTT & 2005+ & Аводей, Воеводский & Унивалентность + типы тождества \\
Кубическая HoTT & 2015+ & Коэн--Кокан и др. & Вычислительная унивалентность \\
Универс-полиморфная HoTT & 2013+ & Воеводский & Полиморфизм универсумов \\
Линейная логика + ! & 1987 & Жирар & Ресурсная чувствительность \\
NBG + AFA & 1988 & Акцель & Необоснованные множества \\
Теория $(\infty, 1)$-топосов & 2009 & Люри & $(\infty, 1)$-топосы \\
Некоммутативная геометрия & 1980-е & Конн & Спектральные тройки \\
Когезивный $(\infty, 1)$-топос & 2013 & Шрайбер & Модальности когезии \\
Мотивная $\mathrm{SH}(k)$ & 1999 & Воеводский, Морель & Мотивная гомотопия \\
Реализуемость (Eff) & 1982 & Хайланд & Эффективный топос \\
Конструктивная математика Маркова & 1950-е--1980-е & А.\,А.~Марков, Нагорный & Конструктивный анализ + принцип Маркова \\
Конструктивный анализ Бишопа & 1967 & Бишоп, Бриджес & Бесвыборный конструктивный вещественный анализ \\
Синтетическая дифф.~геометрия & 1970-е & Кок, Ловер & Инфинитезималии как аксиомы \\
Элементарный высший топос & 2021 & Расех & Элементарный $(\infty, 1)$-топос \\
\bottomrule
\end{longtable}

\paragraph{Страта $\LCls$.} Классификаторы (мета-рамки). $\LCls$ распадается на максимальный подкласс $\Meta_{\mathrm{Cls}}^{\top}$ (категоричный с точностью до $(\infty, \infty)$-эквивалентности по Теореме~\ref{thm:meta-cat}) и его дополнение из \emph{частичных} (немаксимальных) классификаторов, классифицирующих строгие под-стеки $\fM$, а не весь $\fM$. Все перечисленные ниже классификаторы немаксимальны: каждый удовлетворяет (M1)--(M5), но нарушает (Max-1), о чём свидетельствует строго-подмножественный характер его классификационного фокуса.

\begin{longtable}{@{}p{3.3cm}l p{3cm} p{4.4cm}c@{}}
\toprule
Мета-рамка & Год & Автор & Классификационный фокус & В $\Meta_{\mathrm{Cls}}^{\top}$? \\ \midrule
$\infty$-космои & 2022 & Риль--Верити & $(\infty, 1)$-категорные теории & Нет \\
Унивалентные основания & 2010+ & Воеводский & Расширения HoTT & Нет \\
Когезивная рамка & 2013 & Шрайбер & Когезивные $\infty$-топосы & Нет \\
Программа Higher Algebra & 2017+ & Люри & Стабильные $\infty$-кат.\ + операды & Нет \\
Синтетическая математика & 2000-е+ & Тейлор, Шульман & Аксиоматические синтетические основания & Нет \\
\bottomrule
\end{longtable}

\paragraph{Страта $\LAbs$.} Пуста, по AFN-T.

\subsection{Две мета-операции}\label{sec:indexing-scheme}

\paragraph{Горизонтальная и вертикальная мета.} На страте $\mathcal{L}$ действуют две категорно различные мета-операции:
\begin{itemize}
\item \textbf{Горизонтальная (классифицирующая) мета.} $\mathrm{Cls}(\mathcal{L})$ — класс $2$-категорий $\mathbf{F}$, \emph{классифицирующих}, но не \emph{порождающих} члены $\mathcal{L}$: локально малых, снабжённых доступным классификационным функтором в модули $\mathcal{L}$, непорождающих и параметризованных метатеорией (ср.\ (M1)--(M5)).
\item \textbf{Вертикальная (порождающая) мета.} $\mathrm{Gen}(\mathcal{L})$ — класс рамок, \emph{порождающих} члены $\mathcal{L}$ как внутренние выводы --- кандидаты в абсолютные основания.
\end{itemize}

\begin{proposition}[Неколлапс горизонтальной меты на $\LFnd$]\label{prop:no-collapse}
$\mathrm{Cls}(\LFnd)$ \emph{не} вкладывается в $\LFnd$ и не совпадает с $\LAbs$. Следовательно, это отдельный формальный страт, который я обозначаю $\LCls$.
\end{proposition}

\begin{proof}[Схема доказательства.]
Член $\mathbf{F} \in \mathrm{Cls}(\LFnd)$ — $2$-категория Rich-оснований с доступным классификационным функтором в $\fM$. Сама она не является Rich-основанием (не удовлетворяет (R1)--(R5) как порождающая теория), поэтому $\mathbf{F} \notin \LFnd$. Она не в $\LAbs$, поскольку (M4) требует непорождаемости, тогда как $\LAbs$ требует максимальной порождающей способности $(\Pi_{3\text{-max}})$.
\end{proof}

\paragraph{Горизонтальная стабилизация.} Теорема~\ref{thm:meta-stab} утверждает, что итерация $\mathrm{Cls}$ на $\LCls$ стабилизируется \emph{на теоретическом уровне}: дважды итерированные модули реализуют ту же $(\infty, \infty)$-теорию, что и однажды итерированные (пишется $\simeq_2$, с оговоркой о восхождении по универсумам из Теоремы~\ref{thm:meta-stab}(b)), так что $\mathrm{Cls}(\LCls) \simeq_2 \LCls$. Никакого нового горизонтального страта, отличного от $\LCls$, не возникает.

\paragraph{Вертикальный шаг к запретной страте.} $\LAbs \supseteq \mathrm{Gen}(\LFnd)$: член $\LAbs$ был бы максимально порождающим основанием, из которого выводится всё $\LFnd$. Категорно это свежая операция, а не итерация $\mathrm{Cls}$. По AFN-T, $\LAbs = \emptyset$.

\paragraph{Сводка.} Формальная последовательность читается так:
\[
\underbrace{\LFnd}_{\text{(R1)--(R5)}} \; \xrightarrow{\mathrm{Cls}} \; \underbrace{\LCls}_{\text{(M1)--(M5)}} \; \supsetneq \; \underbrace{\LClsMax}_{(\mathrm{Max}\text{-1})\text{--}(\mathrm{Max}\text{-4})} \; \xrightarrow{\mathrm{Gen}} \; \underbrace{\LAbs = \emptyset}_{\text{AFN-T}}
\]

\section{Разумные Rich-метатеории}\label{sec:rs}

\subsection{Условия Rich-метатеории}

\begin{definition}[Разумная Rich-метатеория]\label{def:rs}
Теория $S$, сформулированная в первопорядковом мета-синтаксисе, — \textbf{разумная Rich-метатеория} (или \textbf{R-S}), если она удовлетворяет:
\begin{itemize}
\item \textbf{(R1) Арифметическая полнота}: $S$ интерпретирует арифметику Робинсона $\mathsf{Q}$.
\item \textbf{(R2) Рекурсивная перечислимость}: множество аксиом $S$ рекурсивно перечислимо.
\item \textbf{(R3) Существование модели}: $S$ имеет хотя бы одну модель $\cM$ в объемлющем теоретико-множественном мета-универсуме $\mathbf{U}_2$ (внешний универсум Гротендика Соглашения~\ref{conv:zfc-inacc}, достаточно большой, чтобы вместить модели всех перечисленных R-S-примеров, включая $\ZFC$ + недостижимые кардиналы, чьи наименьшие модели имеют размер $\kappa_1$ и потому лежат в $\mathbf{U}_2 \setminus \mathbf{U}_1$).
\item \textbf{(R4) Гёделево кодирование}: существует рекурсивная инъекция $\ulcorner \cdot \urcorner : \mathrm{Formulas}(S) \to \mathbb{N}$, для которой предикаты «$n$ — гёделев номер доказуемой формулы» и «$n$ — гёделев номер правильно построенной формулы» $S$-определимы.
\item \textbf{(R5) Категорная семантика}: существует $n_S \in \mathbb{N} \cup \{\infty\}$, \emph{внутренняя категорная глубина} $S$, такая что
\begin{itemize}
\item[\textbf{(R5a)}] синтаксическая категория $\Syn(S)$ теории $S$ (объекты: контексты; $1$-морфизмы: доказуемые-из-контекста подстановки; высшие морфизмы: доказательства равенства подстановок) — существенно $\mathbf{U}_1$-малая $(\infty, n_S)$-категория, чья категория моделей $\Mod(S) \simeq \Fun^{\mathrm{lex}}(\Syn(S), -)$ $\kappa_S$-доступна --- эквивалентно, чьё $\mathrm{Ind}_{\kappa_S}$-пополнение $\mathrm{Ind}_{\kappa_S}(\Syn(S))$ есть $\kappa_S$-доступная $(\infty, n_S)$-категория (Адамек--Росицкий~\cite{AdamekRosicky1994}, теорема~2.26; все утверждения о доступности $\Syn(S)$ ниже читаются в этом $\mathrm{Ind}$-пополненном смысле, а утверждения о размере относятся к самой $\Syn(S)$), --- где $\kappa_S \in \{\kappa_1, \kappa_2\}$ определяется теорией $S$: для $S$ силы непротиворечивости $\leq \kappa_1$ (например, $\ZFC$, $\PA$, $\HoTT$, $\MLTT$, $\CIC$) положим $\kappa_S = \kappa_1$; для $S$ силы, требующей второго недостижимого (например, $\ZFC$ + недостижимый), положим $\kappa_S = \kappa_2$. Параметр доступности следует силе непротиворечивости теории; все теоремы ниже формулируются равномерно для произвольного $\kappa_S \in \{\kappa_1, \kappa_2\}$, если конкретный случай не оговорён;
\item[\textbf{(R5b)}] синтаксико-семантическое вычисление $\ev : \Syn(S) \to \cM$ для каждой модели $\cM \models S$ есть $S$-сохраняющий $(\infty, n_S)$-функтор, и пара функторов $(\Syn, \Mod)$ образует сопряжение типа Ламбека--Скотта, единица которого на $\Syn(S)$ является $(\infty, n_S)$-эквивалентностью на \emph{начальную $S$-структурированную $(\infty, n_S)$-категорию}, то есть начальный объект $(\infty, n_S + 1)$-категории $S$-структурированных $(\infty, n_S)$-категорий и $S$-сохраняющих $(\infty, n_S)$-функторов между ними.
\end{itemize}
Сопряжение Ламбека--Скотта (R5a)--(R5b) на уровне $n_S$ устанавливается композицией: Ламбек--Скотт~\cite{LambekScott1986} для $n_S = 1$; Капулкин--Ламсдейн~\cite{KapulkinLumsdaine2021} и Аводей--Уоррен~\cite{AwodeyWarren2009} для $n_S = \infty$ в смысле HoTT; унитарность Барвика--Шоммер-Приса~\cite{BarwickSchommerPries2021} вместе со сравнением моделей Бергнер--Резка~\cite{BergnerRezk2013} и $(\infty, 2)$-выпрямлением Ганьи--Харпаза--Ланари~\cite{GagnaHarpazLanari2021} для общего $n_S \geq 2$. Ниже используются только структурные следствия --- $2$-функториальность $\Syn$ и единственность начального $S$-структурированного объекта с точностью до канонической $(\infty, n_S)$-эквивалентности.
\emph{Замечание о внутренней категорной глубине $n_S$.} Параметр $n_S$ означает наибольшее $k$, на котором $\Syn(S)$ несёт структурно нетривиальные $k$-морфизмы до того, как все высшие клетки становятся обратимыми. Для первопорядковой $S$ имеем $n_S = 1$. Для теории типов Мартина-Лёфа и HoTT с унивалентностью синтаксическая категория есть \emph{некая} $(\infty, 1)$-категория в стандартной терминологии (Капулкин--Ламсдейн~\cite{KapulkinLumsdaine2021}); обозначение $n_S = \infty$ в этой статье относится к \emph{неограниченной итерации образования типов тождества}, дающей нетривиальное содержание высших равенств на каждом конечном уровне $k \in \mathbb{N}$, то есть отсутствию такого $k$, на котором вся структура становится $(k-1)$-усечённой. Это чтение «неограниченной глубины» — итеративное понятие, и оно \emph{не} совпадает с $(\infty, \infty)$-категорным понятием Резка / Барвика--Шоммер-Приса~\cite{BarwickSchommerPries2021} (где $k$-морфизмы на каждом $k$ \emph{необратимы}). Оба чтения согласуются в поведении AFN-T (Теорема~\ref{thm:vertical}), и ни один аргумент статьи не зависит от этого различия.

Примеры, верифицирующие (R5a)--(R5b): первопорядковые $S$ при $n_S = 1$ (Ламбек--Скотт~\cite{LambekScott1986}); теория типов Мартина-Лёфа с $n_S = \infty$ (Аводей--Уоррен~\cite{AwodeyWarren2009}); HoTT с $n_S = \infty$ (Капулкин--Ламсдейн~\cite{KapulkinLumsdaine2021}); зависимые теории типов (Джейкобс~\cite{Jacobs1999}).
\end{itemize}
\end{definition}

\subsection{Примеры R-S}

Ниже перечислены разумные Rich-метатеории. Каждый пункт — \emph{формальная теория} $S$ (аксиоматическая система с р.п.\ аксиоматизацией (R2), интерпретирующая Q (R1)); категорные структуры вроде топосов выступают как \emph{модели} этих теорий, а не как сами теории, и цитируются через свои внутренние языки.
\begin{itemize}
\item $\ZFC$; $\ZFC$ + недостижимые кардиналы; $\NBG$; Морс--Келли; $\NBG + \AFA$ (Акцель).
\item Интуиционистские теории типов $\MLTT$ и $\CIC$; $\HoTT$ с унивалентностью; универс-полиморфная $\HoTT$.
\item \emph{Линейная логика с экспоненциалом} $!$, снабжённая арифметическими константами, достаточными для интерпретации Q (то есть интуиционистская линейная арифметика $\mathsf{ILL}_{\mathsf{PA}}$ или ограниченный линейно-арифметический фрагмент лудики Жирара~\cite{Girard2001}); чисто пропозициональная линейная логика (R1) \emph{не} удовлетворяет и исключена.
\item \emph{Внутренняя зависимая теория типов $\infty$-топосов Люри} (аксиоматическое представление $(\infty, 1)$-топоса как модели гомотопической теории типов с унивалентными универсумами и высшими модальностями, Люри HTT~\cite{LurieHTT} \S 5--\S 6, через соответствие внутренних языков Шульмана~\cite{Shulman2019} и Аводея--Уоррена~\cite{AwodeyWarren2009}); сам $(\infty, 1)$-топос — \emph{каноническая модель} этой теории, а не отдельная теория.
\item \emph{Внутренний язык когезивного $\infty$-топоса} (Шрайбер~\cite{Schreiber2013}): $\HoTT$, расширенная сопряжёнными тройками модальных операторов $(\Pi \dashv \flat \dashv \sharp)$, соответствующих аксиомам когезии, оформленная как теоретико-типовое расширение.
\item \emph{Мотивно-расширенные основания}: $\ZFC$ (или $\HoTT$), расширенная аксиомами существования мотивной стабильной $\infty$-категории $\mathrm{SH}(k)$ и её $\mathrm{A}^1$-гомотопической инвариантности как формальное Rich-расширение; сама мотивная $\infty$-категория — снова модельно-теоретический конструкт, не самостоятельная теория.
\item \emph{Внутренний язык топоса реализуемости} $\mathrm{Eff}$: интуиционистская арифметика высших порядков с реализуемостью (Хайланд~\cite{Hyland1982}), канонической моделью которой служит $\mathrm{Eff}$; ниже используется сокращение «$\mathrm{Eff}$-теория» для этого внутреннего языка.
\item \emph{Конструктивная математика Маркова} (Марков, Нагорный~\cite{MarkovNagornyj1988}: $\mathsf{HA}$ + принцип Маркова, с категорной семантикой в $\mathrm{Eff}$).
\item \emph{Конструктивный анализ Бишопа}~\cite{Bishop1967,BridgesVita2006} (бесвыборный конструктивный вещественный анализ с полной категорной семантикой в любом топосе с объектом натуральных чисел).
\item \emph{Предикативный анализ Фефермана}~\cite{Feferman1964} (системы силы $\mathsf{ATR}_0$ с предикативным свёртыванием).
\end{itemize}

\subsection{Не-примеры}

\begin{itemize}
\item \emph{Аффинная логика без экспоненциала $!$}: нарушает (R1) --- индукция Пеано требует контракции на гипотезе индукции.
\item \emph{Противоречивые теории}: нарушают (R3).
\item \emph{Второпорядковая арифметика с полной семантикой}: нарушает (R2) --- множество общезначимых формул не рекурсивно перечислимо.
\item \emph{Интуитивные «теории всего»} в физике без формальной аксиоматизации: нарушают (R1) и (R2).
\item \emph{Строгий философский ультрафинитизм} (Есенин-Вольпин~\cite{EseninVolpin1961,EseninVolpin1970}): нарушает (R1) при строжайшем чтении --- тотальность функции следования отвергается для «неосуществимо больших» аргументов, так что полная индукция Пеано недоступна. Ограниченные варианты (см.\ пограничные случаи ниже) принадлежат другому страту.
\end{itemize}

\begin{proposition}[Независимость (R1)--(R5)]\label{prop:rs-independence}
Пять условий (R1)--(R5), определяющих класс $\RS$, попарно независимы: для каждого $i \in \{1, 2, 3, 4, 5\}$ существует теория $T_i$, удовлетворяющая всем условиям, кроме (R$i$).
\end{proposition}

\begin{proof}[Свидетели]
\begin{itemize}[nosep]
\item \emph{Нарушение (R1) при выполнении (R2)--(R5)}: \emph{аффинная логика без экспоненциала $!$} удовлетворяет р.п.-аксиоматизируемости (R2), имеет непротиворечивые модели в $\mathbf{U}_1$ (R3), допускает гёделево кодирование своего множества формул (R4) и обладает стандартной категорной семантикой в симметрично-моноидальных категориях (R5 при $n_S = 1$); нарушено (R1): без контракции полная индукция Пеано на гипотезе индукции недоступна, так что арифметика Робинсона Q не интерпретируется в стандартном смысле.
\item \emph{Нарушение (R2) при выполнении (R1), (R3)--(R5)}: \emph{второпорядковая арифметика с полной (внешне-хенкинской) семантикой}. Она тривиально интерпретирует Q (R1), имеет полносемантические модели (R3), допускает гёделево кодирование (R4) и обладает категорной семантикой в $2$-категориях (R5). Но (R2) нарушено: множество теорем, общезначимых при полной семантике, не р.п.\ (следствие второпорядковой неполноты).
\item \emph{Нарушение (R3) при выполнении (R1), (R2), (R4), (R5)}: \emph{противоречивое} р.п.\ расширение Q, например $T_3 := \mathsf{Q} + (0 = 1)$. Над объемлющей метатеорией (R3) для теорий с (R2) \emph{эквивалентно} непротиворечивости: по теореме полноты всякая непротиворечивая р.п.\ теория имеет счётную модель, лежащую в $\mathbf{U}_1 \subseteq \mathbf{U}_2$; обратно, модель свидетельствует непротиворечивость. Поэтому никакая \emph{непротиворечивая} теория не может служить (R3)-свидетелем, и $T_3$ — канонический способ нарушения. Остальные условия для $T_3$ выполнены: (R1), (R2), (R4) синтаксичны и немедленны; для (R5) синтаксическая категория $\Syn(T_3)$ — терминальная (доказуемо коллапсировавшая) структура, $\Mod(T_3) = \emptyset$, и сопряжение Ламбека--Скотта между ними выполняется тривиально (универсально квантифицированные клаузы (R5b) удовлетворены вакуумно, а условие единицы выполняется для коллапсировавшего начального объекта). Тем самым $T_3$ изолирует (R3) чисто.
\item \emph{Нарушение (R4) при выполнении (R1)--(R3), (R5)}: искусственная теория $T$, расширяющая Q в языке, гёделева нумерация которого невычислима (например, через нерекурсивную биекцию между формулами и натуральными числами). (R1)--(R3) и (R5) сохраняются после переигрывания Q-интерпретации; (R4) нарушается, поскольку биекция нерекурсивна.
\item \emph{Нарушение (R5) при выполнении (R1)--(R4)}: \emph{чистая первопорядковая арифметика Пеано $\mathsf{PA}$, снабжённая нестандартной схемой свёртывания}, делающей синтаксическую категорию недоступной (например, раздуванием объектов-контекстов до собственного класса через $\mathbf{U}_2$-размерные области квантификации без соответствующей $(\infty, n)$-функториальной семантики). (R1)--(R4) наследуются от $\mathsf{PA}$; (R5a) нарушается, поскольку $\Syn(T)$ более не $\kappa_1$-доступна. Более канонично: свидетелем служит любая теория, чей формальный язык допускает свёртывание вне натурально-числового страта, но не допускает сопряжения Ламбека--Скотта с категорией моделей.
\end{itemize}
\end{proof}

\begin{itemize}
\item \emph{Ограниченная арифметика и исчисления осуществимости.} Системы $\mathsf{I\Delta_0}$, $\mathsf{S}_2^i$ / $\mathsf{T}_2^i$ Басса~\cite{Buss1986}, $\mathsf{V}_0$, полиномиально-вычислимые арифметики и калькулюсы теоретико-доказательственной осуществимости формализуют ограниченные фрагменты ультрафинитистской интуиции. Они удовлетворяют (R2)--(R5), но реализуют (R1) лишь в ограниченной форме: объемлющая арифметика интерпретирует фрагмент $\mathsf{PA}$ с ограниченными кванторами, а не полную $\mathsf{PA}$. Они допускаются в $\RS$ как \emph{слабые Rich-метатеории}, классифицирующие соответствующий под-стек $\fM$ на ограниченной внутренней категорной глубине $n_S$. Их Морита-отношение к стандартным $\RS$-точкам изучается средствами обратной математики~\cite{Simpson2009}.
\item \emph{Строгий ультрафинитизм (традиция Есенина-Вольпина).} Строжайший философский ультрафинитизм, отвергающий тотальность следования для «неосуществимо больших» аргументов, выпадает из $\RS$ при любом стандартном чтении (R1) (выше, Не-примеры). Ограниченно-арифметические и полиномиальные варианты выше дают формальные реконструкции, допускаемые в $\RS$ как слабые Rich-метатеории.
\item \emph{Паранепротиворечивые теории.} Рассматриваются отдельно в Приложении~\ref{app:paraconsistent}: дополнительная конструкция \emph{классического ядра} допускает паранепротиворечивую $S$ в $\RS$ через семантику Раутли--Мейера, с пере-выводом AFN-T в этом расширенном сеттинге.
\end{itemize}

\subsection{\texorpdfstring{Класс $\cS_S$}{Класс S\_S}}

\begin{definition}[$S$-определимые категорные структуры]\label{def:SS}
Пусть $S \in \RS$. Определим индуктивно:
\begin{align*}
\cS_S^{\mathrm{local}} &:= \{(\cM, Y) : \cM \models S,\ Y \in \Ob(\cM)\} \;\cup\; \{(\Syn(S), X) : X \in \Ob(\Syn(S))\}, \\
\cS_S^{(0)} &:= \cS_S^{\mathrm{local}}, \\
\cS_S^{(k+1)} &:= \cS_S^{(k)} \,\cup\, \mathcal{O}\bigl(\cS_S^{(k)}\bigr),
\end{align*}
где $\mathcal{O}$ — замыкание под следующим фиксированным списком $(\infty, n)$-категорных операций, каждая из которых берёт $\cS_S^{(k)}$-объекты на входе и производит объекты в $\cS_S^{(k+1)}$:
\begin{enumerate}[nosep,label=\textrm{(O\arabic*)}]
\item Левые и правые Кан-расширения вдоль $\cS_S^{(k)}$-морфизмов;
\item Конструкции Гротендика $\int F$ для картезиановых расслоений $F : I \to \cS_S^{(k)}$ с $I \in \cS_S^{(k)}$;
\item Пределы и копределы малых диаграмм (всюду «малость» диаграммы означает: индексирована $(\infty, n)$-категорией, являющейся $\mathbf{U}_2$-малой) со значениями в $\cS_S^{(k)}$;
\item Категории срезов/запятых/противоположные, построенные из $\cS_S^{(k)}$-данных;
\item Производные функторы (в смысле Квиллена) между $\cS_S^{(k)}$-объектами, снабжёнными модельно-категорной структурой (когда применимо; вакуумно на объектах без модельной структуры);
\item Симметрично-моноидальные тензорные произведения и внутренний hom в замкнутых моноидальных $\cS_S^{(k)}$-объектах.
\end{enumerate}
Положим $\cS_S^{\mathrm{global}} := \bigcup_{k < \omega} \cS_S^{(k)}$ и $\cS_S := \cS_S^{\mathrm{local}} \cup \cS_S^{\mathrm{global}}$. Все операции (O1)--(O6) $\kappa_1$-доступны (Адамек--Росицкий~\cite{AdamekRosicky1994}); $\cS_S^{\mathrm{global}}$ $\kappa_1$-доступен и $\mathbf{U}_2$-мал.
\end{definition}

\subsection{\texorpdfstring{Ключевая лемма: $S$-определимость влечёт членство в $\cS_S$}{Ключевая лемма: S-определимость влечёт членство в S\_S}}

\begin{lemma}[Лемма об $S$-определимости]\label{lem:SS-membership}
Если $X$ — формально $S$-определимый объект, то есть $X$ задан формулой $\phi_X \in S$, выражающей его характеристические свойства, то $X \in \cS_S$.
\end{lemma}

\begin{proof}
По Соглашению~\ref{conv:notation} (клауза Морита-редуцируемости) объект $X$, представленный формулой $\phi_X$, рассматривается в $\mathbf{StrCat}_{S, n}$ как заострённая структура $(\Syn(S), [\phi_X])$. Поскольку $[\phi_X] \in \Ob(\Syn(S))$, эта пара лежит в $\cS_S^{\mathrm{local}}$ \emph{по первой клаузе Определения~\ref{def:SS}} --- никакой операции замыкания не требуется. Следовательно, $X \in \cS_S^{\mathrm{local}} \subseteq \cS_S$. (Если $\phi_X$ родна собственной под-теории $\iota_F : F \hookrightarrow S$, индуцированный функтор интерпретации $\Syn(\iota_F) : \Syn(F) \to \Syn(S)$ переносит $[\phi_X]_F$ в $[\phi_X]$, так что пойнтинг можно эквивалентно взять над $F$; вывод не меняется.)
\end{proof}

\subsection{Стратификация по размеру}\label{sec:size-stratification}

\paragraph{Классы размера внутри $\cS_S$.}
\begin{itemize}
\item \textbf{Малые} (элементы $\mathbf{U}_1$): множества ранга $< \kappa_1$.
\item \textbf{Большие} (элементы $\mathbf{U}_2 \setminus \mathbf{U}_1$): множества ранга в $[\kappa_1, \kappa_2)$.
\item \textbf{За пределами $\mathbf{U}_2$}: множества ранга $\geq \kappa_2$, включая $\mathbf{U}_2$-собственные классы.
\end{itemize}

\paragraph{Кардинальные границы объемлющих структур.}
\begin{itemize}
\item Класс $S$-определимости $\cS_S = \cS_S^{\mathrm{local}} \cup \cS_S^{\mathrm{global}}$ индексирован $\mathbf{U}_2$: $\cS_S^{\mathrm{local}} \subseteq \mathbf{U}_1$ (объекты $\mathbf{U}_1$-размерных моделей) и $\cS_S^{\mathrm{global}} \subseteq \mathbf{U}_2$ (производные конструкции над $\cS_S^{\mathrm{local}}$ через $\mathbf{U}_2$-ограниченные функториальные операции).
\item $|\Ob(\cF)/\!\simeq| \leq 2^{\aleph_0} < \kappa_1$; $|\cF| \leq \kappa_2$.
\item $|\Ob(\mathbf{StrCat}_{S, n})/\!\simeq| \leq 2^{\kappa_1} < \kappa_2$.
\item $|\fM / \!\simeq_2| \leq \kappa_2$; $\fM^{(\mathrm{Cls}\cdot k)}$ на уровне $\kappa_k$ (Теорема~\ref{thm:meta-stab}(b)).
\end{itemize}

\section{\texorpdfstring{Условия страты $\LAbs$}{Условия страты L\_Abs}}\label{sec:conditions}

\begin{definition}[$(F_S)$]\label{def:F}
$(\infty, n)$-категория $X$ удовлетворяет $(F_S)$ для $S \in \RS$ тогда и только тогда, когда $\exists F \in \cF$ с:
\begin{itemize}
\item[\textbf{(F-1)}] $\iota_F : F \hookrightarrow S$ (включение под-теории);
\item[\textbf{(F-2)}] $\iota_F^* : \Mod(S) \to \Mod(F)$ (индуцированный модельный редукт);
\item[\textbf{(F-3)}] $\phi_X \in F$ с $\Syn(F)[\phi_X] \simeq_{(\infty, n)} X$.
\end{itemize}
\end{definition}

\begin{definition}[$(\Pi_{4, S, n})$]\label{def:pi4}
$(\Pi_{4, S, n})(X)$ выполняется тогда и только тогда, когда $\forall Y \in \cS_S,\ \forall\ (\infty, n)$-функтора $F : X \to Y$, $F$ не является $(\infty, n)$-эквивалентностью на свой образ.
\end{definition}

\begin{definition}[$(\Pi_{3\text{-max}, S, n})$]\label{def:pi3max}
$(\Pi_{3\text{-max}, S, n})(X)$ выполняется тогда и только тогда, когда $\forall F \in \cF\ \exists\ (\infty, n)$-функтор $F \to X$, эквивалентный на свой образ.
\end{definition}

\begin{definition}[$\LAbs$]\label{def:level6}
$X \in \LAbs$ тогда и только тогда, когда $X \models (F_S) \wedge (\Pi_{4, S, n}) \wedge (\Pi_{3\text{-max}, S, n})$ для некоторого $(S, n) \in \RS \times (\mathbb{N} \cup \{\infty\})$.
\end{definition}

\section{\texorpdfstring{Граничная лемма: пустота $\LAbs$ ($(\alpha)$-часть AFN-T)}{Граничная лемма: пустота L\_Abs ((alpha)-часть AFN-T)}}\label{sec:afnt-alpha}

\begin{theorem}[Граничная лемма: несовместимость $(F_S) \wedge (\Pi_{4, S, n})$; $(\alpha)$-часть AFN-T]\label{thm:afnt-alpha}
Для каждой Rich-метатеории $S \in \RS$ и каждого категорного уровня $n \in \mathbb{N} \cup \{\infty\}$ не существует $(\infty, n)$-категорного объекта $X$, удовлетворяющего тройке $\LAbs$
\[
(F_S) \wedge (\Pi_{4, S, n}) \wedge (\Pi_{3\text{-max}, S, n}).
\]
\end{theorem}

\begin{corollary}[Пустота $\LAbs$]\label{cor:level6-empty-alpha}
$\LAbs = \emptyset$.
\end{corollary}

\begin{proof}
Немедленно из Теоремы~\ref{thm:afnt-alpha} и Определения~\ref{def:level6}.
\end{proof}

\begin{proof}[Доказательство Теоремы~\ref{thm:afnt-alpha}]
Пусть $X$ удовлетворяет $(F_S) \wedge (\Pi_{4, S, n})$. Тогда $X \in \cS_S^{\mathrm{global}} \subseteq \cS_S$ по Лемме~\ref{lem:SS-membership}, и $\id_X : X \to X$ — $(\infty, n)$-эквивалентность с целью в $\cS_S$, что нарушает $(\Pi_{4, S, n})$.
\end{proof}

\section{\texorpdfstring{Граничная лемма: невозможность побега через предел ($(\beta)$-часть AFN-T)}{Граничная лемма: невозможность побега через предел ((beta)-часть AFN-T)}}\label{sec:afnt-beta}

\begin{theorem}[Граничная лемма: невозможность побега через предел; $(\beta)$-часть AFN-T]\label{thm:afnt-beta}
Пусть $\lambda$ — ординал с $\mathrm{cf}(\lambda) < \kappa_2$, регулярным кардиналом (так что конфинальная подпоследовательность $\{A_{\kappa_i}\}_{i < \mathrm{cf}(\lambda)}$, по которой вычисляется копредел, индексирована в $\mathbf{U}_2$; это рабочая граница для категорной техники, никаких ограничений на сам $\lambda$, кроме существования такой конфинальности), и пусть $\{A_\kappa\}_{\kappa < \lambda}$ — трансфинитная последовательность объектов, удовлетворяющая:
\begin{enumerate}[label=(\roman*)]
\item Каждый $A_\kappa$ формально специфицирован внутри некоторого Rich-основания $F_\kappa$ над $S$.
\item Последовательность монотонна в подходящем $(\infty, n)$-категорном смысле (существуют когерентные функторы перехода $A_\kappa \to A_{\kappa'}$ для $\kappa < \kappa'$).
\item Предел $A_\infty := \colim_{\kappa < \lambda} A_\kappa$ существует в объемлющей $(\infty, n)$-категории.
\end{enumerate}
Тогда $A_\infty \in \cS_S^{\mathrm{global}}$, и, следовательно, $A_\infty$ Морита-редуцируем к элементу $\cS_S$. В частности, $A_\infty$ не удовлетворяет $(\Pi_{4, S, n})$.
\end{theorem}

\begin{proof}
Каждый $A_\kappa \in \cS_S$ по Лемме~\ref{lem:SS-membership}; $S$-определимость функторов перехода делает $D : \lambda \to \cS_S$ (со значениями в заострённых структурах Определения~\ref{def:SS}) $S$-индексированной диаграммой. По Определению~\ref{def:SS}, $\cS_S^{\mathrm{global}}$ замкнут под конструкциями Гротендика и Кан-расширениями вдоль $S$-определимых морфизмов; следовательно, $A_\infty = \colim D \in \cS_S^{\mathrm{global}}$. Доступность: по стандартным свойствам замкнутости доступных категорий относительно фильтрованных копределов (Адамек--Росицкий~\cite{AdamekRosicky1994}, в частности замкнутость $\kappa$-доступных категорий под $\lambda$-фильтрованными копределами для регулярного $\lambda \leq \kappa$ и сохранение доступности при доступной переиндексации) копредел наследует $\kappa_1$-доступность от диаграммы. Морита-редукция к каноническому модельно-категорному копределу противоречит $(\Pi_4)$.
\end{proof}

\subsection{Башни собственных классов не дают побега}

\begin{proposition}[Башни собственных классов: дихотомия $\RS$-замыкания]\label{prop:proper-class}
Пусть $\{A_\alpha\}_{\alpha \in \mathrm{Ord}}$ — башня, индексированная собственным классом, где каждый $A_\alpha$ формально специфицирован внутри Rich-основания $F_\alpha$ над метатеорией $S$, и пусть $A_\infty := \colim_{\alpha \in \mathrm{Ord}} A_\alpha$ обозначает подразумеваемый предел. Выполняется ровно одно из следующего.
\begin{itemize}
\item[\textbf{(a)}] \emph{Равномерно порождённая башня.} Существует единая рекурсивная порождающая схема $\psi(\alpha, A)$ в некоторой метатеории $S_{\mathrm{tow}} \supseteq S$, специфицирующая $\alpha \mapsto (F_\alpha, \phi_{A_\alpha})$ равномерно для всех $\alpha \in \mathrm{Ord}$. Если предел $A_\infty$ допускает формулу $\phi_{A_\infty}$ в $S_{\mathrm{tow}}$ и $S_{\mathrm{tow}} \in \RS$, то $A_\infty$ удовлетворяет $(F_{S_{\mathrm{tow}}})$, и Теорема~\ref{thm:afnt-alpha}, применённая к $S_{\mathrm{tow}}$, закрывает Уровень~6 на $A_\infty$.
\item[\textbf{(b)}] \emph{Неравномерно порождённая башня.} Нет единой рекурсивной схемы ни в каком $S' \in \RS$, порождающей башню; сопоставление спецификаций $\alpha \mapsto (F_\alpha, \phi_{A_\alpha})$ не является $S'$-рекурсивно перечислимым ни для какого $S' \in \RS$.
\end{itemize}
В случае (b) башня выпадает из области действия Определения~\ref{def:level6}: предел $A_\infty$ не может свидетельствовать $(F_{S'})$ ни для какого $S' \in \RS$, поскольку $(F_{S'})$ по (F-3) требует формулы $S'$, а (R2) вынуждает р.п.-аксиоматизацию любой такой формулы.
\end{proposition}

\begin{proof}
\textbf{(a)} Рекурсивная схема $\psi$ — формула $S_{\mathrm{tow}}$; предел специфицируется формулой $\phi_{A_\infty}(A) \equiv \forall \alpha\,\psi(\alpha, A) \Rightarrow A = A_\alpha\text{-предел}$. Если $S_{\mathrm{tow}} \in \RS$ (что сохраняется р.п.-расширением, когда не нарушается существование модели (R3)), Определение~\ref{def:F} даёт $(F_{S_{\mathrm{tow}}})(A_\infty)$, и Теорема~\ref{thm:afnt-alpha} закрывает конъюнкцию $(F_{S_{\mathrm{tow}}}) \wedge (\Pi_{4, S_{\mathrm{tow}}, n})$. Башни с $S_{\mathrm{tow}} \notin \RS$ попадают в случай (b).

\textbf{(b)} Предположим, к противоречию, что $A_\infty$ удовлетворяет $(F_{S'})$ для некоторого $S' \in \RS$. По (F-3) существует формула $\phi_{A_\infty} \in S'$, классифицирующая $A_\infty$. По (R2) $S'$ р.п.-аксиоматизирована, так что деривационная структура $\phi_{A_\infty}$ $S'$-рекурсивно перечислима. Развёртка $\phi_{A_\infty}$ по конструкции башни даёт $S'$-р.п.\ порождающую формулу для последовательности $\alpha \mapsto (F_\alpha, \phi_{A_\alpha})$, что противоречит неравномерности. Следовательно, $A_\infty$ не удовлетворяет никакому $(F_{S'})$ для $S' \in \RS$ и не является кандидатом в $\LAbs$.
\end{proof}

\subsection{Конкретные башни}

\paragraph{Категорная башня.} $A_0 = $ категория «шагающая стрелка» $\mathbf{D}$; $A_1 = \Cat$; $A_2 = 2\text{-}\Cat$; $\ldots$; $A_n = (\infty, n)$-$\Cat$. Предел $A_\infty = (\infty, \infty)$-$\Cat$ существует, и по теореме унитарности Барвика--Шоммер-Приса~\cite{BarwickSchommerPries2021}, в сочетании со сравнением моделей Бергнер--Резка~\cite{BergnerRezk2013} и глобулярно-кубической эквивалентностью Ара--Мальциниотиса~\cite{AraMaltsiniotis2017}, башня стабилизируется: каноническое вложение $(\infty, \infty)\text{-}\Cat \hookrightarrow (\infty, \infty+1)\text{-}\Cat$ — эквивалентность. Это вертикальная стабилизация $(\infty, n)$-иерархии.

\paragraph{Топосная башня.} $A_0 = $ $1$-топос; $A_1 = 2$-топос (Шульман~\cite{Shulman2008}); $A_n = (\infty, n)$-топос. Предел $A_\infty$ — стабилизированный $(\infty, \infty)$-топос — определяется здесь в смысле унитарности Барвика--Шоммер-Приса~\cite{BarwickSchommerPries2021}, применённой к башне $A_n$; $(\infty, 1)$-страт следует Люри~\cite{LurieHTT}, $(\infty, 2)$-страт — Шульману~\cite{Shulman2008}, высшие страты — продолжающаяся программа.

\paragraph{Спектральная башня.} $A_0 = C^*$-алгебра; $A_1 = $ спектральная тройка (Конн); $A_n = (\infty, n)$-версия; каждый уровень — производная конструкция над предыдущим.

\subsection{Пространство траекторий}

Более тонкий вариант башенного подхода вводит выбор на каждом шаге: $A_{\kappa+1}$ не определяется однозначно по $A_\kappa$, а зависит от структурного выбора. Пространство $\mathbb{T}$ всех возможных траекторий — большой категорный объект: дерево, симплициальное множество или группоид путей.

\begin{proposition}\label{prop:trajectory-space}
Пространство траекторных выборов допускает (по меньшей мере) четыре представления в разных категорных разрешениях --- это \emph{не} взаимно эквивалентные объекты, но каждое захватывает траекторные данные на своём уровне усечения, и \emph{каждое} лежит в $\cS_S^{\mathrm{global}}$:
\begin{itemize}
\item как группоид путей $\Pi_1(\fM)$, где $\fM$ — классифицирующий $2$-стек Rich-оснований;
\item как симплициальное множество — нерв $\mathrm{N}(\fM)$;
\item как дерево выборов — дерево уточняющих выборов над $\fM$;
\item как категория отношений $\mathrm{Rel}(\fM)$.
\end{itemize}
В каком бы разрешении ни было взято $\mathbb{T}$, вывод ниже применим.
\end{proposition}

\begin{proof}[Доказательство (схема)]
Каждое из четырёх представлений — стандартная производная конструкция над $(\infty, 1)$-топосом Rich-оснований. См.\ Люри~\cite{LurieHTT} \S 3.2 для отождествления группоидов путей с симплициальными множествами; Пронк~\cite{Pronk1996} для интерпретации через бикатегорию дробей; остальные эквивалентности классические.
\end{proof}

\begin{corollary}
Предел любого траекторного башенного подхода лежит в $\cS_S^{\mathrm{global}}$ и потому Морита-редуцируем внутри классифицирующего стека $\fM$. Побега в $\LAbs$ не происходит.
\end{corollary}

\section{Пятиосевая абсолютность}\label{sec:five-axis}

Теорема~\ref{thm:afnt-alpha} устанавливает AFN-T для фиксированной пары $(S, n)$. Теперь я устанавливаю робастность теоремы: она выполняется для всех $S \in \RS$ и всех $n$, а сверх того — для мета-итераций, альтернативных категорных упорядочений и соображений полноты.

Я называю это свойством \emph{пятиосевой абсолютности} AFN-T. Рисунок~\ref{fig:five-axis} суммирует пять осей и их структурные зависимости.

\begin{figure}[ht]
\centering
\begin{tikzpicture}[
  axis/.style={-{Stealth[length=2.5mm]}, thick},
  dep/.style={-{Stealth[length=2mm]}, thin, gray, densely dotted},
  label/.style={font=\footnotesize, align=center, inner sep=2pt},
  center/.style={draw, circle, thick, minimum size=1.8cm, align=center, font=\bfseries\small, fill=violet!8}
]
\node[center] (C) at (0,0) {AFN-T};
\node[label, above=1.9cm of C]          (H) {Горизонтальная \\[1pt] \scriptsize\itshape по $S \in \RS$};
\node[label, below=1.9cm of C]          (V) {Вертикальная \\[1pt] \scriptsize\itshape по $n \in \mathbb{N} \cup \{\infty\}$};
\node[label, left=2.1cm of C]           (M) {Мета-вертикальная \\[1pt] \scriptsize\itshape мета-итерации};
\node[label, right=2.1cm of C]          (L) {Латеральная \\[1pt] \scriptsize\itshape альт.\ упорядочения};
\node[label] (K) at (3.3,2.3) {Полнота \\[1pt] \scriptsize\itshape внутри $\cL\cS$};
\draw[axis] (C) -- (H);
\draw[axis] (C) -- (V);
\draw[axis] (C) -- (M);
\draw[axis] (C) -- (L);
\draw[axis] (C) -- (K);
\draw[dep] (K.west) to[bend right=12] node[pos=0.35, above right=2pt and 6pt, font=\scriptsize, text=gray, inner sep=1pt] {поглощает} (H.east);
\draw[dep] (L.south) to[bend left=20] node[midway, right, font=\scriptsize, text=gray, inner sep=1pt] {$\sim$ вертикальная} (V.east);
\draw[dep] (M.south) to[bend right=20] node[midway, left, font=\scriptsize, text=gray, inner sep=1pt] {$\mu$-замыкание} (V.west);
\end{tikzpicture}
\caption{Пятиосевая абсолютность AFN-T (Теоремы~\ref{thm:horizontal}--\ref{thm:completeness}). Сплошные стрелки — пять осей, вдоль которых AFN-T абсолютна: горизонтальная (Теорема~\ref{thm:horizontal}), вертикальная (Теорема~\ref{thm:vertical}), мета-вертикальная (Теорема~\ref{thm:meta-vertical}), латеральная (Теорема~\ref{thm:lateral}) и --- условно на Lawvere-скоуп --- полнота (Теорема~\ref{thm:completeness}). Пунктирные стрелки отмечают установленные структурные зависимости: ось полноты поглощает горизонтальную; латеральная ось редуцируется к вертикальной через Ара--Мальциниотиса и Барвика--Шоммер-Приса; мета-вертикальная ось есть $\mu$-замыкание вертикальной под $(\infty, \infty)$-стабилизацией. Минимальное логически независимое представление содержит только горизонтальную и вертикальную.}
\label{fig:five-axis}
\end{figure}

\begin{theorem}[Горизонтальная абсолютность]\label{thm:horizontal}
AFN-T (Теорема~\ref{thm:afnt-alpha}) выполняется для каждой $S \in \RS$ равномерно: существует единое доказательство, работающее для всех Rich-метатеорий, удовлетворяющих (R1)--(R5).
\end{theorem}

\begin{proof}
Доказательство Теоремы~\ref{thm:afnt-alpha} состоит из одного применения Леммы~\ref{lem:SS-membership} (синтаксико-семантический мост) вместе с наблюдением, что $\id_X : X \to X$ реализует Морита-редукцию, нарушающую $(\Pi_{4, S, n})$. Лемма~\ref{lem:SS-membership} использует только (R2) и (R5a) (синтаксическая категория существует, и её объекты поддерживают конструкцию пойнтинга Соглашения~\ref{conv:notation}); наблюдение о Морита-редукции использует только категорный тождественный морфизм. Никакие данные $S$ сверх (R1)--(R5) не привлекаются. Следовательно, доказательство применяется равномерно ко всем $S \in \RS$.
\end{proof}

\begin{theorem}[Вертикальная абсолютность]\label{thm:vertical}
AFN-T (Теорема~\ref{thm:afnt-alpha}) выполняется для каждого категорного уровня $n \in \mathbb{N} \cup \{\infty\}$.
\end{theorem}

\begin{proof}
Синтаксико-семантический мост Теоремы~\ref{thm:afnt-alpha} опирается на (R5a) через Лемму~\ref{lem:SS-membership}: $\Syn(F)$ — $S$-определимая $(\infty, n_F)$-категория для каждого $n$. Сопряжение Ламбека--Скотта выполняется равномерно на всех $(\infty, n)$-уровнях (классически для $n = 0, 1$; Люри~\cite{LurieHTT} \S 5--\S 6 для $n = 1$; для общего $(\infty, n)$ — через Барвика--Шоммер-Приса~\cite{BarwickSchommerPries2021} со сравнением моделей Бергнер--Резка~\cite{BergnerRezk2013}).

Для предельного $n = \infty$: $\Syn(F)$ стабилизируется через
\[ (\infty, \infty)\text{-}\Cat = \colim_{n < \infty} (\infty, n)\text{-}\Cat, \]
и отождествление $\Syn(F)$ как $S$-определимой производной конструкции сохраняется в пределе. Морита-редукция $\id_X : X \to X$ — категорный тождественный морфизм — тривиально переносится на любой $(\infty, n)$-уровень.
\end{proof}

\begin{theorem}[Мета-вертикальная стабилизация]\label{thm:meta-vertical}
На категорном уровне $(\infty, \infty)$ любая мета-итерация перечисленных типов стабилизируется: $(\infty, \infty + 1)\text{-}\Cat = (\infty, \infty)\text{-}\Cat$. Следовательно, мета-итерации не дают побега из AFN-T.
\end{theorem}

\begin{proof}
По теореме унитарности Барвика--Шоммер-Приса~\cite{BarwickSchommerPries2021} вместе со сравнениями моделей Бергнер--Резка~\cite{BergnerRezk2013} и Ара--Мальциниотиса~\cite{AraMaltsiniotis2017} башня $(\infty, n)$-категорий стабилизируется на $n = \infty$: дальнейшая итерация не добавляет структуры. Конкретно: каноническое вложение
\[ (\infty, \infty)\text{-}\Cat \hookrightarrow (\infty, \infty + 1)\text{-}\Cat \]
есть эквивалентность.

Для итерированных категорий функторов: $\Fun^k(\cC) \subseteq (\infty, \infty)\text{-}\Cat$ для любого $k$; предел — снова $(\infty, \infty)$-$\Cat$.

Для нестандартных моделей: они требуют либо (a) обогащения теории за пределы R-S (нарушая (R2)), либо (b) редукции через анти-фундированную аксиому Акцеля (AFA) к стандартным необоснованным множествам (Акцель~\cite{Aczel1988}), что не добавляет структурных осей.

Следовательно, все мета-итерации стабилизируются внутри существующих пяти осей.
\end{proof}

\begin{theorem}[Латеральная абсолютность]\label{thm:lateral}
Все альтернативные категорные упорядочения редуцируются к $(\infty, n)$-иерархии (или обогащают её) и потому удовлетворяют AFN-T.
\end{theorem}

\begin{proof}
Каждое из перечисленных упорядочений допускает канонический перевод в $(\infty, n)$-$\Cat$ для подходящего $n$:
\begin{itemize}
\item Операды $\to$ симметрично-моноидальные $(\infty, 1)$-категории (Люри HA~\cite{LurieHA} \S 5).
\item Двойные категории $\to$ $(\infty, 2)$-категории (Грандис).
\item Глобулярные/кубические/опетопные $\omega$-категории $\simeq (\infty, \infty)$-категории (Ара--Мальциниотис~\cite{AraMaltsiniotis2017}).
\item Стабильные $(\infty, 1)$-категории $\to (\infty, 1)$-категории через спектрификацию.
\item Фьюжн-категории $\to$ конечные $(\infty, 1)$-категории.
\end{itemize}
В каждом случае перевод сохраняет формальную определимость: объекты альтернативного упорядочения соответствуют объектам $(\infty, n)$-иерархии на нужном уровне. Следовательно, AFN-T применяется равномерно по всем упорядочениям.
\end{proof}

\begin{definition}[Lawvere-скоуп]\label{def:law-scope}
Основополагающая теория $F$ принадлежит \textbf{Lawvere-скоупу} $\cL\cS$, если выполняются следующие структурные свойства (метатеория, в которой строятся $\Syn(F)$ и $\Mod(F)$, произвольна, но фиксирована на протяжении всех четырёх клауз):
\begin{itemize}
\item \textbf{(L1)} $F$ имеет синтаксическую категорию $\Syn(F)$ --- малую категорию контекстов и доказуемых подстановок.
\item \textbf{(L2)} $F$ имеет семантическую $2$-категорию $\Mod(F)$ --- $2$-категорию моделей.
\item \textbf{(L3)} Существует каноническое сопряжение $\Syn \dashv \Mod$ (Ловер~\cite{Lawvere1969}).
\item \textbf{(L4)} $\Mod(F)$ реализуема как $(\infty, n)$-категория для некоторого $n$ (традиция Гротендика--Люри).
\end{itemize}
Членство в $\cL\cS$ — структурное свойство $F$ и не предполагает никакой конкретной R-S. В частности, каждая R-S $S$ сама лежит в $\cL\cS$ (взять тавтологические $\Syn = S$, $\Mod = \Mod(S)$ из (R5)), так что понятие свободно от циркулярности для применения в Теореме~\ref{thm:completeness}.
\end{definition}

\begin{theorem}[Ось полноты, условно на Lawvere-скоуп]\label{thm:completeness}
Пусть $F \in \cL\cS$. Тогда каждая структурная вариация $F$ редуцируется к одной из четырёх осей (горизонтальной, вертикальной, мета-вертикальной, латеральной). В частности, AFN-T пятиосево абсолютна внутри $\cL\cS$.
\end{theorem}

\begin{proof}
Пусть $F \in \cL\cS$. По Определению~\ref{def:law-scope}, $F$ характеризуется четвёркой:
\begin{itemize}
\item \textbf{(A) Синтаксическая структура $\Syn(F)$}: сигнатура + аксиомы + правила вывода --- \emph{горизонтальная ось} (по R-S).
\item \textbf{(B) Семантическая глубина}: уровень $n$ категории $\Mod(F)$ как $(\infty, n)$-категории --- \emph{вертикальная ось}.
\item \textbf{(C) Мета-рефлексия}: итерации гёделева кодирования и принципов рефлексии --- \emph{мета-вертикальная ось}.
\item \textbf{(D) Альтернативные семантики}: операдные, глобулярные, кубические, опетопные и т.\,д.\ --- \emph{латеральная ось}.
\end{itemize}
По Определению~\ref{def:law-scope} членство в $\cL\cS$ \emph{означает}, что $F$ представлена ровно данными (L1)--(L4), то есть четвёркой (A)--(D); именно здесь живёт условность теоремы --- внутри $\cL\cS$ исчерпываемость (A)--(D) выполняется по определению скоупа, а не как дополнительное утверждение. Значит, любая мнимая пятая ось $\eta$ — подкомпонента одной из (A)--(D). Либо $\eta \subseteq$ (A), и тогда AFN-T для $\eta$ редуцируется к горизонтальной абсолютности (Теорема~\ref{thm:horizontal}); либо $\eta \subseteq$ (B) — редукция к вертикальной (Теорема~\ref{thm:vertical}); либо $\eta \subseteq$ (C) — к мета-вертикальной (Теорема~\ref{thm:meta-vertical}); либо $\eta \subseteq$ (D) — к латеральной (Теорема~\ref{thm:lateral}). Независимой пятой оси внутри $\cL\cS$ нет.
\end{proof}

\section{Три стандартных обходных пути}\label{sec:bypass}

\begin{theorem}[Абсолютность полиморфизма универсумов]\label{thm:universe}
Универс-полиморфные структуры в любой Rich-метатеории $S$ Морита-редуцируемы к производным конструкциям над $(\infty, 1)$-топосом малых категорий. Формально: если $X$ задан как
\[ X = \prod_{\ell : \mathrm{Level}} F(\ell) \]
для семейства функторов $F(\ell) : \cU_\ell \to \cU_\ell$, то
\[ X \simeq \lim_\ell \Fun(\cU_\ell, \cU_\ell) \in \cS_S^{\mathrm{global}}. \]
В частности, $X \in \cS_S$, и универс-полиморфные структуры удовлетворяют $\neg (\Pi_4)$.
\end{theorem}

\begin{proof}
Универс-полиморфная структура $X = \prod_\ell F(\ell)$ — сечение расслоения Гротендика $\pi_{\mathrm{End}} : \int \mathrm{End}(\cU) \to \mathrm{Level}$, ассоциированного (Люри HTT~\cite{LurieHTT} \S 3.2) с \emph{диаграммой эндофункторов} $\ell \mapsto \Fun(\cU_\ell, \cU_\ell)$, где переходные отображения — ограничения вдоль вложений универсумов $\cU_\ell \hookrightarrow \cU_{\ell'}$. Сечения этого расслоения соответствуют совместимым семействам $\{F(\ell)\}_\ell$, и пространство таких сечений есть
\[ \Gamma(\pi) = \lim_\ell \Fun(\cU_\ell, \cU_\ell). \]
Это производная конструкция (индексированный предел), а значит, она принадлежит $\cS_S^{\mathrm{global}}$.

Для $(\infty, n)$-версий: конструкция распространяется через Люри HTT \S 4.2 (пределы в $(\infty, n)$-категориях).
\end{proof}

\paragraph{Рефлексивная башня непротиворечивости.} Второй стандартный обход пытается сбежать от AFN-T через итерированные утверждения непротиворечивости. Полагают
\[ S_{\kappa+1} = S_\kappa + \Con(S_\kappa) \]
и надеются, что предел $S_\infty$ захватывает нечто за пределами отправной точки.

\begin{theorem}[Ограниченность рефлексивной башни]\label{thm:reflective}
Для каждой Rich-метатеории $S$ рефлексивная башня $\{S_\kappa\}$ стабилизируется на границе недостижимого кардинала:
\[ \Con(\text{рефлексивная башня над }S) = \Con(S) + \kappa_{\mathrm{inacc}}, \]
ровно один дополнительный недостижимый кардинал. Побега в $\LAbs$ по этому пути не происходит.
\end{theorem}

\begin{proof}
Доказательство использует два результата.

\textbf{Шаг 1} (верхняя граница). Башня $\{S_\kappa\}$ имеет теоретико-доказательственный ординал, ограниченный первым недостижимым кардиналом над $S$. Это следует из стандартного анализа принципов рефлексии: принцип рефлексии $\Con(S)$ добавляет к силе непротиворечивости $S$ не более одного недостижимого кардинала (Ратьен~\cite{Rathjen1999}, Феферман~\cite{Feferman1964}).

\textbf{Шаг 2} (стабилизация). По $(\infty, \infty)$-мета-вертикальной стабилизации (Теорема~\ref{thm:meta-vertical}) башня не продолжается неограниченно на категорном уровне: $(\infty, \infty + 1)\text{-}\Cat = (\infty, \infty)\text{-}\Cat$.

Сочетая: башня ограничена сверху $\Con(S) + \kappa_{\mathrm{inacc}}$ и стабилизируется на $(\infty, \infty)$-уровне. Любой предел $S_\infty$ остаётся внутри Lawvere-скоупа и внутри $\cS_S$.
\end{proof}

Это самый тонкий из трёх обходных путей: интенсиональное уточнение действительно уточняет классификацию базового уровня. Я разбираю его подробно через построение функтора интенсионального уточнения $\II$ и доказательство его срез-локальности.

Начну с определения целевой $2$-категории display-$2$-категорий.

\begin{definition}[Display-$2$-класс]\label{def:display-class}
Пусть $\cC$ — $2$-категория с $2$-пулбэками. \textbf{Display-класс} $\cD \subseteq \Mor_1(\cC)$ — класс $1$-морфизмов, удовлетворяющий:
\begin{itemize}
\item \textbf{(D1) Содержит эквивалентности}: каждая $1$-эквивалентность в $\cC$ лежит в $\cD$.
\item \textbf{(D2) Устойчивость к пулбэкам}: для $d : B \to A$ из $\cD$ и любого $f : A' \to A$ $2$-пулбэк $f^*d : f^*B \to A'$ существует и лежит в $\cD$.
\item \textbf{(D3) Замкнутость под композицией}: если $d_1 : C \to B \in \cD$ и $d_2 : B \to A \in \cD$, то $d_2 \circ d_1 \in \cD$.
\item \textbf{(D4) $2$-клеточная когерентность}: для $2$-морфизма $\phi : f \Rightarrow g : A' \to A$ и $d \in \cD$ индуцированный $2$-морфизм $\phi^* d : f^* d \Rightarrow g^* d$ — изоморфизм в срезе над $A'$.
\end{itemize}
\end{definition}

\begin{definition}[Display-$2$-категория]\label{def:display-2-cat}
\textbf{Display-$2$-категория} — тройка $(\cC, \cD, \iota)$, где:
\begin{itemize}
\item $\cC$ — локально малая $2$-категория с $2$-пулбэками.
\item $\cD$ — display-класс в $\cC$.
\item $\iota : \cD \hookrightarrow \Mor_1(\cC)$ — каноническое вложение.
\end{itemize}
\end{definition}

\begin{definition}[$2$-категория $\Sint$]\label{def:Sint}
$2$-категория $\Sint$ display-$2$-категорий имеет:
\begin{itemize}
\item \textbf{Объекты}: display-$2$-категории $(\cC, \cD, \iota)$.
\item \textbf{$1$-морфизмы}: пары $(F, F_\cD) : (\cC_1, \cD_1, \iota_1) \to (\cC_2, \cD_2, \iota_2)$, где $F : \cC_1 \to \cC_2$ — $2$-функтор, а $F_\cD : \cD_1 \to \cD_2$ — ограничение $F$, такие что $F \circ \iota_1 = \iota_2 \circ F_\cD$ и $2$-пулбэки сохраняются с точностью до когерентных $2$-изоморфизмов.
\item \textbf{$2$-морфизмы}: естественные $2$-преобразования $\eta : F \Rightarrow G$, ограничения которых $\eta|_\cD : F_\cD \Rightarrow G_\cD$ когерентно согласованы.
\end{itemize}
\end{definition}

\begin{theorem}[Существование $\II$]\label{thm:I-existence}
Существует контравариантный $2$-функтор
\[ \II : \cF^{\mathrm{op}} \to \Sint \]
из $2$-категории Rich-оснований ($\LFnd$) в $\Sint$ со следующими свойствами:
\begin{enumerate}
\item[\textbf{(I-1)}] \textbf{Гомотопическая инвариантность}: каждая $2$-эквивалентность $\phi : f \simeq g$ в $\cF$ индуцирует $2$-изоморфизм $\II(\phi) : \II(f) \simeq \II(g)$ в $\Sint$.
\item[\textbf{(I-2)}] \textbf{Калибровочная ковариантность}: для калибровочного преобразования $\tau : F_1 \to F_2$ морфизм $\II(\tau) : \II(F_2) \to \II(F_1)$ — $1$-морфизм в $\Sint$, являющийся $1$-эквивалентностью, когда $\tau$ обратимо.
\item[\textbf{(I-3)}] \textbf{Строгое уточнение Морита}: существуют $F_1, F_2 \in \cF$ с $F_1 \sim_M F_2$ (Морита-эквивалентны), но $\II(F_1) \not\simeq \II(F_2)$.
\item[\textbf{(I-4)}] \textbf{Морита как $2$-локализация}: забывающий $2$-функтор $\cU : \Sint \to 2\text{-}\Cat$, $(\cC, \cD, \iota) \mapsto \cC$, удовлетворяет $\cU \circ \II \simeq \rho$ (классификационный функтор); Морита-эквивалентность восстанавливается $2$-локализацией $\II$ по классу $\mathcal{W}_\cU = \cU^{-1}(1\text{-equiv}(2\text{-}\Cat))$.
\end{enumerate}
\end{theorem}

\begin{proof}

\textbf{Шаг 1} (построение на объектах). Для $F \in \cF$ строю $\II(F) = (\cC_F, \cD_F, \iota_F) \in \Ob(\Sint)$:
\begin{itemize}
\item $\cC_F := \cF / F$ — $2$-категория среза: объекты — пары $(F', p : F' \to F)$; $1$-морфизмы — когерентные треугольники над $F$; $2$-морфизмы — $2$-клетки, совместимые с проекцией на $F$.
\item $\cD_F$ — \textbf{канонически минимальный display-класс}: наименьший класс $1$-морфизмов в $\cC_F$, содержащий все $1$-эквивалентности и замкнутый под $2$-пулбэками, композицией и $2$-клетками. Формально,
\[ \cD_F := \bigcap \{\cD' \subseteq \Mor_1(\cC_F) : \cD' \text{ — display-класс, содержащий все } 1\text{-эквивалентности}\}. \]
Пересечение непусто (содержит класс всех $1$-морфизмов) и само является display-классом по аргументам о замкнутости пересечений.

\textbf{Явное описание $\cD_F$ через индуктивную конструкцию}:
\begin{itemize}
\item $\cD_F^{(0)}$: все $1$-эквивалентности.
\item $\cD_F^{(n+1)}$: объекты, получаемые из $\cD_F^{(n)}$ $2$-пулбэком вдоль произвольных морфизмов, композицией элементов $\cD_F^{(n)}$ или переносом вдоль $2$-клеток.
\item $\cD_F := \bigcup_{n < \omega} \cD_F^{(n)}$.
\end{itemize}
Это определение даёт класс, автоматически удовлетворяющий (D1)--(D4) по построению, не требуя от объемлющей $\cC_F$ локальной картезиановой замкнутости или экспонент.

\item $\iota_F : \cD_F \hookrightarrow \Mor_1(\cC_F)$ — каноническое вложение.
\end{itemize}

\textbf{Шаг 2} (построение на $1$-морфизмах). Для $f : F_1 \to F_2$ в $\cF$ положим
\[ \II(f) := (f^*, f^*|_\cD) : (\cF / F_2, \cD_{F_2}, \iota_{F_2}) \to (\cF / F_1, \cD_{F_1}, \iota_{F_1}), \]
где $f^* : \cF / F_2 \to \cF / F_1$ — пулбэк-$2$-функтор. Направление контравариантно: $\II$ переводит $f$ в $1$-морфизм противоположного направления.

То, что $f^*$ сохраняет элементы display-класса, следует индукцией по порождению класса: базовый класс — $1$-эквивалентности, сохраняемые пулбэком; индуктивные шаги замыкания сохраняются пулбэком по свойству пулбэка-от-пулбэка.

\textbf{Шаг 3} (построение на $2$-морфизмах). Для $\phi : f \Rightarrow g$ в $\cF$ индуцированное естественное $2$-преобразование $\phi^* : f^* \Rightarrow g^*$ (структура Бека--Шевалле) когерентно ограничивается на display-классы по (D4). Это даёт
\[ \II(\phi) : \II(f) \Rightarrow \II(g) \]
в $\Sint$.

\textbf{Шаг 4} (функториальность). $\II(\id_F) = \id_{\II(F)}$ немедленно. $\II(g \circ f) \simeq \II(f) \circ \II(g)$ следует из когерентности Бека--Шевалле композиции пулбэков.

\textbf{Шаг 5} (проверка (I-1)). $2$-эквивалентность $\phi : f \simeq g$ имеет $2$-обратную $\phi^{-1}$. Индуцированное $\phi^* : f^* \Rightarrow g^*$ обратимо через $(\phi^{-1})^*$. Следовательно, $\II(\phi)$ — $2$-изоморфизм в $\Sint$.

\textbf{Шаг 6} (проверка (I-2)). Для калибровочного преобразования $\tau : F_1 \to F_2$ ($1$-морфизма с выделенной $2$-эквивалентностью $\rho$-проекций) пулбэк $\tau^*$ индуцирует $\II(\tau) : \II(F_2) \to \II(F_1)$. Если $\tau$ обратимо, $\II(\tau)$ — $1$-эквивалентность по Шагу 5.

\textbf{Шаг 7} (проверка (I-3)). Я предъявляю классический пример экстенсионального коллапса.

\emph{Вычислительная рамка.} Чтобы точно сформулировать инвариант $\tau$, я фиксирую модель вычислимости: работаю внутри \emph{эффективного топоса} $\mathrm{Eff}$ (Хайланд~\cite{Hyland1982}), рассматриваемого как эталонная внутренняя модель конструктивной вычислимости. Все $2$-функторы и $2$-эквивалентности в под-$2$-категории $\Sint^{\mathrm{eff}} \subseteq \Sint$ обязаны быть внутренне вычислимыми в $\mathrm{Eff}$; эквивалентно, и функтор, и его квазиобратный реализуются рекурсивными морфизмами. Это ограничение консервативно для нижеследующего аргумента, поскольку Морита-эквивалентность Хофмана~\cite{Hofmann1995} строится явно примитивно-рекурсивными трансляциями и потому лежит в $\Sint^{\mathrm{eff}}$.

Пусть $F_1 := \MLTT$ (теория типов Мартина-Лёфа с пропозициональными типами тождества) и $F_2 := \ETT$ (экстенсиональная теория типов с рефлексией дефинициального равенства).

\emph{Морита-эквивалентность}: $F_1 \sim_M F_2$ на уровне $(1, 1)$-категорных моделей следует из теоремы консервативности Хофмана (Хофман~\cite{Hofmann1995}, глава~3): синтаксическая трансляция между $\MLTT + \mathsf{UIP}$ и $\ETT$ доказуемо консервативна и поднимается до Морита-эквивалентности их $(1, 1)$-категорных категорий моделей на классе локально картезианово замкнутых категорий, где корректно определена соответствующая семантика типов тождества.

\emph{Интенсиональное различие}: определим типизационный инвариант $\tau : \Ob(\Sint^{\mathrm{eff}}) \to \{0, 1\}$ формулой
\[ \tau(\cC, \cD, \iota) := \begin{cases} 1 & \text{если существует $\mathrm{Eff}$-внутренняя процедура нормализации для всех } d \in \cD, \\ 0 & \text{иначе.} \end{cases} \]

Инвариант $\tau$ — $2$-инвариант $\Sint^{\mathrm{eff}}$-эквивалентности: любая вычислимая $2$-эквивалентность $(F, F_\cD) : (\cC_1, \cD_1) \simeq (\cC_2, \cD_2)$ в $\Sint^{\mathrm{eff}}$ переносит процедуру нормализации вдоль реализатора $F_\cD^{-1}$. Ограничение вычислимостью существенно: без него абстрактная $2$-эквивалентность могла бы, в принципе, отождествить нормализующий display-класс с ненормализующим через невычислимую биекцию. Внутри $\mathrm{Eff}$ такая биекция не существует как внутренний морфизм.

Применяя к $\II(F_1), \II(F_2)$:
\begin{itemize}
\item $\tau(\II(F_1)) = 1$: $\MLTT$ имеет сильно нормализующую проверку типов (Штрайхер~\cite{Streicher1991}, теорема 3.3; Вернер~\cite{Werner1994}).
\item $\tau(\II(F_2)) = 0$: $\ETT$ с правилом рефлексии имеет неразрешимую проверку типов (Хофман~\cite{Hofmann1995}, глава 3), и $\mathrm{Eff}$-внутреннего нормализатора не существует.
\end{itemize}

Следовательно, $\tau(\II(F_1)) \neq \tau(\II(F_2))$, так что $\II(F_1) \not\simeq \II(F_2)$ в $\Sint^{\mathrm{eff}}$, несмотря на $F_1 \sim_M F_2$.

\textbf{Шаг 8} (проверка (I-4)). Для $F \in \cF$,
\[ \cU(\II(F)) = \cC_F = \cF / F. \]
По $2$-категорной лемме Йонеды (Келли~\cite{Kelly1982}, теорема 2.4.1), $\cF / F \simeq \rho(F)$ с точностью до $2$-эквивалентности. Значит, $\cU \circ \II \simeq \rho$.

Для Морита-$2$-локализации: положим
\[ \mathcal{W}_\cU := \{(F, F_\cD) \in \Mor_1(\Sint) : F \in 1\text{-equiv}(2\text{-}\Cat)\}. \]
Класс $\mathcal{W}_\cU$ удовлетворяет условиям $2$-исчисления дробей (Пронк~\cite{Pronk1996}, Келли--Стрит~\cite{KellyStreet1974}): $2$-замкнутость под композицией, содержание тождеств, совместимость с $2$-пулбэками. По теореме Пронк $2$-локализация $\Sint[\mathcal{W}_\cU^{-1}]$ существует с универсальным свойством:
\[ \Sint[\mathcal{W}_\cU^{-1}] \simeq \image(\cU)_{\mathrm{gauge}} = \fM \]
(последняя эквивалентность следует из калибровочного отождествления $\fM$ как $\Trace(\mathsf{A})/\mathrm{gauge}$ в смысле стека модулей).

Тем самым Морита-эквивалентность восстанавливается из $\II$ через $2$-локализацию.
\end{proof}

\begin{theorem}[Срез-локальность $\II$]\label{thm:slice-locality}
Пусть $\II : \cF^{\mathrm{op}} \to \Sint$ — функтор Теоремы~\ref{thm:I-existence}, и пусть $\pi : \cF \to \fM$ — калибровочный фактор. Тогда:
\begin{enumerate}
\item[\textbf{(a)}] Существует $2$-функтор $\tpi : \Sint \to \fM$, делающий диаграмму
\[
\begin{tikzcd}
\cF \arrow[r, "\II"] \arrow[d, "\pi"'] & \Sint \arrow[d, "\tpi"] \\
\fM \arrow[r, equal] & \fM
\end{tikzcd}
\]
$2$-коммутативной.
\item[\textbf{(b)}] Слой $\tpi^{-1}([F])$ над точкой $[F] \in \fM$ — $2$-категория $\mathrm{Int}([F])$, \emph{интенсиональный слой над калибровочным классом}.
\item[\textbf{(c)}] Интенсиональное уточнение не производит новых точек в $\fM$: для каждого объекта $X$ из образа $\II$ точка $\tpi(X) \in \fM$ уже существует.
\end{enumerate}
\end{theorem}

\begin{proof}

\textbf{(a)} Положим $\tpi := [\, \cdot \,]_{\mathrm{gauge}} \circ \cU$, где $[\, \cdot \,]_{\mathrm{gauge}} : 2\text{-}\Cat \to \fM$ — калибровочный фактор на объемлющих $2$-категориях. По (I-4), $\tpi \circ \II = [\rho(\cdot)]_{\mathrm{gauge}} = \pi$, поскольку калибровочные классы $F$ и $\rho(F)$ совпадают по определению.

\textbf{(b)} Для $[F] \in \fM$,
\[ \mathrm{Int}([F]) := \tpi^{-1}([F]) = \{(\cC, \cD, \iota) \in \Sint : [\cC]_{\mathrm{gauge}} = [F]_{\mathrm{gauge}}\}. \]
Это полная $2$-подкатегория $\Sint$: её $2$-морфизмы наследуются. Нетривиальность $\mathrm{Int}([F])$ следует из (I-3): $\II(\MLTT)$ и $\II(\ETT)$ лежат в $\mathrm{Int}([\MLTT])$, но не эквивалентны.

Структура $2$-расслоения: для $[F] \in \fM$ и $f : [F'] \to [F]$ картезиановы поднятия существуют через $2$-пулбэк в $\Sint$; картезиановы поднятия компонуются с точностью до когерентного $2$-изо (стандартная теория $2$-расслоений Гротендика; Эрмида~\cite{Hermida1999}, Бакович~\cite{Bakovic2010}).

\textbf{(c)} Для $\alpha_{\mathrm{new}} \in \image(\II)$ имеем $\alpha_{\mathrm{new}} = \II(F)$ для некоторого $F \in \cF$. Тогда $\tpi(\alpha_{\mathrm{new}}) = \pi(F) = [F]_{\mathrm{gauge}} \in \fM$ --- существующая точка. Новой точки базы не возникает.
\end{proof}

\begin{corollary}[Интенсиональное уточнение — уровня среза]\label{cor:slice-level}
Интенсиональное уточнение не добавляет AFN-T новой структурной оси: база ($\fM$) не затронута, и TH-абсолютность выполняется без модификаций.
\end{corollary}

\begin{proof}
Пятиосевая абсолютность (Теоремы~\ref{thm:horizontal}--\ref{thm:completeness}) касается точек $\fM$. Теорема~\ref{thm:slice-locality}(c) показывает, что интенсиональное уточнение действует лишь на слоях, не на точках базы. Новой оси вариации не вводится.
\end{proof}

\subsection{Интенсиональная градуировка}\label{sec:intensional-grading}

Теорема о срез-локальности представляет каждый слой $\mathrm{Int}([F])$ как голую $2$-категорию. Индуктивная конструкция канонического display-класса в Теореме~\ref{thm:I-existence} (Шаг 1) на деле наделяет каждый слой более тонкой, функториальной структурой: $\mathbb{N}$-индексированной градуировкой, измеряющей, \emph{насколько далеко от эквивалентности} находится display-данное. Этот подраздел фиксирует градуировку и два её структурных свойства; оба доказательства коротки, поскольку порождающая индукция Теоремы~\ref{thm:I-existence} уже несёт их неявно.

\begin{definition}[Display-грейд]\label{def:display-grade}
Пусть $F \in \cF$ и $\cD_F = \bigcup_{n < \omega} \cD_F^{(n)}$ — канонически минимальный display-класс с его порождающей фильтрацией (Теорема~\ref{thm:I-existence}, Шаг 1). Для $d \in \cD_F$ определим \emph{грейд}
\[ \gr(d) := \min\{\, n : d \in \cD_F^{(n)} \,\} \in \mathbb{N}. \]
Так, $\gr(d) = 0$ тогда и только тогда, когда $d$ — $1$-эквивалентность, и $\gr(d) = n+1$ тогда и только тогда, когда $d$ впервые возникает на стадии $n+1$, то есть как $2$-пулбэк, композит или $2$-клеточный перенос данных грейда $\leq n$, но само не имеет грейда $\leq n$. Для объекта $\II(F) = (\cC_F, \cD_F, \iota_F) \in \Sint$ обозначим $\mathrm{Int}^{(n)}([F]) \subseteq \mathrm{Int}([F])$ полную под-$2$-категорию слоя (Теорема~\ref{thm:slice-locality}(b)) на display-данных грейда $\leq n$.
\end{definition}

\begin{proposition}[Градуировка структурна]\label{prop:grading}
\begin{enumerate}[label=(\alph*)]
\item \emph{Исчерпывающая фильтрация}: $\mathrm{Int}^{(0)}([F]) \subseteq \mathrm{Int}^{(1)}([F]) \subseteq \cdots$ с $\mathrm{Int}([F]) = \colim_{n} \mathrm{Int}^{(n)}([F])$.
\item \emph{Функториальность (невозрастание грейда)}: для каждого $1$-морфизма $f : F_1 \to F_2$ в $\cF$ пулбэк-компонента $f^*|_{\cD}$ морфизма $\II(f)$ удовлетворяет $\gr(f^* d) \leq \gr(d)$ для всех $d \in \cD_{F_2}$; следовательно, $\II(f)$ ограничивается на каждую стадию градуировки, $\II(f) : \mathrm{Int}^{(n)}([F_2]) \to \mathrm{Int}^{(n)}([F_1])$.
\item \emph{Субаддитивность}: $\gr(d_2 \circ d_1) \leq \max(\gr(d_1), \gr(d_2)) + 1$ и $\gr(\phi^* d) \leq \gr(d) + 1$ для $2$-клеточных переносов $\phi$.
\end{enumerate}
\end{proposition}

\begin{proof}
(a) немедленно из $\cD_F = \bigcup_n \cD_F^{(n)}$. (b) — содержание Шага 2 Теоремы~\ref{thm:I-existence}: сохранение display-данных под $f^*$ доказывается там \emph{индукцией по порождению класса} --- базовый класс ($1$-эквивалентности, грейд $0$) сохраняется пулбэком, и каждый индуктивный шаг замыкания сохраняется постадийно по свойству пулбэка-от-пулбэка; чтение той индукции с удержанием индекса стадии даёт $f^*(\cD_{F_2}^{(n)}) \subseteq \cD_{F_1}^{(n)}$ для каждого $n$, что и есть утверждение. (c) — определение порождающего шага: композиты и $2$-клеточные переносы данных, доступных на стадии $\max(\gr(d_1), \gr(d_2))$ (соотв.\ $\gr(d)$), зачисляются на следующей стадии.
\end{proof}

\begin{corollary}[Экстенсиональность — слепота к грейду]\label{cor:grade-collapse}
Морита-$2$-локализация из (I-4) обращает ровно слой грейда $0$: локализующий класс $\mathcal{W}_\cU$ состоит из морфизмов, чей подлежащий $2$-функтор — эквивалентность, а на display-данных это элементы грейда $0$. Следовательно, экстенсиональный фактор $\Sint[\mathcal{W}_\cU^{-1}] \simeq \fM$ отождествляет все грейды: композит $\mathrm{Int}^{(n)}([F]) \hookrightarrow \mathrm{Int}([F]) \to \{[F]\}$ постоянен по $n$. Интенсиональное содержание основания — в точности башня $\{\mathrm{Int}^{(n)}([F])\}_{n}$, которую забывает экстенсиональная проекция.
\end{corollary}

\begin{remark}[Чтение и открытый вопрос]\label{rem:grading-open}
Градуировка даёт слоям среза Теоремы~\ref{thm:slice-locality} ту структуру, которую пример экстенсионального коллапса (Теорема~\ref{thm:I-existence}, Шаг 7) неявно эксплуатирует: типизационный инвариант $\tau$ вычисляется на display-данных \emph{ограниченного грейда}, и различие $\MLTT$/$\ETT$ видно уже на конечных стадиях башни. Открытыми остаются два естественных вопроса: (i) \emph{строгость} --- для каких $F$ фильтрация $\mathrm{Int}^{(n)}([F])$ строго возрастает по $n$ (без коллапса грейда на конечной стадии)? Коллапс на стадии $n_0$ выражал бы сильное свойство нормальной формы display-геометрии $F$; (ii) \emph{грейд как инвариант Морита-уточнения} --- отделяет ли минимальная стадия коллапса $n_0(F) \in \mathbb{N} \cup \{\infty\}$, инвариантная относительно $\Sint$-эквивалентности по Предложению~\ref{prop:grading}(b), Морита-эквивалентные основания за пределами бинарного инварианта $\tau$ Шага 7, давая количественную шкалу уточнения между экстенсиональной тождественностью и интенсиональной различимостью.
\end{remark}

\section{\texorpdfstring{Мета-классификация $\LCls$}{Мета-классификация классификаторов}}\label{sec:meta}

\subsection{\texorpdfstring{Класс $\Meta_{\mathrm{Cls}}$}{Класс Meta\_Cls}}

\begin{definition}[$\LCls$-мета-структура]\label{def:meta}
Объект $\mathbf{F}$ — \textbf{$\LCls$-мета-структура}, если он удовлетворяет:
\begin{itemize}
\item \textbf{(M1) Локально малая $2$-категория}: $\mathbf{F}$ — локально малая $2$-категория.
\item \textbf{(M2) Классификационный функтор}: есть $2$-функтор $\Cl_{\mathbf{F}} : \mathbf{F} \to \fM$ со значениями в классифицирующем $2$-стеке Rich-оснований; пишем $\fM(\mathbf{F}) := \image(\Cl_{\mathbf{F}}) \subseteq \fM$ для под-стека, высекаемого классификацией $\mathbf{F}$.
\item \textbf{(M3) Доступная рефлексия}: есть доступный эндофунктор $R_{\mathbf{F}} : \mathbf{F} \to \mathbf{F}$.
\item \textbf{(M4) Непорождаемость}: $\mathbf{F}$ не выводит аксиом оснований $F \in \image(\Cl_{\mathbf{F}})$ из своих аксиом; классификация параметрическая, не аксиоматическая.
\item \textbf{(M5) Параметризация метатеорией}: $\mathbf{F}$ определён внутри некоторой $S \in \RS$.
\end{itemize}
Класс таких $\mathbf{F}$ обозначаю $\Meta_{\mathrm{Cls}}$.
\end{definition}

\begin{definition}[Максимальность]\label{def:maximality}
Мета-структура $\mathbf{F} \in \Meta_{\mathrm{Cls}}$ \textbf{максимальна}, если дополнительно:
\begin{itemize}
\item \textbf{(Max-1) Полная классификация}: $\image(\Cl_{\mathbf{F}}) = \fM$.
\item \textbf{(Max-2) Калибровочная полнота}: $\mathrm{Aut}_2(\mathbf{F})$ действует транзитивно на калибровочных классах в $\image(\Cl_{\mathbf{F}})$.
\item \textbf{(Max-3) Стратификация по глубине}: $\mathbf{F}$ допускает \emph{глубинно-индексированную фильтрацию} $\mathbf{F}^{(0)} \hookrightarrow \mathbf{F}^{(1)} \hookrightarrow \cdots \hookrightarrow \mathbf{F}^{(n)} \hookrightarrow \cdots$ доступных $2$-подкатегорий с $\mathbf{F} \simeq_2 \colim_{n < \omega} \mathbf{F}^{(n)}$, снабжённую \emph{функцией глубины} $\dep : \Ob(\mathbf{F}) \to \omega$, сопоставляющей каждому объекту наименьшее $n$, при котором объект лежит в $\mathbf{F}^{(n)}$, и удовлетворяющую следующему условию \emph{предикативности}: для каждого предиката $P$, входящего в определение объекта $A \in \mathbf{F}^{(n)}$ свёртыванием (то есть определение имеет вид «$A :=$ объект, классифицирующий $\{x : P(x)\}$»), всякий объект $Y$, выступающий целью под-выражения в $P$ --- в частности, всякая цель $R_\mathbf{F}$, $\Cl_\mathbf{F}$ или вычислений классификатора подобъектов --- удовлетворяет $\dep(Y) < n$.
\item \textbf{(Max-4) Интенсиональная полнота с консервативностью}: существует $2$-функтор $\II_{\mathbf{F}} : \mathbf{F}^{\mathrm{op}} \to \Sint$, удовлетворяющий свойствам (I-1)--(I-4) Теоремы~\ref{thm:I-existence}, и притом пара $(\Cl_{\mathbf{F}}, \II_{\mathbf{F}})$ \emph{совместно консервативна}: два параллельных $1$-морфизма $\mathbf{F}$, отождествляемые и $\Cl_{\mathbf{F}}$, и $\II_{\mathbf{F}}$, равны (с точностью до обратимой $2$-клетки). Канонический классификатор $\mathbf{F}_S^*$ Теоремы~\ref{thm:meta-cat} удовлетворяет этой клаузе по построению (его морфизмы — пары из базового морфизма и морфизма интенсионального слоя), так что клауза не опустошает класс.
\end{itemize}
Класс максимальных мета-структур обозначаю $\Meta_{\mathrm{Cls}}^{\top}$.
\end{definition}

\subsection{Мета-категоричность}

\begin{theorem}[Мета-категоричность]\label{thm:meta-cat}
Пусть $\mathbf{F}_1, \mathbf{F}_2 \in \Meta_{\mathrm{Cls}}^{\top}$ параметризованы одной Rich-метатеорией $S \in \RS$. Тогда $\mathbf{F}_1 \simeq_{(\infty, \infty)} \mathbf{F}_2$ через каноническую эквивалентность.
\end{theorem}

\begin{proof}
Я показываю, что каждый $\mathbf{F}_i$ канонически $(\infty, \infty)$-эквивалентен общему объекту $\mathbf{F}_S^*$, построенному из $S$, откуда $\mathbf{F}_1 \simeq_{(\infty, \infty)} \mathbf{F}_S^* \simeq_{(\infty, \infty)} \mathbf{F}_2$.

\textbf{Шаг 1: канонический максимальный классификатор.} Пусть $\rho : \cF \to (\infty, n_S)\text{-}\Cat$ — классификационный функтор Соглашения~\ref{conv:notation}. Рассмотрим гротендиково выпрямление, применённое к картезианову расслоению $\int \rho \to S$, с калибровочным фактором:
\[ \mathbf{F}_S^* := \mathbf{St}(\rho) / \mathrm{gauge}, \]
где $\mathbf{St}(\rho)$ — выпрямление $\rho$.

\textbf{Шаг 2: построение $\Phi_\mathbf{F}$.} Любая $\mathbf{F} \in \Meta_{\mathrm{Cls}}^{\top}$ несёт четыре структуры: классификационный функтор $\Cl_{\mathbf{F}} : \mathbf{F} \to \fM$, существенно сюръективный на $\fM$ по (Max-1); транзитивное действие $\mathrm{Aut}_2$ на калибровочных классах по (Max-2); каноническую глубинную фильтрацию $\mathbf{F}^{(\bullet)}$ по (Max-3); каноническое интенсиональное пополнение $\II_{\mathbf{F}} : \mathbf{F}^{\mathrm{op}} \to \Sint$ с (I-1)--(I-4) по (Max-4).

Каждая структура допускает универсальное представление на $\mathbf{F}_S^*$: классификация — тождество на $\fM$; калибровка канонична; глубинная фильтрация — $\mathbf{F}_S^{*(n)} := \{[\alpha] : \dep_S^*([\alpha]) < n\}$, где $\dep_S^*([\alpha])$ — кванторный ранг определяющей формулы $F_{[\alpha]}$ в $\Syn(S)$ (конечен по (R2), (F-3); предикативность (Max-3) следует индукцией по арифметической иерархии Клини); $\II_{\mathbf{F}_S^*}$ — вложение Йонеды, факторизующееся через $\tilde\pi$ (Теорема~\ref{thm:slice-locality}). По универсальному свойству гротендикова выпрямления (Люри HTT~\cite{LurieHTT}, теорема 3.2.0.1) в $\mathbf{StrCat}_{S, 2}$ существует единственный (с точностью до канонического $(\infty, 2)$-изоморфизма) $2$-функтор
\[ \Phi_\mathbf{F} : \mathbf{F} \to \mathbf{F}_S^*, \]
сохраняющий все четыре структуры.

\textbf{Шаг 3: $\Phi_\mathbf{F}$ существенно сюръективен.} Для каждого калибровочного класса $[\alpha] \in \fM$ условие (Max-1) даёт $A \in \mathbf{F}$ с $\Cl_\mathbf{F}(A) = [\alpha]$, и $\Phi_\mathbf{F}(A)$ проецируется в тот же класс в $\mathbf{F}_S^*$. Транзитивность действия $\mathrm{Aut}_2(\mathbf{F})$ на калибровочных классах (Max-2) обеспечивает покрытие всего $\mathbf{F}_S^*$.

\textbf{Шаг 4: $\Phi_\mathbf{F}$ вполне унивалентен.} Я доказываю, что для любых $A, B \in \mathbf{F}$ индуцированное отображение
\[ \Phi_\mathbf{F} : \Hom_\mathbf{F}(A, B) \longrightarrow \Hom_{\mathbf{F}_S^*}(\Phi_\mathbf{F}(A), \Phi_\mathbf{F}(B)) \]
есть эквивалентность Hom-пространств.

\emph{Точность}: пусть $f, g : A \to B$ в $\mathbf{F}$ имеют $\Phi_\mathbf{F}(f) = \Phi_\mathbf{F}(g)$. Применение $\Cl_\mathbf{F}$ к обеим сторонам даёт $\Cl_\mathbf{F}(f) = \Cl_\mathbf{F}(g)$ в $\fM$, так что $f$ и $g$ различаются не более чем калибровкой. Применение $\II_\mathbf{F}$ (удовлетворяющего (I-1) гомотопической инвариантности и (I-2) калибровочной ковариантности Теоремы~\ref{thm:I-existence}) даёт $\II_\mathbf{F}(f) = \II_\mathbf{F}(g)$ в $\Sint$. По клаузе совместной консервативности (Max-4) пара $(\Cl_\mathbf{F}, \II_\mathbf{F})$ разделяет параллельные морфизмы $\mathbf{F}$. Следовательно, $f = g$. (Совместная консервативность — аксиома максимальности, а не следствие одних (I-3)--(I-4): (I-3) — различение на уровне объектов и само по себе морфизмов не разделяет.)

\emph{Полнота}: дан $h : \Phi_\mathbf{F}(A) \to \Phi_\mathbf{F}(B)$ в $\mathbf{F}_S^*$; поднимаю $h$ в $\mathbf{F}$. Морфизм $h$ факторизуется как $h_{\mathrm{ext}} \in \Hom_\fM(\Cl_\mathbf{F}(A), \Cl_\mathbf{F}(B))$ (экстенсиональная часть, через каноническую проекцию $\mathbf{F}_S^* \to \fM$) вместе с $h_{\mathrm{int}} \in \Hom_{\Sint}(\II_{\mathbf{F}_S^*}(B), \II_{\mathbf{F}_S^*}(A))$ (интенсиональный слой, через $\II_{\mathbf{F}_S^*}$). По (Max-1) $h_{\mathrm{ext}}$ поднимается до калибровочного класса морфизмов $A \to B$ в $\cF$ на классифицированном уровне; транзитивность (Max-2) выбирает представителя в $\mathbf{F}$. По (Max-4) $h_{\mathrm{int}}$ поднимается единственно в интенсиональное пополнение $\mathbf{F}$, которое по (I-4) «Морита как $2$-локализация» совпадает с интенсиональными данными $\II_\mathbf{F}$. Комбинируя, получаю морфизм $\tilde h : A \to B$ в $\mathbf{F}$ с $\Phi_\mathbf{F}(\tilde h) = h$.

\emph{$2$-морфизмы}: тот же аргумент этажом выше --- с использованием (M3)-доступности рефлексии и (Max-4)-интенсиональной $2$-функториальности --- устанавливает Hom-эквивалентность на уровне $2$-клеток. Полная когерентная $(\infty, 2)$-эквивалентность следует индукцией по уровням усечения.

Итак, $\Phi_\mathbf{F}$ — $2$-эквивалентность в $\mathbf{StrCat}_{S, 2}$.

\textbf{Шаг 5: поднятие до $(\infty, \infty)$.} $2$-эквивалентность $\Phi_\mathbf{F}$, установленная в Шагах 3--4, распространяется поуровнево: на каждом усечении $\tau_{\leq n}$ для $n \in \mathbb{N}$ я получаю $(\infty, n)$-эквивалентность по (M3)-доступности всех задействованных структур в сочетании с теоремой унитарности Барвика--Шоммер-Приса~\cite{BarwickSchommerPries2021}. Семейство $\{\tau_{\leq n}(\Phi_\mathbf{F})\}_{n \in \mathbb{N}}$ совместимо с переходами усечения по критерию Уайтхеда (Лемма~\ref{lem:whitehead}). Взятие копредела через $(\infty, \infty)$-стабилизацию (Теорема~\ref{thm:bergner-lurie-stab}) даёт $(\infty, \infty)$-эквивалентность $\Phi_\mathbf{F} : \mathbf{F} \xrightarrow{\simeq_{(\infty, \infty)}} \mathbf{F}_S^*$.

\textbf{Заключение.} Применяя Шаг~5 к обоим $\mathbf{F}_1$ и $\mathbf{F}_2$: композиция
\[ \mathbf{F}_1 \xrightarrow{\Phi_{\mathbf{F}_1}} \mathbf{F}_S^* \xrightarrow{\Phi_{\mathbf{F}_2}^{-1}} \mathbf{F}_2 \]
— каноническая $(\infty, \infty)$-эквивалентность.
\end{proof}

\subsection{Структурная множественность}

\begin{theorem}[Структурная множественность]\label{thm:meta-mult}
В $\Meta_{\mathrm{Cls}} \setminus \Meta_{\mathrm{Cls}}^{\top}$ существуют по меньшей мере три попарно не-$2$-эквивалентные мета-структуры:
\begin{itemize}
\item $\mathbf{F}_{\mathrm{univalent}}$: программа Унивалентных оснований (библиотека Coq \emph{Foundations} Воеводского~\cite{VoevodskyFoundations2010}; HoTT Book~\cite{HoTTBook2013}; Капулкин--Ламсдейн~\cite{KapulkinLumsdaine2021}; Шульман~\cite{Shulman2019} — унивалентные универсумы в каждом гротендиковом $\infty$-топосе).
\item $\mathbf{F}_{\mathrm{cosmoi}}$: $\infty$-космои Риля--Верити~\cite{RiehlVerity2022}.
\item $\mathbf{F}_{\mathrm{cohesive}}$: когезивная выс\-ше-топосная рамка Шрайбера~\cite{Schreiber2013}.
\end{itemize}
\end{theorem}

\begin{proof}

\textbf{Шаг 1} (каждая — в $\Meta_{\mathrm{Cls}}$). Каждая из $\mathbf{F}_{\mathrm{univalent}}$, $\mathbf{F}_{\mathrm{cosmoi}}$, $\mathbf{F}_{\mathrm{cohesive}}$ удовлетворяет (M1)--(M5) Определения~\ref{def:meta}. Проверяю:
\begin{itemize}
\item \textbf{(M1)}: каждая — локально малая $2$-категория (стандартно из соответствующих определений).
\item \textbf{(M2)}: каждая имеет классификационный функтор, образ которого — конкретный под-$2$-стек $\fM$:
\begin{align*}
\image(\Cl_{\mathbf{F}_{\mathrm{univalent}}}) &= \{[\alpha_F] : F \supseteq \MLTT + \text{Унивалентность}\}, \\
\image(\Cl_{\mathbf{F}_{\mathrm{cosmoi}}}) &= \{[\alpha_F] : F \text{ — } (\infty, 1)\text{-категорная теория}\}, \\
\image(\Cl_{\mathbf{F}_{\mathrm{cohesive}}}) &= \{[\alpha_F] : F \text{ поддерживает модальности когезии}\}.
\end{align*}
\item \textbf{(M3)}: каждая поддерживает доступную рефлексию --- усечение $\tau_{\leq n}$ для $\mathbf{F}_{\mathrm{univalent}}$, гомотопическую рефлексию для $\mathbf{F}_{\mathrm{cosmoi}}$, модальности когезии для $\mathbf{F}_{\mathrm{cohesive}}$.
\item \textbf{(M4)}: ни одна не выводит аксиом классифицируемых ею оснований; все параметрические.
\item \textbf{(M5)}: каждая параметризована подходящей $S \in \RS$.
\end{itemize}

\textbf{Шаг 2} (попарная неэквивалентность). Я предъявляю конкретные различия классификационных образов:
\begin{itemize}
\item $\mathbf{F}_{\mathrm{univalent}}$ против $\mathbf{F}_{\mathrm{cosmoi}}$: $[\alpha_{\mathrm{NCG}}] \in \image(\Cl_{\mathbf{F}_{\mathrm{cosmoi}}})$ (NCG — $(\infty, 1)$-теория через спектральные тройки), но $[\alpha_{\mathrm{NCG}}] \notin \image(\Cl_{\mathbf{F}_{\mathrm{univalent}}})$ (NCG — не HoTT-расширение).
\item $\mathbf{F}_{\mathrm{univalent}}$ против $\mathbf{F}_{\mathrm{cohesive}}$: стандартная HoTT без когезии лежит в $\image(\Cl_{\mathbf{F}_{\mathrm{univalent}}})$, но не в $\image(\Cl_{\mathbf{F}_{\mathrm{cohesive}}})$.
\item $\mathbf{F}_{\mathrm{cosmoi}}$ против $\mathbf{F}_{\mathrm{cohesive}}$: некогезивные $(\infty, 1)$-теории лежат в $\image(\Cl_{\mathbf{F}_{\mathrm{cosmoi}}})$, но не в $\image(\Cl_{\mathbf{F}_{\mathrm{cohesive}}})$.
\end{itemize}
Неэквивалентность утверждается в $\Meta_{\mathrm{Cls}}$ \emph{над} $\fM$: морфизмы мета-структур обязаны коммутировать с классификационными функторами (с точностью до обратимой $2$-клетки). Будь любые две из трёх $2$-эквивалентны над $\fM$, их классификационные образы совпадали бы как под-стеки $\fM$, что противоречит этим наблюдениям. (Как голые $2$-категории, абстрактные эквивалентности, игнорирующие $\Cl$, не исключены — и нерелевантны для утверждения о плюрализме.)

\textbf{Шаг 3} (ни одна не максимальна). Ни одна из трёх не удовлетворяет (Max-1) (полная классификация), поскольку их классификационные образы — строгие подмножества $\fM$. Следовательно, $\mathbf{F}_{\mathrm{univalent}}, \mathbf{F}_{\mathrm{cosmoi}}, \mathbf{F}_{\mathrm{cohesive}} \notin \Meta_{\mathrm{Cls}}^{\top}$, и Теорема~\ref{thm:meta-cat} к ним не применяется.
\end{proof}

\begin{corollary}[Плюрализм в $\LCls$]\label{cor:plurality}
$\LCls$ структурно плюралистичен: $\Meta_{\mathrm{Cls}}$ содержит несколько неэквивалентных членов.
\end{corollary}

\subsection{Мета-классификационная стабилизация}

\begin{theorem}[Мета-классификационная стабилизация]\label{thm:meta-stab}
Определим $\fM^{(\mathrm{Cls})}$ как классифицирующий $2$-стек $\LCls$-мета-структур:
\[ \fM^{(\mathrm{Cls})} := \Meta_{\mathrm{Cls}} / \text{мета-калибровка}. \]
Тогда:
\begin{enumerate}
\item[\textbf{(a)}] $\fM^{(\mathrm{Cls})} \in \Meta_{\mathrm{Cls}}$ (классифицирующий $2$-стек сам принадлежит $\LCls$).
\item[\textbf{(b)}] \textbf{Идемпотентность на теоретическом уровне}: $\fM^{(\mathrm{Cls}\cdot 2)} := \Meta_{\mathrm{Cls}}^{(\text{для }\fM^{(\mathrm{Cls})})} / \text{мета-калибровка}$ эквивалентен $\fM^{(\mathrm{Cls})}$ \emph{как $(\infty, \infty)$-теория} --- то есть они реализуют одни и те же $(\infty, \infty)$-теоретические модули в смысле унитарности Барвика--Шоммер-Приса~\cite{BarwickSchommerPries2021}. Два стека \emph{не} отождествляются как теоретико-множественные объекты: $\fM^{(\mathrm{Cls}\cdot 2)}$ живёт на один шаг гротендикова универсума выше $\fM^{(\mathrm{Cls})}$, и башня мета-итераций $\{\fM^{(\mathrm{Cls}\cdot k)}\}_{k \geq 1}$ восходит по иерархии недостижимых $\kappa_1 < \kappa_2 < \cdots$; эквивалентность $\simeq_2$ фиксирует, что каждый уровень инстанциирует одну и ту же $(\infty, \infty)$-теорию, а не что уровни теоретико-множественно тождественны.
\item[\textbf{(c)}] Никакой эскалации в $\LAbs$ через мета-итерацию не происходит.
\end{enumerate}
\end{theorem}

\begin{proof}

\textbf{(a)}: $\fM^{(\mathrm{Cls})}$ — фактор-$2$-стек над $\Meta_{\mathrm{Cls}}$. Его условия:
\begin{itemize}
\item \textbf{(M1)}: локально мал (наследуется от $\Meta_{\mathrm{Cls}}$).
\item \textbf{(M2)}: классификационный функтор — тождество на калибровочных классах (мета-структуры классифицируют сами себя).
\item \textbf{(M3)}: доступная рефлексия $R_{\fM^{(\mathrm{Cls})}}$ индуцируется из $R_{\mathbf{F}}$ покомпонентно.
\item \textbf{(M4)}: непорождаемость по наследству.
\item \textbf{(M5)}: параметризован той же $S \in \RS$.
\end{itemize}
Следовательно, $\fM^{(\mathrm{Cls})} \in \Meta_{\mathrm{Cls}}$.

\textbf{(b)}: Я редуцирую мета-стабилизацию к теореме $(\infty, \infty)$-стабилизации (Теорема~\ref{thm:bergner-lurie-stab}) через формальное вложение по категорной глубине, а не по аналогии.

\emph{Шаг (b.1): категорная глубина $\Meta_{\mathrm{Cls}}$.} Каждая $\mathbf{F} \in \Meta_{\mathrm{Cls}}$ есть, по (M1)--(M2), локально малая $2$-категория вместе с классификационным функтором $\Cl_{\mathbf{F}}$ в $\fM$. Поскольку $\fM$ — $(\infty, \infty)$-стек (Соглашение~\ref{conv:notation}), а $\Cl_{\mathbf{F}}$ обязан быть доступным (M3), $\mathbf{F}$ лежит в $(\infty, \infty)$-категории доступных $(\infty, \infty)$-предпучков на $\fM$. Обозначим эту объемлющую категорию $\Pi_{(\infty, \infty)}$. Тогда $\Meta_{\mathrm{Cls}} \hookrightarrow \Pi_{(\infty, \infty)}$.

\emph{Шаг (b.2): классифицирующий стек $\Meta_{\mathrm{Cls}}$ живёт уровнем выше.} Фактор $\fM^{(\mathrm{Cls})} := \Meta_{\mathrm{Cls}} / \mathrm{gauge}_{\mathrm{meta}}$ — по конструкции Гротендика--Люри (Люри HTT~\cite{LurieHTT} \S 3.2) — $(\infty, \infty)$-стек над $\fM$; рассматриваемый как объект, классифицирующий объекты $\Pi_{(\infty, \infty)}$, он лежит в категории на уровень выше, $\Pi_{(\infty, \infty + 1)}$.

\emph{Шаг (b.3): итерация добавляет ещё один уровень.} По той же конструкции $\fM^{(\mathrm{Cls}\cdot 2)} := \Meta_{\mathrm{Cls}}^{(\text{для } \fM^{(\mathrm{Cls})})} / \mathrm{gauge}_{\mathrm{meta}}$ — классифицирующий стек уровнем выше $\fM^{(\mathrm{Cls})}$, то есть объект $\Pi_{(\infty, \infty + 2)}$.

\emph{Шаг (b.4): применение $(\infty, \infty)$-стабилизации.} По Теореме~\ref{thm:bergner-lurie-stab} (Барвик--Шоммер-Прис~\cite{BarwickSchommerPries2021}) канонические вложения
\[ \Pi_{(\infty, \infty)} \hookrightarrow \Pi_{(\infty, \infty + 1)} \hookrightarrow \Pi_{(\infty, \infty + 2)} \]
суть эквивалентности $(\infty, \infty)$-категорий: поскольку категории доступных функторов сохраняют эквивалентности в целевой $(\infty, \infty)$-категории (а ограничение доступности инвариантно относительно скачков усечения $(\infty, n) \hookrightarrow (\infty, n+1)$), эквивалентность Теоремы~\ref{thm:bergner-lurie-stab} на самой $(\infty, \infty)$-$\Cat$ переносится на категорию доступных предпучков $\Pi_{(\infty, \infty)} = \Fun^{\mathrm{acc}}(\fM, (\infty, \infty)\text{-}\Cat)$. Поэтому канонические $2$-функторы
\[ \iota_{\mathrm{meta}} : \fM^{(\mathrm{Cls})} \hookrightarrow \fM^{(\mathrm{Cls}\cdot 2)}, \quad \pi_{\mathrm{meta}} : \fM^{(\mathrm{Cls}\cdot 2)} \to \fM^{(\mathrm{Cls})} \]
удовлетворяют $\pi_{\mathrm{meta}} \circ \iota_{\mathrm{meta}} = \id$ и $\iota_{\mathrm{meta}} \circ \pi_{\mathrm{meta}} \simeq_2 \id$, где последнее — образ стабилизационной эквивалентности под конструкцией Гротендика--Люри.

\emph{Шаг (b.5): заключение.} $\fM^{(\mathrm{Cls}\cdot 2)} \simeq_2 \fM^{(\mathrm{Cls})}$, что и утверждалось.

\textbf{(c)}: Предположим, к противоречию, что $\fM^{(\mathrm{Cls}\cdot k)}$ (для некоторого $k \geq 1$) удовлетворяет условиям $\LAbs$ $(F_S) \wedge (\Pi_{4, S, n}) \wedge (\Pi_{3\text{-max}, S, n})$ (Определение~\ref{def:level6}). По (b) $\fM^{(\mathrm{Cls}\cdot k)}$ реализует ту же $(\infty, \infty)$-теорию, что и $\fM^{(\mathrm{Cls})} \in \Meta_{\mathrm{Cls}}$ (хотя и на более высоком уровне гротендикова универсума $\kappa_k$). Как $\LCls$-мета-структура, $\fM^{(\mathrm{Cls}\cdot k)}$ классифицирует, но не порождает Rich-основания (M4), что противоречит $\LAbs$-максимальной порождаемости $(\Pi_{3\text{-max}, S, n})$. Противоречие показывает: эскалации в $\LAbs$ не происходит.
\end{proof}

\begin{corollary}[Замкнутость итераций]\label{cor:iteration-closure}
Для каждого $k \geq 1$ стек $\fM^{(\mathrm{Cls}\cdot k)}$ реализует ту же $(\infty, \infty)$-теорию, что и $\fM^{(\mathrm{Cls})}$ (обозначение $\simeq_2$ зарезервировано за этой теоретико-уровневой эквивалентностью Теоремы~\ref{thm:meta-stab}(b)). Мета-итерация стабилизируется на $\LCls$ \emph{как теория}; её теоретико-множественная реализация восходит по иерархии гротендиковых универсумов $\kappa_k$ с ростом $k$.
\end{corollary}

\begin{corollary}[Замкнутость $\Meta_{\mathrm{Cls}}$ под мета-классификацией]\label{cor:meta-closure}
Класс $\Meta_{\mathrm{Cls}}$ замкнут под итерацией классификации: итерированная мета-классификация $\LCls$-мета-структур остаётся в $\LCls$.
\end{corollary}

\section{AC/OC-двойственность и дуальная граничная лемма}\label{sec:ac-oc-duality}

Сопряжение Ламбека--Скотта $(\Syn, \Mod)$ из (R5b), собранное по $\RS$ и над $n \in \mathbb{N} \cup \{\infty\}$, индуцирует $(\infty, \infty)$-эквивалентность между теоретико-центричными модулями $\fM$ и акт-центричными модулями $\fM_\cE$ (Теорема~\ref{thm:ac-oc-duality}). Отсюда следует дуальная граничная лемма $\LAbsE = \emptyset$ (Теорема~\ref{thm:dual-afnt}).

\begin{figure}[ht]
\centering
\begin{tikzpicture}[
  cat/.style={draw, rounded corners=2pt, minimum width=3.6cm, minimum height=0.85cm, align=center, font=\small},
  stack/.style={draw, rounded corners=2pt, minimum width=3.6cm, minimum height=0.85cm, align=center, font=\small, fill=blue!6},
  strat/.style={draw, rounded corners=1.5pt, minimum width=3.2cm, minimum height=0.55cm, align=center, font=\footnotesize, fill=violet!8},
  emptybox/.style={draw, dashed, rounded corners=1.5pt, minimum width=3.2cm, minimum height=0.55cm, align=center, font=\footnotesize, fill=red!6},
  func/.style={-{Stealth[length=2.2mm]}, thick},
  equiv/.style={thick, double, double distance=1.5pt},
  node distance=0.4cm
]
% Top row: 2-categories
\node[cat, fill=yellow!12] (F) at (-3.6,3.2) {$\cF$ \\[-1pt] \scriptsize\itshape Rich-основания (R1--R5)};
\node[cat, fill=green!12] (E) at (3.6,3.2) {$\cE$ \\[-1pt] \scriptsize\itshape заострённые рефлексивные акты};
\draw[func] ([yshift=5pt]F.east) -- node[midway, above, font=\footnotesize] {$\varepsilon$} ([yshift=5pt]E.west);
\draw[func] ([yshift=-5pt]E.west) -- node[midway, below, font=\footnotesize] {$\alpha$} ([yshift=-5pt]F.east);
\node[font=\scriptsize] at (0, 3.2) {$\dashv$};

% Middle row: moduli stacks
\node[stack] (M) at (-3.6, 1.6) {$\fM$ \\[-1pt] \scriptsize\itshape $= \cF / \mathrm{gauge}$};
\node[stack] (ME) at (3.6, 1.6) {$\fME$ \\[-1pt] \scriptsize\itshape $= \cE / \mathrm{gauge}_\cE$};
\draw[func, dashed, gray!60!black] (F) -- node[midway, left, font=\scriptsize, gray!60!black] {фактор} (M);
\draw[func, dashed, gray!60!black] (E) -- node[midway, right, font=\scriptsize, gray!60!black] {фактор} (ME);
\draw[equiv] (M.east) -- node[midway, above, font=\footnotesize] {$\simeq_{(\infty, \infty)}$} node[midway, below, font=\scriptsize, gray!60!black] {Теор.~\ref{thm:ac-oc-duality}} (ME.west);

% Bottom rows: strata correspondence
\node[strat] (L5F) at (-3.6, 0.3) {$\LFnd$};
\node[strat] (L5E) at (3.6, 0.3) {$\LFndE$};
\node[strat] (LCF) at (-3.6, -0.4) {$\LCls$};
\node[strat] (LCE) at (3.6, -0.4) {$\LClsE$};
\node[strat] (LMF) at (-3.6, -1.1) {$\LClsMax$};
\node[strat] (LME) at (3.6, -1.1) {$\LClsMaxE$};
\node[emptybox] (L6F) at (-3.6, -1.8) {$\LAbs = \emptyset$};
\node[emptybox] (L6E) at (3.6, -1.8) {$\LAbsE = \emptyset$};

\draw[equiv, gray!50!black, shorten >=2pt, shorten <=2pt] (L5F) -- (L5E);
\draw[equiv, gray!50!black, shorten >=2pt, shorten <=2pt] (LCF) -- (LCE);
\draw[equiv, gray!50!black, shorten >=2pt, shorten <=2pt] (LMF) -- (LME);
\draw[equiv, red!60!black, shorten >=2pt, shorten <=2pt] (L6F) -- node[midway, above, font=\scriptsize, red!60!black] {Теор.~\ref{thm:dual-afnt}} (L6E);

% Connect stacks to strata
\draw[gray!40, densely dotted] (M.south) -- (L5F.north);
\draw[gray!40, densely dotted] (ME.south) -- (L5E.north);
\draw[gray!40, densely dotted] (L5F.south) -- (LCF.north);
\draw[gray!40, densely dotted] (L5E.south) -- (LCE.north);
\draw[gray!40, densely dotted] (LCF.south) -- (LMF.north);
\draw[gray!40, densely dotted] (LCE.south) -- (LME.north);
\draw[gray!40, densely dotted] (LMF.south) -- (L6F.north);
\draw[gray!40, densely dotted] (LME.south) -- (L6E.north);

\end{tikzpicture}
\caption{AC/OC-двойственность и индуцированная биекция стратов. На уровне $2$-категорий (верхний ряд) сопряжённая пара $\varepsilon \dashv \alpha$ (Определение~\ref{def:enactments}) связывает теоретико-центричную $\cF$ с акт-центричной $\cE$: $\varepsilon$ — секция синтаксического само-акта $F \mapsto (F, \Syn(F), \id, \id)$; $\alpha$ — забывающий функтор $(F, \cC, \iota, r) \mapsto F$. Переход к калибровочным классам (средний ряд) даёт $(\infty, \infty)$-эквивалентность классифицирующих $2$-стеков $\fM \simeq \fME$ (Теорема~\ref{thm:ac-oc-duality}) — через полную унивалентность $\varepsilon$ и существенную сюръективность посредством выбранного рефлектора $r$. При этой эквивалентности четыре страты соответствуют биективно (нижние ряды): $\LFnd \leftrightarrow \LFndE$, $\LCls \leftrightarrow \LClsE$, $\LClsMax \leftrightarrow \LClsMaxE$ и $\LAbs = \emptyset \leftrightarrow \LAbsE = \emptyset$ (Теорема~\ref{thm:dual-afnt}). Симметрия формализует примирение объект-центричного и действие-центричного взглядов на основания математики: ни один не первичен; оба описывают один и тот же классифицирующий стек.}
\label{fig:ac-oc-duality}
\end{figure}

\subsection{\texorpdfstring{$2$-категория Rich-актов}{2-категория Rich-актов}}

\begin{definition}[$2$-категория актов $\cE$]\label{def:enactments}
\emph{$2$-категория Rich-актов} $\cE$ имеет:
\begin{itemize}
\item \textbf{Объекты}: четвёрки $(F, \cC, \iota, r)$, называемые \emph{заострёнными рефлексивными актами над $F$}, где $F \in \cF$, $\cC \in \mathbf{StrCat}_{S, n_F}$ (Определение~\ref{def:strcat}) — $\kappa_1$-доступная $S$-модельная категория с $F$-модельной структурой, $\iota : \Syn(F) \to \cC$ — $F$-сохраняющий вполне унивалентный $(\infty, n_F)$-функтор на реплетную под-$(\infty, n_F)$-категорию, и $r : \cC \to \Syn(F)$ — \emph{выбранный} левый сопряжённый (рефлектор) к $\iota$ с единицей $\eta_r : \id_\cC \Rightarrow \iota \circ r$ и коединицей $\varepsilon_r : r \circ \iota \Rightarrow \id_{\Syn(F)}$, удовлетворяющими треугольным тождествам в $\mathbf{StrCat}_{S, n_F}$. Рефлектор $r$ — часть данных; он однозначно определён по $\iota$ с точностью до единственной обратимой $2$-клетки (Риль--Верити~\cite{RiehlVerity2022}, предл.~2.1.11; Адамек--Росицкий~\cite{AdamekRosicky1994}, теорема~1.26).
\item \textbf{$1$-морфизмы} $(\phi, \Phi) : (F_1, \cC_1, \iota_1, r_1) \to (F_2, \cC_2, \iota_2, r_2)$: пара из $1$-морфизма $\phi : F_1 \to F_2$ в $\cF$ (точная интерпретация) и $S$-сохраняющего $(\infty, n)$-функтора $\Phi : \cC_1 \to \cC_2$, снабжённого обратимой $2$-клеткой $\Phi \circ \iota_1 \Rightarrow \iota_2 \circ \Syn(\phi)$;
\item \textbf{$2$-морфизмы}: когерентные естественные преобразования $(\phi, \Phi) \Rrightarrow (\phi', \Phi')$, совместимые со структурными клетками, то есть пары $(\tau, T)$ с $\tau : \phi \Rightarrow \phi'$ в $\cF$ и $T : \Phi \Rrightarrow \Phi'$, удовлетворяющие очевидному пятиугольнику с $\iota_1, \iota_2, \Syn(\tau)$.
\end{itemize}
Имеются канонический забывающий $2$-функтор и каноническая секция:
\[
  \alpha : \cE \to \cF, \quad (F, \cC, \iota, r) \mapsto F; \qquad
  \varepsilon : \cF \to \cE, \quad F \mapsto \bigl(F,\, \Syn(F),\, \id_{\Syn(F)},\, \id_{\Syn(F)}\bigr).
\]
По построению $\alpha \circ \varepsilon = \id_\cF$ строго. Для читаемости я иногда пишу $(F, \cC, \iota)$ вместо $(F, \cC, \iota, r)$, когда рефлектор подразумевается.
\end{definition}

\begin{definition}[Калибровка актов]\label{def:gauge-E}
Пусть $\mathcal{R}_\cE \subseteq \Mor_1(\cE)$ — класс \emph{рефлексионных морфизмов}: $1$-морфизмов вида $(\tau, \Phi)$, в которых $\tau$ — калибровочное преобразование в $\cF$ (Соглашение~\ref{conv:notation}), а композит $r_2 \circ \Phi \circ \iota_1 : \Syn(F_1) \to \Syn(F_2)$ — $(\infty, n)$-эквивалентность. Класс $\mathcal{R}_\cE$ содержит все тождества и все $1$-эквивалентности, замкнут под композицией (рефлекторы компонуются вдоль обратимых коединиц), устойчив под существующими в $\cE$ $2$-пулбэками (покомпонентно, поскольку калибровочные преобразования и эквивалентности пулбэк-устойчивы) и насыщен; следовательно, он удовлетворяет условиям исчисления дробей (BF1)--(BF5) Леммы~\ref{lem:pronk}. Положим
\[ \mathrm{gauge}_\cE := \text{класс $2$-локализации, порождённый } \mathcal{R}_\cE, \qquad \fM_\cE := \cE[\mathcal{R}_\cE^{-1}] \]
через бикатегорию дробей Пронк (Лемма~\ref{lem:pronk}). Два акта \emph{gauge$_\cE$-эквивалентны} тогда и только тогда, когда они соединены зигзагом рефлексионных морфизмов. В частности, рефлексия $(\id_F, r) : (F, \cC, \iota, r) \to \varepsilon(F)$ обращается в $\fM_\cE$ --- это, и только это, добавляет «переход к калибровке»: единица $\eta_r$ собственной рефлексии \emph{не} обратима в самой $\cE$.
\end{definition}

\subsection{Акт-синтаксико-семантическая лемма}

\begin{definition}[Класс $\SSE$]\label{def:SSE}
Для $S \in \RS$ класс $\SSE \subseteq \Ob \cE$ \emph{$S$-определимых заострённых актов} — замыкание базового класса
\[
  \SSEbase = \bigl\{ (F', \cM, \iota_\cM, r_\cM) : \cM \models S,\ F' \in \cF,\ \iota_\cM := \ev_\cM,\ r_\cM := (\ev_\cM)^L \bigr\},
\]
где $\ev_\cM : \Syn(F') \to \cM$ — синтаксико-семантическое вычисление из (R5b), а $(\ev_\cM)^L$ — его левый сопряжённый (Ламбек--Скотт), относительно покомпонентного применения операций Определения~\ref{def:SS}: $F$-компонента — по Определению~\ref{def:SS}; $\cC$-компонента — через забывающий $\mathbf{StrCat}_{S, n} \to (\infty, n)\text{-}\Cat$; $(\iota, r)$-компоненты — по функториальности сопряжений под $S$-сохраняющими $(\infty, n)$-функторами. Под-класс $\SSEglobal \subseteq \SSE$ — соответствующее замыкание $\SSEbase$ только теми операциями Определения~\ref{def:SS}, что определяют $\cS_S^{\mathrm{global}}$ (Кан-расширения, конструкции Гротендика, производные функторы).
\end{definition}

\begin{lemma}[Акт-синтаксико-семантическая лемма]\label{lem:SS-membership-E}
Пусть $(F, \cC, \iota, r) \in \cE$ — заострённый рефлексивный акт, четвёрка которого формально $S$-определима, то есть существуют формулы $\phi_F, \psi_\cC, \chi_\iota, \chi_r$ в некотором $F' \in \cF$ с $F' \hookrightarrow S$, специфицирующие соответственно $F, \cC, \iota, r$, причём треугольные тождества и рефлексивность $\iota$ доказуемы в $F'$. Тогда $(F, \cC, \iota, r) \in \SSEglobal$.
\end{lemma}

\begin{proof}
По условию $F$ формально $S$-определим через $\phi_F$; Лемма~\ref{lem:SS-membership} даёт $F \in \cS_S^{\mathrm{global}}$. Модельная категория $\cC$ специфицирована $\psi_\cC$ внутри $F'$; как $S$-структурированная $(\infty, n)$-категория, подлежащая $(\infty, n)$-категория $\cC$ лежит в $\cS_S^{\mathrm{global}}$ по тому же аргументу, применённому к $\cS_S^{\mathrm{global}}$-замыканию Определения~\ref{def:SS}. Пойнтинг $\iota$ и рефлектор $r$, специфицированные $\chi_\iota$ и $\chi_r$ в $F'$, — морфизмы между $\cS_S^{\mathrm{global}}$-объектами, определимые в $F' \hookrightarrow S$. Сопряжённая пара $(\iota, r)$ реализуется в $\SSEglobal$ следующей явной конструкцией: $\iota$ получается как Кан-расширение (Определение~\ref{def:SS}, O1) вдоль вложения $\Syn(F') \hookrightarrow \cC$, специфицированного $\chi_\iota$; существование \emph{левого} рефлектора $r$ (левого сопряжённого к $\iota$ по Определению~\ref{def:enactments}) следует затем из теоремы существования сопряжённого функтора для доступных функторов между локально $\kappa_S$-представимыми категориями (Адамек--Росицкий~\cite{AdamekRosicky1994}, теорема существования сопряжённого функтора в кополном случае), с переносом треугольных тождеств из $\chi_r$. Собранная покомпонентно, четвёрка $(F, \cC, \iota, r) \in \SSEglobal$.
\end{proof}

\subsection{AC/OC Морита-двойственность}

\begin{theorem}[AC/OC Морита-двойственность --- формально-организационная форма]\label{thm:ac-oc-duality}

Сопряжённая пара $\varepsilon \dashv \alpha$ Определения~\ref{def:enactments} индуцирует $(\infty, \infty)$-эквивалентность классифицирующих $2$-стеков
\[
  \fM \;\simeq_{(\infty, \infty)}\; \fM_\cE, \qquad \fM_\cE := \cE / \mathrm{gauge}_\cE,
\]
где левый стек модулей — классифицирующий $2$-стек $\cF$ из Соглашения~\ref{conv:notation}, а $\fM_\cE$ — соответствующий классифицирующий $2$-стек $\cE$. Конкретно:
\begin{enumerate}[label=(\alph*)]
\item $\varepsilon : \cF \to \cE$ вполне унивалентен;
\item единица $\id_\cF \Rightarrow \alpha \circ \varepsilon$ — равенство; коединица $\varepsilon \circ \alpha \Rightarrow \id_\cE$ — калибровочная эквивалентность, то есть эквивалентность после перехода к $\fM_\cE$;
\item страты $\LFnd, \LCls, \LClsMax, \LAbs$ стека $\fM$ биективно соответствуют стратам $\LFnd^\cE, \LCls^\cE, (\LClsMax)^\cE, \LAbs^\cE$ стека $\fM_\cE$ через индуцированную эквивалентность с сохранением мета-операций $\mathrm{Cls}$ и $\mathrm{Gen}$.
\end{enumerate}
\end{theorem}

\begin{proof}
\emph{(a) Полная унивалентность $\varepsilon$.} Пусть $F_1, F_2 \in \cF$. Развёртывая Определение~\ref{def:enactments}, $1$-морфизм $\varepsilon(F_1) \to \varepsilon(F_2)$ — это пара $(\phi, \Phi)$ с $\phi : F_1 \to F_2$ в $\cF$ и $\Phi : \Syn(F_1) \to \Syn(F_2)$ — $S$-сохраняющим $(\infty, n)$-функтором с обратимой $2$-клеткой $\Phi \Rightarrow \Syn(\phi)$. \emph{Классифицирующее свойство} синтаксических категорий (функториальная семантика: $S$-сохраняющие $(\infty, n)$-функторы $\Syn(F_1) \to \cC$ естественно соответствуют $F_1$-модельным структурам на $\cC$; Ламбек--Скотт~\cite{LambekScott1986} для $n = 1$, Капулкин--Ламсдейн~\cite{KapulkinLumsdaine2021} для теоретико-типового $(\infty, 1)$-случая), применённое с $\cC = \Syn(F_2)$, отождествляет $S$-сохраняющие функторы $\Syn(F_1) \to \Syn(F_2)$ с $F_1$-моделями в $\Syn(F_2)$, то есть с интерпретациями $\phi : F_1 \to F_2$: каждый такой функтор возникает, с точностью до обратимой $2$-клетки, как $\Syn(\phi)$ для единственного $\phi$. Следовательно, забывающее отображение hom-$2$-категорий
\[
  \Hom_{\cE}(\varepsilon F_1, \varepsilon F_2) \longrightarrow \Hom_{\cF}(F_1, F_2), \qquad (\phi, \Phi) \mapsto \phi,
\]
есть $(\infty, 2)$-эквивалентность. Это устанавливает полную унивалентность $\varepsilon$ на каждой hom-$2$-категории.

\emph{(b) Существенная сюръективность на калибровочном уровне.} Пусть $(F, \cC, \iota, r) \in \cE$. Рефлексионный морфизм $(\id_F, r) : (F, \cC, \iota, r) \to \varepsilon(F) = (F, \Syn(F), \id, \id)$ --- структурная клетка которого — обратимая коединица $\varepsilon_r : r \circ \iota \Rightarrow \id_{\Syn(F)}$ (левый сопряжённый к вполне унивалентному функтору; Риль--Верити~\cite{RiehlVerity2022}, предл.~2.1.11) --- лежит в классе $\mathcal{R}_\cE$ Определения~\ref{def:gauge-E} (здесь $r_2 \circ \Phi \circ \iota_1 = \id \circ r \circ \iota \simeq \id_{\Syn(F)}$). По построению $\fM_\cE = \cE[\mathcal{R}_\cE^{-1}]$ этот морфизм становится обратимым в $\fM_\cE$, давая каноническую gauge$_\cE$-эквивалентность
\[
  (F, \cC, \iota, r) \;\simeq_{\mathrm{gauge}_\cE}\; \varepsilon(F) \quad\text{в } \fM_\cE.
\]
Подчеркнём, что \emph{не} утверждается: единица $\eta_r : \id_\cC \Rightarrow \iota \circ r$ не обратима в $\cE$, если только $\iota$ не существенно сюръективен; эквивалентность выше создаётся локализацией — ровно так же, как кольцевая Морита-эквивалентность создаётся обращением бимодульных эквивалентностей. Единственность $r$ с точностью до единственной обратимой $2$-клетки при данном $\iota$ делает эквивалентность естественной по объекту. Следовательно, $\varepsilon$ существенно сюръективен после перехода к калибровочным классам, и теорема обретает точное содержание: \emph{модули актов, локализованные по рефлексиям, эквивалентны модулям теорий}.

В сочетании с (a) пара $(\varepsilon, \alpha)$ индуцирует эквивалентность $\fM \simeq \fM_\cE$ на $(\infty, 2)$-уровне.

\emph{Поднятие до $(\infty, \infty)$.} Аргумент выше параметричен по $n \in \mathbb{N} \cup \{\infty\}$: он использует лишь (R5b), $2$-функториальность $\Syn$ по $n$ и единственность рефлекторов на $(\infty, n)$. Первое — это (R5b); второе — Капулкин--Ламсдейн~\cite{KapulkinLumsdaine2021} вместе с Барвиком--Шоммер-Присом~\cite{BarwickSchommerPries2021}; третье — Адамек--Росицкий~\cite{AdamekRosicky1994}, теорема~1.26, применённая по срезам через Люри HTT~\cite{LurieHTT} \S 5.4 (сохранение $\lambda$-фильтрованных копределов доступными сопряжениями). Стабилизация эквивалентности при $n = \infty$ следует из Теоремы~\ref{thm:bergner-lurie-stab}.

\emph{(c) Соответствие стратов.} Классификационная мета-операция $\mathrm{Cls}$ на $\fM$ переводит $F \in \LFnd$ в его класс классифицируемых под-стеков. При эквивалентности $\fM \simeq \fM_\cE$ акт $\varepsilon(F) = (F, \Syn(F), \id, \id)$ наследует классификационную структуру через $\iota = \id$, то есть $\mathrm{Cls}^\cE(\varepsilon F)$ совпадает с $\varepsilon$, применённым к $\mathrm{Cls}(F)$. Аналогично для $\mathrm{Gen}$. Страты $\LClsE \supsetneq \LClsMaxE \supseteq \LAbsE$ на $\fM_\cE$ определяются $\cE$-аналогами Определения~\ref{def:hierarchy}: калибровочный класс $[(F, \cC, \iota, r)] \in \LClsE$ тогда и только тогда, когда $[F] \in \LCls$; аналогично для более тонких стратов. При этой трансляции эквивалентность из (a)--(b) ограничивается до биекции стратов.
\end{proof}

\begin{corollary}[Консервативность; сила непротиворечивости]\label{cor:ac-oc-conservativity}
Расширение рамки MSFS категорией $\cE$ и структурными отображениями $(\alpha, \varepsilon)$ — консервативное расширение $\cF$-основанной рамки: каждое утверждение о $\fM_\cE$ эквивалентно, через Теорему~\ref{thm:ac-oc-duality}, утверждению о $\fM$. В частности, объемлющая метатеория, необходимая для построения $\cE$ и рассуждений о ней, та же, что и для $\cF$:
\[
  \Con(\text{MSFS с } \cE) \;=\; \Con(\text{MSFS}) \;=\; \Con(\ZFC + 2\text{-inacc}).
\]
\end{corollary}

\begin{proof}
По Теореме~\ref{thm:ac-oc-duality} каждое утверждение о $\fM_\cE$ переводится в эквивалентное утверждение о $\fM$ через эквивалентность. Новые аксиомы не вводятся; по построению $\cE$ $\mathbf{U}_2$-мала и строится внутри $\ZFC + 2\text{-inacc}$ из $\cF$ и $\mathbf{StrCat}_{S, n}$ (Определение~\ref{def:strcat}), используя лишь сопряжения Ламбека--Скотта, гарантированные (R5b), и рефлекторы, данные как объектные данные (Определение~\ref{def:enactments}).
\end{proof}

Эквивалентность Теоремы~\ref{thm:ac-oc-duality} переносит условия $\LAbs$ из \S\ref{sec:conditions} на акт-сторону. Неформально: \emph{акт-абсолютная} тройка $(F, \cC, \iota)$ — та, чьё акт-стороннее описание одновременно допускает формальную спецификацию, не допускает нетривиальной координации с другими актами $S$-определимого класса и максимально восприимчиво (всякий другой акт факторизуется через него).

\begin{definition}[Условия $\LAbsE$]\label{def:level6-cE}
Для $S \in \RS$ и $n \in \mathbb{N} \cup \{\infty\}$ заострённый рефлексивный акт $(F, \cC, \iota, r) \in \cE$ \textbf{акт-абсолютен}, если он удовлетворяет тройке:
\begin{itemize}
\item[$(\tilde{F}_S)$] \emph{Акт-определимость}: четвёрка $(F, \cC, \iota, r)$ формально $S$-определима в смысле Леммы~\ref{lem:SS-membership-E}, то есть специфицирована формулами $(\phi_F, \psi_\cC, \chi_\iota, \chi_r)$ в некотором $F' \in \cF$ с $F' \hookrightarrow S$ (это уточняет $(F_S)$, дополнительно специфицируя модельно-категорное представление, пойнтинг и рефлектор).
\item[$(\tilde{\Pi}_{4, S, n})$] \emph{Некоординируемость}: не существует $1$-морфизма $(\phi, \Phi) : (F, \cC, \iota, r) \to (F', \cC', \iota', r')$ в $\cE$ с целевой четвёркой из $\SSE$ (Определение~\ref{def:SSE}), в котором $\Phi$ — $(\infty, n)$-эквивалентность на свой образ.
\item[$(\tilde{\Pi}_{3\text{-max}, S, n})$] \emph{Максимальная восприимчивость}: для каждого акта $(F', \cC', \iota', r') \in \cE$ существует $1$-морфизм $(\phi, \Phi) : (F', \cC', \iota', r') \to (F, \cC, \iota, r)$, реализующий Морита-редукцию на $(\infty, n)$-уровне.
\end{itemize}
\emph{Страта $\LAbsE$} — класс $\mathrm{gauge}_\cE$-классов эквивалентности в $\fM_\cE$, имеющих акт-абсолютного представителя.
\end{definition}

\subsection{Дуальная граничная лемма}

\begin{theorem}[Дуальная граничная лемма; $(\alpha)$-часть Dual-AFN-T]\label{thm:dual-afnt}
Для каждой Rich-метатеории $S \in \RS$ и каждого категорного уровня $n \in \mathbb{N} \cup \{\infty\}$ ни один заострённый рефлексивный акт $(F, \cC, \iota, r) \in \cE$ не удовлетворяет тройке $\LAbsE$ $(\tilde{F}_S) \wedge (\tilde{\Pi}_{4, S, n}) \wedge (\tilde{\Pi}_{3\text{-max}, S, n})$. Уже под-пара $(\tilde{F}_S) \wedge (\tilde{\Pi}_{4, S, n})$ структурно самопротиворечива.
\end{theorem}

\begin{proof}
Пусть $(F, \cC, \iota, r) \in \cE$ удовлетворяет $(\tilde{F}_S) \wedge (\tilde{\Pi}_{4, S, n})$. По Лемме~\ref{lem:SS-membership-E}, $(F, \cC, \iota, r) \in \SSEglobal \subseteq \SSE$. Тождественный $1$-морфизм $(\id_F, \id_\cC)$ имеет $\Phi = \id_\cC$ — $(\infty, n)$-эквивалентность на свой образ с целью в $\SSE$. Такая координация запрещена $(\tilde{\Pi}_{4, S, n})$. Противоречие. Следовательно, $\LAbsE = \emptyset$.
\end{proof}

\begin{corollary}[Пустота $\LAbsE$]\label{cor:level6-empty-cE}
$\LAbsE = \emptyset$.
\end{corollary}

\begin{theorem}[Дуальная пятиосевая абсолютность]\label{thm:dual-five-axis}
Теоремы~\ref{thm:horizontal}--\ref{thm:completeness} переносятся на $\LAbs^\cE$: пустота $\LAbs^\cE$ абсолютна вдоль пяти структурных осей
\begin{center}
\emph{горизонтальной} (по $\RS$), \quad \emph{вертикальной} (по $n$), \quad \emph{мета-вертикальной}, \quad \emph{латеральной}, \quad \emph{полноты} (внутри Lawvere-скоупа).
\end{center}
В частности, три стандартных обходных пути (Теоремы~\ref{thm:universe}, \ref{thm:reflective}, \ref{thm:slice-locality}) закрывают свои $\cE$-аналоги.
\end{theorem}

\begin{proof}
Каждая ось переносится через Теорему~\ref{thm:ac-oc-duality} покомпонентно: пять механизмов Теорем~\ref{thm:horizontal}--\ref{thm:completeness} действуют параметрически на $F$-компоненте или естественно на $\cC$-компоненте четвёрки $(F, \cC, \iota, r)$, и пара $(\alpha, \varepsilon)$ их сохраняет. Lawvere-скоуп $\mathrm{LS}(\cE)$ определяется как акты с $F \in \mathrm{LS}(\cF)$ и $\cC$, замкнутой симметрично-моноидальной (картезианово замкнутой или $*$-автономной — покрывая линейную логику и квантово-категорные акты); обобщённая диагональ Яновского~\cite{Yanofsky2003} применяется в $\cC$ и переносится на $\Syn(F)$ через $\iota \dashv r$. Три обходных пути переносятся аналогично покомпонентным действием $(\alpha, \varepsilon)$.
\end{proof}

\section{Место в серии no-go}\label{sec:series}

AFN-T занимает позицию в структурной серии фундаментальных no-go-теорем. Все такие результаты инстанциируют одну категорную схему: \emph{система, достаточно выразительная, чтобы сформулировать кандидата в максимальный объект, не может содержать этот объект без противоречия}, — реализуемую самореферентным вложением, превращающим мнимый максимальный объект в элемент того самого класса, который он должен был превзойти. Каждый специализирует схему своим аспектом максимальности и диагональным шагом:

\begin{center}
\begin{tabular}{@{}l l p{3.2cm} p{5.7cm}@{}}
\toprule
Теорема & Год & Аспект максимальности & Диагональный шаг \\ \midrule
Кантор~\cite{Cantor1891} & 1891 & кардинальный & $x \notin f(x)$ для $f : X \to \cP(X)$ \\
Рассел~\cite{Russell1903} & 1903 & свёртывание & $R = \{x : x \notin x\}$, $R \in R \iff R \notin R$ \\
Гёдель (I, II)~\cite{Godel1931} & 1931 & полнота / непротиворечивость & $G = $ «$G$ недоказуемо» \\
Тарский~\cite{Tarski1936} & 1936 & семантическая истина & «это предложение не истинно» \\
Ловер FP~\cite{Lawvere1969} & 1969 & категорная неподвижная точка & картезианова замкнутость + слабая сюръективность $\Rightarrow$ у каждого $f$ есть неподвижная точка \\
Эрнст~\cite{Ernst2015} & 2015 & неограниченная теория категорий & диагональ Ловера при $n = 1$ под Феферманом (R1)--(R3) \\
AFN-T & --- & фундаментальный (определимость + нередуцируемость + порождаемость) & синтаксико-семантический мост: $(F_S) \Rightarrow X \in \cS_S^{\mathrm{global}}$ по Лемме~\ref{lem:SS-membership}, противоречие с $(\Pi_4)$ через $\id_X$ (Теорема~\ref{thm:afnt-alpha}) \\
Dual-AFN-T & --- & акт-абсолютный (определимость + некоординируемость + восприимчивость) & тождественный акт $(\id_F, \id_\cC)$ с целью в $\SSE$ (Лемма~\ref{lem:SS-membership-E}) противоречит $(\tilde{\Pi}_4)$ (Теорема~\ref{thm:dual-afnt}); эквивалентно через $\alpha$ + AFN-T \\
\bottomrule
\end{tabular}
\end{center}

\subsection{Подчинение}

\begin{theorem}[Подчинение через специализацию (уровень рамки)]\label{thm:subsumption}
Для каждой классической no-go-теоремы $T$ из списка ниже существует тройка $(S_T, n_T, X_T)$ с $S_T \in \RS$, $n_T \in \mathbb{N} \cup \{\infty\}$ и $X_T$ — мнимым максимальным объектом $T$, — такая что (i) $X_T \models (F_{S_T})$; (ii) $X_T$ при гипотезе $T$ удовлетворял бы $(\Pi_{4, S_T, n_T})$; (iii) Теорема~\ref{thm:afnt-alpha}, применённая в $(S_T, n_T)$, даёт структурное противоречие с существованием $X_T$ как $\LAbs$-кандидата. Теорема даёт \emph{рамку}, в которой обитает каждый классический no-go; построчные диагонали (канторовская $|X|<|\cP(X)|$, расселовская $R \in R \leftrightarrow R \notin R$ и т.\,д.) — независимые классические доказательства, уточняющие абстрактный $\id_{X_T}$-шаблон до конкретного свидетеля, используемого $T$, и они не выводятся из одной AFN-T~$\alpha$.
\begin{center}\footnotesize
\begin{tabular}{@{}lllcl@{}}
\toprule
$T$ & $S_T$ & $X_T$ & $n_T$ & $(F_{S_T})$-свидетель \\ \midrule
Кантор~\cite{Cantor1891} & $\ZFC$ & универсальный $\cP : \mathbf{Set} \to \mathbf{Set}$ & $1$ & $\phi_{\cP}(X,Y) \equiv Y = \cP(X)$ \\
Рассел~\cite{Russell1903} & $\ZFC$ & множество $\{x : x \notin x\}$ & $1$ & $\phi_R(x) \equiv x \notin x$ \\
Гёдель I~\cite{Godel1931} & $\mathsf{PA}$ & $G$ (недоказуемое предложение) & $1$ & $\phi_G \equiv \neg \mathrm{Prov}(\ulcorner G \urcorner)$ \\
Гёдель II~\cite{Godel1931} & $\mathsf{PA}$ & внутреннее доказательство $\mathrm{Con}(\mathsf{PA})$ & $1$ & $\phi_{\mathrm{Con}}$ через $\mathrm{Prov}$ \\
Тарский~\cite{Tarski1936} & $\mathsf{PA}$ & арифметический предикат истины & $1$ & $T(\phi) \equiv \phi$ \\
Ловер FP~\cite{Lawvere1969} & CCC-расш.\ $\ZFC$ & сл.\ сюръ.\ $X$, $f:X \to X^X$ & $1$ & $\phi_X$ в CCC-синтаксисе \\
Эрнст~\cite{Ernst2015} & Феферман (R1)--(R3) & UCT-объект & $1$ & неогр.\ схема Фефермана \\
\bottomrule
\end{tabular}
\end{center}
В каждой строке AFN-T в $(S_T, n_T)$ даёт $\LAbs(S_T, n_T) = \emptyset$ в аспекте максимальности строки; классическая диагональ $T$ — конкретное свидетельское уточнение, устанавливающее $T$ внутри рамки.
\end{theorem}

\begin{proof}
$X_T$ каждой строки $S_T$-определим перечисленной формулой (i), так что $X_T \in \cS_{S_T}^{\mathrm{global}}$ по Лемме~\ref{lem:SS-membership}. $\id_{X_T} : X_T \to X_T$ нарушает $(\Pi_{4, S_T, n_T})$ согласно Теореме~\ref{thm:afnt-alpha}, так что $X_T \notin \LAbs$. Построчная диагональ реконструирует: Кантор — через $|X| < |\cP(X)|$, уточняющее $\id_{\cP}$; Рассел — через $R \in R \leftrightarrow R \notin R$, уточняющее $\id_R$; Гёдель — через $G$-предложение $\leftrightarrow \neg\mathrm{Prov}(\ulcorner G \urcorner)$; лжец Тарского; формула неподвижной точки Ловера $f(x_0) = x_0$ через $\lambda x. f(x)(x)$; Эрнст — через неограниченную аксиомную схему Фефермана. Каждое уточнение специализирует универсальный $\id_{X_T}$-шаблон Теоремы~\ref{thm:afnt-alpha} до конкретного формального свидетеля строки.
\end{proof}

\section{Следствия и открытые вопросы}\label{sec:consequences}

\subsection{Следствия для программ оснований}

AFN-T налагает равномерную структурную границу на каждую действующую программу оснований: каждая обитает в $\LFnd$ или в $\LCls$ (возможно, с членством в максимальном подклассе $\LClsMax \subseteq \LCls$), при пустом $\LAbs$. Конкретно:
\begin{itemize}\setlength{\itemsep}{1pt}
\item \textbf{Унивалентные основания} (Воеводский~\cite{VoevodskyFoundations2010}; HoTT Book~\cite{HoTTBook2013}) — частичная $\LCls$-мета-рамка (Теорема~\ref{thm:meta-mult}), классифицирующая $\HoTT$-расширения; возвысить её до $\LAbs$ нельзя.
\item \textbf{Теория высших топосов} (Люри~\cite{LurieHTT}) и \textbf{когезивная рамка} (Шрайбер~\cite{Schreiber2013}) занимают соответственно уровни 5 и 5+.
\item \textbf{Конденсированная математика} (Клаузен--Шольце~\cite{ClausenScholze2019}), \textbf{перфектоидная геометрия} и \textbf{$\infty$-космои} (Риль--Верити) ограничены аналогично; эскалация в $\LAbs$ невозможна.
\item \textbf{Пруф-ассистенты} (Coq, Lean, Agda) реализуют $\LFnd$-основания ($\CIC$, $\CIC + \mathrm{Univalence}$, $\MLTT$); множественность таких систем — множественность $\LFnd$.
\item \textbf{Категорная логика}: классификация систем оснований $(\infty, n)$-категориями и их мета-структурами имеет верхнюю границу в $\LCls$, в согласии с унитарностью Барвика--Шоммер-Приса~\cite{BarwickSchommerPries2021}.
\end{itemize}

\paragraph{Диагностика: какое из условий $\LAbs$ нарушает конкретный кандидат?} Для каждого кандидата $X$, выдвигавшегося в литературе на роль абсолютного мета-основания, точный способ провала внутри формальной рамки AFN-T устанавливается проверкой, какой член тройки $(F_S) \wedge (\Pi_{4, S, n}) \wedge (\Pi_{3\text{-max}, S, n})$ нарушен.
\begin{itemize}\setlength{\itemsep}{1pt}
\item \emph{Унивалентные основания (UF)}: проходят $(F_S)$ (формально определимы как HoTT + схема аксиомы унивалентности), условно проходят $(\Pi_{4, S, n})$ (не Морита-редуцируемы к $\ZFC$-внутренним структурам в прямом направлении), но \emph{нарушают $(\Pi_{3\text{-max}, S, n})$}: UF не классифицируют Rich-основания вне HoTT-смежного семейства (например, не классифицируют Морита-редукцией ZFC-специфическую кардинальную арифметику или некоммутативно-геометрические спектральные тройки). Значит, UF $\in \LCls \setminus \LClsMax$.
\item \emph{Теория высших топосов (Люри)}: проходит $(F_S)$, нарушает $(\Pi_{3\text{-max}, S, n})$: классифицирует категорные структуры HoTT-совместимой формы, но не допускает классификационного функтора на Rich-основания вне этого класса (например, на первопорядковые теории множеств без HoTT-стильного внутреннего языка).
\item \emph{Когезивная рамка (Шрайбер)}: проходит $(F_S)$, нарушает $(\Pi_{3\text{-max}})$; вдобавок аксиомы когезии специализируются на геометрические контексты, дополнительно сужая классификационный образ.
\item \emph{$\infty$-космои (Риль--Верити)}: проходят $(F_S)$ (аксиоматизированы в $\ZFC + 2\text{-inacc}$), нарушают $(\Pi_{3\text{-max}})$; они классифицируют $(\infty, 1)$-категорные структуры, а не весь $\fM$.
\item \emph{Любой гипотетический кандидат $X$, удовлетворяющий $(F_S) \wedge (\Pi_{3\text{-max}})$}: по Теореме~\ref{thm:afnt-alpha} $X$ тогда нарушает $(\Pi_{4, S, n})$: будь $X$ одновременно формально определим (через $(F_S)$) и максимально порождающ (через $(\Pi_{3\text{-max}})$), то по Лемме~\ref{lem:SS-membership} $X \in \cS_S^{\mathrm{global}}$, и $\id_X$ реализовал бы Морита-редукцию к $\cS_S$-цели, противореча $(\Pi_{4, S, n})$.
\end{itemize}
Итак, три условия не могут выполняться одновременно: нынешние кандидаты нарушают $(\Pi_{3\text{-max}})$; любой гипотетический более сильный кандидат с $(F_S) \wedge (\Pi_{3\text{-max}})$ нарушил бы $(\Pi_{4, S, n})$. Диагностика указывает точное препятствие в каждом случае, а не просто провозглашает невозможность.

\subsection{Открытые вопросы}

\paragraph{(Q1) Существование максимальной мета-рамки.} Теорема~\ref{thm:meta-cat} утверждает категоричность $\Meta_{\mathrm{Cls}}^{\top}$ условно на его непустоту. Существует ли конкретная $\mathbf{F} \in \Meta_{\mathrm{Cls}}^{\top}$, конструируемая в $\ZFC + 2\text{-inacc}$? Положительный ответ отождествил бы каноническую $\LCls$-мета-рамку с точностью до $(\infty, \infty)$-эквивалентности. Построение такого свидетеля лежит вне рамок настоящего препринта; кандидатная рамка — предмет отдельной продолжающейся работы. Непустота $\Meta_{\mathrm{Cls}}^{\top}$ — экзистенциальный вопрос с \emph{четырьмя независимыми структурными препятствиями}, отвечающими четырём условиям максимальности: (Max-1) требует классификатора с классификацией, существенно сюръективной на весь $\fM$; (Max-2) требует транзитивного $\mathrm{Aut}_2$-действия на калибровочных классах; (Max-3) требует $\omega$-глубинной предикативной фильтрации с условием на схему свёртывания; (Max-4) требует канонического интенсионального пополнения $\II_{\mathbf{F}} : \mathbf{F}^{\mathrm{op}} \to \Sint$ с (I-1)--(I-4). Из них (Max-4) — отдельное экзистенциальное препятствие, отличное от (Max-1)--(Max-3): построение $\II_\mathbf{F}$ как канонического пополнения (не просто существующего, а канонически определяемого структурой классификатора) нетривиально даже при данности первых трёх условий. Решение Q1, соответственно, сводится к независимому построению каждой из четырёх структурных данностей на общем подлежащем классификаторе, а не к одному достаточному условию.

\paragraph{(Q2) Полнота списка мета-рамок.} Насколько полон список $\{\mathbf{F}_{\mathrm{univalent}},\allowbreak \mathbf{F}_{\mathrm{cosmoi}},\allowbreak \mathbf{F}_{\mathrm{cohesive}},\allowbreak \mathbf{F}_{\mathrm{HA}},\allowbreak \mathbf{F}_{\mathrm{synthetic}}\}$ идентифицированных на сегодня $\LCls$-мета-структур ($\mathbf{F}_{\mathrm{HA}}$ — программа Higher Algebra Люри~\cite{LurieHA})? Существуют ли дальнейшие естественные мета-рамки?

\paragraph{(Q3) Исчерпываемость обходных путей.} Теоремы~\ref{thm:universe}, \ref{thm:reflective}, \ref{thm:I-existence}, \ref{thm:slice-locality} закрывают три стандартных обходных пути, а Приложение~\ref{app:paraconsistent} разбирает паранепротиворечивый случай. Кандидаты в дополнительные независимые пути: модальные или темпоральные логические расширения, субструктурные логики без экспоненциала $!$, реализуемость и чисто конструктивные основания, основания на неклассических моделях вычислимости. Каждый из них, по-видимому, редуцируется к одному из закрытых путей внутри R-S; формальная теорема исчерпываемости здесь не даётся и остаётся открытой.

\paragraph{(Q4) Минимальность силы непротиворечивости.} Доказательство использует $\ZFC + 2\text{-inacc}$ (Соглашение~\ref{conv:zfc-inacc}). Достаточно ли более слабой силы непротиворечивости для полной теоремы, или мета-классификационная стабилизация (Теорема~\ref{thm:meta-stab}) подлинно требует второго недостижимого?

\paragraph{(Q5) Под-стек слабых Rich-метатеорий.} Образуют ли ограниченно-арифметические теории ($\mathsf{I\Delta_0}$, $\mathsf{S}_2^i$), удовлетворяющие (R2)--(R5) и ограниченному (R1), отдельный страт $\LFnd^{\mathrm{weak}} \subset \LFnd$, или они Морита-плотны в $\LFnd$ средствами обратной математики~\cite{Simpson2009}?

\appendix

\section{Категорные предварительные сведения}\label{app:categorical}

Здесь я собираю категорные инструменты и стандартные теоремы, привлекаемые в доказательствах выше.

\subsection{\texorpdfstring{$2$-категории}{2-категории}}

\textbf{$2$-категория} $\cC$ состоит из:
\begin{itemize}
\item Класса объектов $\Ob(\cC)$.
\item Для каждой пары $A, B \in \Ob(\cC)$ — малой категории $\Hom_\cC(A, B)$ из $1$-морфизмов и $2$-морфизмов.
\item Функторов композиции и тождественных $1$-морфизмов, удовлетворяющих стандартной $2$-категорной когерентности.
\end{itemize}

$2$-категория \textbf{локально мала}, если каждая $\Hom_\cC(A, B)$ — малая категория. Все $2$-категории в этой статье локально малы, если не оговорено иное.

Ссылки: Келли~\cite{Kelly1982} — обогащённые и $2$-категорные основания; Бенабу~\cite{Benabou1967} — бикатегории.

\subsection{\texorpdfstring{$(\infty, n)$-категории}{(infinity,n)-категории}}

$(\infty, n)$-категория — обобщение теории категорий, в котором:
\begin{itemize}
\item Морфизмы имеют $k$-морфизмы для $k \leq n$.
\item $k$-морфизмы при $k > n$ все обратимы (эквивалентности).
\end{itemize}

Модели: квазикатегории, $n$-категории Сегала, полные пространства Сегала, $\Theta_n$-пространства. Эквивалентность этих моделей установлена Ара--Мальциниотисом~\cite{AraMaltsiniotis2017}.

$(\infty, n)$-иерархия стабилизируется при $n = \infty$: каноническое вложение $(\infty, \infty)\text{-}\Cat \hookrightarrow (\infty, \infty + 1)\text{-}\Cat$ — эквивалентность (Барвик--Шоммер-Прис~\cite{BarwickSchommerPries2021}, итеративно поверх сравнения моделей Бергнер--Резка~\cite{BergnerRezk2013}).

Ссылки: Люри HTT~\cite{LurieHTT} — $(\infty, 1)$-категории; Риль--Верити~\cite{RiehlVerity2022} — синтетическая теория $(\infty, 1)$-категорий.

\subsection{Доступные функторы}

Функтор $F : \cC \to \cD$ между локально малыми категориями \textbf{$\lambda$-доступен} (для регулярного кардинала $\lambda$), если $F$ сохраняет $\lambda$-фильтрованные копределы.

Категория $\cC$ \textbf{локально $\lambda$-представима}, если она кополна и имеет малую плотную подкатегорию $\lambda$-представимых объектов.

Ссылка: Адамек--Росицкий~\cite{AdamekRosicky1994} — теория доступности.

\begin{lemma}[Представление Лэра--Маккаи--Паре]\label{lem:lair}
Каждая доступная категория $2$-категорно эквивалентна $2$-категории моделей малого скетча; эквивалентно, $2$-категории моделей доступной $2$-теории.
\end{lemma}

\noindent\emph{Ссылки для доказательства.} Лэр~\cite{Lair1981} ($1$-категорные скетчи, оригинальный французский источник); Маккаи--Паре~\cite{MakkaiPare1989}, глава~3 (систематическое развитие, включая $2$-категорное соответствие скетч/теория); Адамек--Росицкий~\cite{AdamekRosicky1994}, глава~2 (трактовка через доступность). Лемма используется в статье только через эти ссылки.

\subsection{Категории display-отображений}

\begin{definition}[Категория display-отображений, $1$-категорная версия]\label{def:display-1cat}
\textbf{Категория display-отображений} — пара $(\cC, \cD)$, где $\cC$ — категория с пулбэками, а $\cD \subseteq \Mor(\cC)$ — класс морфизмов, устойчивый под пулбэками и содержащий изоморфизмы.
\end{definition}

Джейкобс~\cite{Jacobs1999} устанавливает, что категории display-отображений дают категорную семантику интенсиональных теорий типов, где display-морфизмы соответствуют семействам зависимых типов.

Гамбино--Гарнер~\cite{GambinoGarner2008} распространяет это на $2$-категорный сеттинг, что даёт Определения~\ref{def:display-class} и~\ref{def:display-2-cat}.

\subsection{Бикатегория дробей}

\begin{lemma}[Бикатегория дробей Пронк]\label{lem:pronk}
Даны $2$-категория $\cC$ и класс $\mathcal{W}$ $1$-морфизмов, удовлетворяющий условиям исчисления дробей Бусфилда--Пронк (BF1)--(BF5); тогда $2$-локализация $\cC[\mathcal{W}^{-1}]$ существует как бикатегория дробей.
\end{lemma}

Ссылка: Пронк~\cite{Pronk1996}.

Условия (BF1)--(BF5):
\begin{itemize}
\item[\textbf{(BF1)}] Содержит тождества.
\item[\textbf{(BF2)}] Замкнут под композицией.
\item[\textbf{(BF3)}] $2$-пулбэк-устойчив.
\item[\textbf{(BF4)}] $2$-право-фильтрован (любой морфизм продолжаем).
\item[\textbf{(BF5)}] $\mathcal{W}$-насыщенность: содержит все $2$-эквивалентности.
\end{itemize}

\subsection{Теорема Ловера о неподвижной точке}

\begin{lemma}[Неподвижная точка Ловера, $2$-категорная]\label{lem:lawvere-2}
В $2$-картезианово замкнутой категории $\cC$ со слабо точечно-сюръективным $1$-морфизмом $e : A \to Y^A$ каждый $f : Y \to Y$ имеет неподвижную точку: существует $y_0 : 1 \to Y$ с $f \circ y_0 = y_0$ с точностью до когерентной $2$-эквивалентности.
\end{lemma}

Ссылки: Ловер~\cite{Lawvere1969}; Яновский~\cite{Yanofsky2003} — $2$-категорная формулировка.

\begin{lemma}[Неподвижная точка Ловера, $(\infty, \infty)$-категорная]\label{lem:lawvere-inf}
Пусть $\cC$ — $(\infty, \infty)$-картезианово замкнутая категория, допускающая объект $A$ и $1$-морфизм $e : A \to Y^A$, \emph{слабо точечно-сюръективный с точностью до когерентной гомотопии} (каждая глобальная точка $1 \to Y^A$ факторизуется через $e$ с точностью до обратимых высших клеток). Тогда каждый $f : Y \to Y$ имеет $(\infty, \infty)$-когерентную неподвижную точку.
\end{lemma}

Это получается распространением Леммы~\ref{lem:lawvere-2} через $\Theta_n$-пространственную модель Резка~\cite{Rezk2010}, переданную в предел $n \to \infty$ теоремой унитарности Барвика--Шоммер-Приса~\cite{BarwickSchommerPries2021}. $(\infty, \infty)$-когерентная неподвижная точка — канонический предел совместимого семейства неподвижных точек, построенных на каждом конечном $n$.

\subsection{Критерий типа Уайтхеда}

\begin{lemma}[Критерий эквивалентности типа Уайтхеда]\label{lem:whitehead}
Пусть $\Phi : \cC_1 \to \cC_2$ — $(\infty, \infty)$-функтор между доступными $(\infty, \infty)$-категориями. Тогда $\Phi$ — $(\infty, \infty)$-эквивалентность тогда и только тогда, когда его усечения $\tau_{\leq n}(\Phi) : \tau_{\leq n}(\cC_1) \to \tau_{\leq n}(\cC_2)$ — $(\infty, n)$-эквивалентности для каждого $n \in \mathbb{N}$.
\end{lemma}

Ссылки: Люри HTT~\cite{LurieHTT} \S 6.5.4 («теорема Уайтхеда» для $(\infty, 1)$-топосов) и \S 5.5.6 (усечённые объекты в $(\infty, 1)$-категориях); Барвик--Шоммер-Прис~\cite{BarwickSchommerPries2021} — $(\infty, n)$-расширение.

\subsection{Стабилизация Бергнер--Люри}

\begin{theorem}[$(\infty, \infty)$-стабилизация]\label{thm:bergner-lurie-stab}
Положим $(\infty, \infty)\text{-}\Cat := \colim_{n < \infty} (\infty, n)\text{-}\Cat$ вдоль канонических вложений и определим $(\infty, \infty + 1)\text{-}\Cat$ как $(\infty, \infty)$-категорию категорий, обогащённых в $(\infty, \infty)\text{-}\Cat$ (один дополнительный шаг обогащения). Тогда каноническое вложение $(\infty, \infty)\text{-}\Cat \hookrightarrow (\infty, \infty + 1)\text{-}\Cat$ — эквивалентность: копредел идемпотентен под обогащением.
\end{theorem}

\begin{proof}[Доказательство (сборка)]
На каждом конечном уровне категории, обогащённые в $(\infty, n)\text{-}\Cat$, суть $(\infty, n+1)$-категории, единственным образом с точностью до эквивалентности по теореме унитарности Барвика--Шоммер-Приса~\cite{BarwickSchommerPries2021} со сравнениями моделей Бергнер--Резка~\cite{BergnerRezk2013}. Обогащение коммутирует с фильтрованным копределом по $n$, так что $(\infty, \infty+1)\text{-}\Cat \simeq \colim_n (\infty, n+1)\text{-}\Cat = (\infty, \infty)\text{-}\Cat$; глобулярная модель Ара--Мальциниотиса~\cite{AraMaltsiniotis2017} реализует ту же идемпотентность для $\omega$-категорий. Это чтение «неограниченной глубины» соответствует конвенции, зафиксированной в (R5); ни одно утверждение статьи не требует альтернативного всюду-необратимого чтения.
\end{proof}

\section{Паранепротиворечивое расширение}\label{app:paraconsistent}

AFN-T в заявленной форме (Теорема~\ref{thm:afnt-alpha}) применяется к классическим Rich-метатеориям. Это приложение распространяет результат на паранепротиворечивые Rich-метатеории \emph{с извлекаемым классическим ядром} --- класс паранепротиворечивых логик, в которых каждая теорема допускает свидетеля с классическим выводом. Расширение \emph{не} покрывает сильно-диалетеические теории (см.\ ниже).

\subsection{Паранепротиворечивые R-S}

\begin{definition}[Паранепротиворечивая R-S с извлекаемым классическим ядром]\label{def:paraconsistent-rs}
Теория $S'$ — \textbf{паранепротиворечивая Rich-метатеория с извлекаемым классическим ядром}, если она удовлетворяет классическим условиям (R1')--(R5'), аналогичным (R1)--(R5), вместе с:
\begin{itemize}
\item \textbf{(Explosion-min)}: $\bot \not\vdash \phi$ для любой формулы $\phi$ (противоречия не влекут всего).
\item \textbf{(Strong-neg)}: существует отрицание $\sim$ такое, что $(p \wedge \sim p) \vdash_\sim \bot$ (конкретный принцип противоречия, хотя и не полный взрыв).
\item \textbf{(Извлекаемость классического ядра)}: наибольший фрагмент $S'^* \subseteq S'$, замкнутый под классическими правилами вывода, сам непротиворечив, непуст и является Rich-метатеорией в $\RS$. Эквивалентно: каждая теорема $S'$ фундаментальной значимости (в частности, определяющая формула любого кандидата в $\LAbs$) допускает свидетеля, выводимого чисто классическим выводом из классических аксиом $S'$.
\end{itemize}
\end{definition}

\subsection{Паранепротиворечивая AFN-T}

\begin{theorem}[Паранепротиворечивая AFN-T]\label{thm:paraconsistent-afnt}
Если $S'$ — паранепротиворечивая Rich-метатеория с извлекаемым классическим ядром в смысле Определения~\ref{def:paraconsistent-rs}, то AFN-T выполняется для $S'$ с тем же формальным содержанием, что и Теорема~\ref{thm:afnt-alpha}.
\end{theorem}

\begin{proof}[Доказательство (схема)]

Пусть $S'^* \subseteq S'$ обозначает \emph{классическое ядро} $S'$: наибольший фрагмент $S'$, замкнутый под классическими правилами вывода, в котором ограничение $(\text{Explosion-min})$ ни разу не привлекается. Эквивалентно, $S'^* = \{\phi \in S' : S' \vdash \phi \text{ только классическими правилами}\}$. Стандартные результаты о звёздной семантике Раутли--Мейера (Прист~\cite{Priest2006}) и о классическом восстановлении для релевантных логик со (Strong-neg) показывают, что $S'^*$ корректно определено, непротиворечиво и само является Rich-метатеорией из $\RS$: оно удовлетворяет (R1)--(R5), поскольку $S'$ удовлетворяет, а ограничение классическим выводом сохраняет р.п.\ аксиоматизацию, существование модели, гёделево кодирование и категорную семантику.

Пусть AFN-T нарушается в $S'$: существует $X$, удовлетворяющий $\LAbs$-паре $(F_{S'}) \wedge (\Pi_{3\text{-max}, S', n})$. Классическая проекция $X$ --- редукт $X^*$ объекта $X$ к $S'^*$, получаемый отбрасыванием паранепротиворечивых выводов, --- даёт мнимого $\LAbs$-свидетеля в $S'^*$.

Проверим, что оба условия сохраняются при классической проекции:
\begin{itemize}
\item $(F_{S'^*})(X^*)$: определяющая формула $\phi_X$ использует только классические связки плюс $\sim$; вычёркивание $\sim$-вхождений сохраняет классическую определимость.
\item $(\Pi_{3\text{-max}, S'^*, n})(X^*)$: каждая $F \in \cF_{S'^*}$ тем более лежит в $\cF_{S'}$; значит, Морита-редукция $F \to X$ существует, и её классическая проекция даёт $F \to X^*$.
\end{itemize}

Итак, $X^*$ — $\LAbs$-свидетель в $S'^* \in \RS$, что противоречит Теореме~\ref{thm:afnt-alpha}.

Следовательно, AFN-T выполняется в $S'$. Полные детали (включая существование классического ядра для любой паранепротиворечивой Rich-метатеории со (Strong-neg) и функториальность классической проекции на $\cS_{S'}$) — в сопутствующей заметке.
\end{proof}

\section{Верификационный леджер Verum}\label{app:verification-ledger}

Каждая теорема, лемма, определение и следствие этой статьи имеет сопутствующее машинно-проверенное формальное доказательство в \textbf{Verum MSFS Corpus}, доступном по адресу \url{https://github.com/gst-st/msfs}. Сопутствующий артефакт снабжён собственным README с полной картой «статья $\to$ корпус», отчётом о доверенной границе, живой аудит-поверхностью (\texttt{verum audit -\!-bundle} выдаёт вердикт одной командой), экспортами в чужие форматы (Coq, Lean~4, Agda, Dedukti, Metamath) и инструкциями воспроизводимости.

\medskip

\noindent\textbf{Доверенная граница.} Верификация MSFS самодостаточна по модулю:
\begin{itemize}
\item ZFC,
\item двух сильно недостижимых кардиналов $\kappa_1 < \kappa_2$,
\item доверенного ядра Verum (\texttt{crates/verum\_kernel/} дерева исходников Verum).
\end{itemize}
Ничто иное не допускается. Граница принуждается ядерным инвариантом
\texttt{verum\_kernel::mechanisation\_roadmap::msfs\_self\_contained()}
и корпусной Verum-теоремой
\texttt{msfs\_self\_containment\_theorem}; оба перепроверяются при каждой сборке и отклоняют любой коммит, вводящий внешнюю зависимость в границах MSFS.

\medskip

\noindent\textbf{Цитаты внешних авторитетов} (Люри HTT 2009, Адамек--Росицкий 1994, Бергнер--Люри 2021, Риль--Верити 2022, Уайтхед, бикатегория дробей Пронк, теорема Ловера о неподвижной точке) представлены в корпусе маркерами \texttt{@framework(name, "citation")} и перечисляются командой \texttt{verum audit -\!-framework-axioms}.


\section*{Благодарности и раскрытия}

Автор благодарит математическое сообщество за фундаментальные результаты Кантора, Рассела, Гёделя, Тарского, Ловера, Ламбека--Скотта, Адамека--Росицкого, Маккаи--Паре, Джейкобса, Пронк, Люри, Риль--Верити, Воеводского, Шрайбера и Клаузена--Шольце, на которых покоится этот синтез. Макс Середа — единственный автор; концептуализация, формальный анализ, доказательства и текст — целиком ответственность автора. Работа выполнена независимо; автор заявляет об отсутствии конкурирующих интересов и внешнего финансирования. Данные не генерировались и не анализировались; полный LaTeX-исходник заархивирован вместе с депозитом препринта.

\begin{thebibliography}{99}

\bibitem{Aczel1988}
P.~Aczel,
\newblock \emph{Non-Well-Founded Sets},
\newblock CSLI Lecture Notes 14, Stanford, 1988.

\bibitem{AdamekRosicky1994}
J.~Ad\'amek and J.~Rosick\'y,
\newblock \emph{Locally Presentable and Accessible Categories},
\newblock London Mathematical Society Lecture Notes 189, Cambridge University Press, 1994.

\bibitem{AraMaltsiniotis2017}
D.~Ara and G.~Maltsiniotis,
\newblock \emph{Homotopy comparison of cubical and globular models of weak $\omega$-categories},
\newblock Preprint, 2017.

\bibitem{AwodeyWarren2009}
S.~Awodey and M.~Warren,
\newblock Homotopy theoretic models of identity types,
\newblock \emph{Mathematical Proceedings of the Cambridge Philosophical Society} 146 (2009), 45--55.

\bibitem{Bakovic2010}
I.~Bakovi\'c,
\newblock \emph{Fibrations of bicategories},
\newblock Preprint, 2010.

\bibitem{Rezk2010}
C.~Rezk,
\newblock A Cartesian presentation of weak $n$-categories,
\newblock \emph{Geometry \& Topology} 14 (2010), no.~1, 521--571; arXiv:0901.3602.

\bibitem{Benabou1967}
J.~B\'enabou,
\newblock Introduction to bicategories,
\newblock in \emph{Reports of the Midwest Category Seminar}, Lecture Notes in Mathematics 47, Springer, 1967, pp.~1--77.

\bibitem{BergnerRezk2013}
J.E.~Bergner and C.~Rezk,
\newblock Comparison of models for $(\infty, n)$-categories, I,
\newblock \emph{Geometry \& Topology} 17 (2013), no.~4, 2163--2202.

\bibitem{BarwickSchommerPries2021}
C.~Barwick and C.~Schommer-Pries,
\newblock On the unicity of the theory of higher categories,
\newblock \emph{Journal of the American Mathematical Society} 34 (2021), 1011--1058; arXiv:1112.0040.

\bibitem{Cantor1891}
G.~Cantor,
\newblock \"Uber eine elementare Frage der Mannigfaltigkeitslehre,
\newblock \emph{Jahresbericht der Deutschen Mathematiker-Vereinigung} 1 (1891), 75--78.

\bibitem{ClausenScholze2019}
D.~Clausen and P.~Scholze,
\newblock \emph{Lectures on Condensed Mathematics},
\newblock Lecture notes, University of Bonn, 2019,
\newblock \url{https://www.math.uni-bonn.de/people/scholze/Condensed.pdf}.

\bibitem{Ernst2015}
M.~Ernst,
\newblock The prospects of unlimited category theory: doing what remains to be done,
\newblock \emph{Review of Symbolic Logic} 8 (2015), no.~2, 306--327.

\bibitem{Feferman1964}
S.~Feferman,
\newblock Systems of predicative analysis,
\newblock \emph{Journal of Symbolic Logic} 29 (1964), no.~1, 1--30.

\bibitem{Shulman2019}
M.A.~Shulman,
\newblock All $(\infty, 1)$-toposes have strict univalent universes,
\newblock preprint, 2019; arXiv:1904.07004.

\bibitem{GagnaHarpazLanari2021}
A.~Gagna, Y.~Harpaz, and E.~Lanari,
\newblock Gray tensor products and lax functors of $(\infty, 2)$-categories,
\newblock \emph{Advances in Mathematics} 391 (2021), 107986; arXiv:2002.06055.

\bibitem{GambinoGarner2008}
N.~Gambino and R.~Garner,
\newblock The identity type weak factorisation system,
\newblock \emph{Theoretical Computer Science} 409 (2008), 94--109.

\bibitem{Godel1931}
K.~G\"odel,
\newblock \"Uber formal unentscheidbare S\"atze der Principia Mathematica und verwandter Systeme I,
\newblock \emph{Monatshefte f\"ur Mathematik und Physik} 38 (1931), 173--198.

\bibitem{Hermida1999}
C.~Hermida,
\newblock Some properties of Fib as a fibred $2$-category,
\newblock \emph{Journal of Pure and Applied Algebra} 134 (1999), 83--109.

\bibitem{Hofmann1995}
M.~Hofmann,
\newblock \emph{Extensional concepts in intensional type theory},
\newblock PhD thesis, University of Edinburgh, 1995.

\bibitem{Hyland1982}
J.M.E.~Hyland,
\newblock The effective topos,
\newblock in \emph{The L.E.J. Brouwer Centenary Symposium}, Studies in Logic and the Foundations of Mathematics 110, North-Holland, 1982, pp.~165--216.

\bibitem{Jacobs1999}
B.~Jacobs,
\newblock \emph{Categorical Logic and Type Theory},
\newblock Studies in Logic 141, Elsevier, 1999.

\bibitem{Bishop1967}
E.~Bishop,
\newblock \emph{Foundations of Constructive Analysis},
\newblock McGraw-Hill, 1967.

\bibitem{BridgesVita2006}
D.S.~Bridges and L.S.~Vita,
\newblock \emph{Techniques of Constructive Analysis},
\newblock Universitext, Springer, 2006.

\bibitem{Buss1986}
S.R.~Buss,
\newblock \emph{Bounded Arithmetic},
\newblock Bibliopolis, Naples, 1986.

\bibitem{EseninVolpin1961}
A.S.~Esenin-Volpin,
\newblock Le programme ultra-intuitionniste des fondements des math\'ematiques,
\newblock in: \emph{Infinitistic Methods}, Pergamon Press / PWN, Warsaw, 1961, pp.~201--223.

\bibitem{EseninVolpin1970}
A.S.~Esenin-Volpin,
\newblock The ultra-intuitionistic criticism and the antitraditional program for foundations of mathematics,
\newblock in: A.~Kino, J.~Myhill, R.E.~Vesley (eds.), \emph{Intuitionism and Proof Theory}, North-Holland, 1970, pp.~3--45.

\bibitem{MarkovNagornyj1988}
A.A.~Markov and N.M.~Nagornyj,
\newblock \emph{The Theory of Algorithms},
\newblock Mathematics and its Applications (Soviet Series), Kluwer, 1988 (English translation of Nauka, Moscow, 1984).

\bibitem{Simpson2009}
S.G.~Simpson,
\newblock \emph{Subsystems of Second Order Arithmetic},
\newblock 2nd ed., Cambridge University Press, 2009.

\bibitem{LambekScott1986}
J.~Lambek and P.J.~Scott,
\newblock \emph{Introduction to Higher-Order Categorical Logic},
\newblock Cambridge Studies in Advanced Mathematics 7, Cambridge University Press, 1986.

\bibitem{KapulkinLumsdaine2021}
K.~Kapulkin and P.L.~Lumsdaine,
\newblock The simplicial model of univalent foundations (after Voevodsky),
\newblock \emph{Journal of the European Mathematical Society} 23 (2021), no.~6, 2071--2126; arXiv:1211.2851.

\bibitem{Lair1981}
C.~Lair,
\newblock Cat\'egories modelables et cat\'egories esquissables,
\newblock \emph{Diagrammes} 6 (1981), L1--L20.

\bibitem{MakkaiPare1989}
M.~Makkai and R.~Par\'e,
\newblock \emph{Accessible Categories: The Foundations of Categorical Model Theory},
\newblock Contemporary Mathematics 104, American Mathematical Society, 1989.

\bibitem{Kelly1982}
G.M.~Kelly,
\newblock \emph{Basic Concepts of Enriched Category Theory},
\newblock London Mathematical Society Lecture Notes 64, Cambridge University Press, 1982.

\bibitem{KellyStreet1974}
G.M.~Kelly and R.~Street,
\newblock Review of the elements of $2$-categories,
\newblock Lecture Notes in Mathematics 420, Springer, 1974, pp.~75--103.

\bibitem{Lawvere1969}
F.W.~Lawvere,
\newblock Adjointness in foundations,
\newblock \emph{Dialectica} 23 (1969), 281--296.

\bibitem{LurieHTT}
J.~Lurie,
\newblock \emph{Higher Topos Theory},
\newblock Annals of Mathematics Studies 170, Princeton University Press, 2009.

\bibitem{LurieHA}
J.~Lurie,
\newblock \emph{Higher Algebra},
\newblock Preprint, 2017.

\bibitem{Girard2001}
J.-Y.~Girard,
\newblock Locus solum: from the rules of logic to the logic of rules,
\newblock \emph{Mathematical Structures in Computer Science} 11 (2001), no.~3, 301--506.

\bibitem{Priest2006}
G.~Priest,
\newblock \emph{In Contradiction: A Study of the Transconsistent},
\newblock 2nd edition, Oxford University Press, 2006.

\bibitem{Pronk1996}
D.~Pronk,
\newblock Etendues and stacks as bicategories of fractions,
\newblock \emph{Compositio Mathematica} 102 (1996), no.~3, 243--303.

\bibitem{Rathjen1999}
M.~Rathjen,
\newblock The realm of ordinal analysis,
\newblock in \emph{Sets and Proofs} (Leeds 1997), London Mathematical Society Lecture Notes 258, Cambridge University Press, 1999, pp.~219--279.

\bibitem{RiehlVerity2022}
E.~Riehl and D.~Verity,
\newblock \emph{Elements of $\infty$-Category Theory},
\newblock Cambridge Studies in Advanced Mathematics 194, Cambridge University Press, 2022.

\bibitem{Russell1903}
B.~Russell,
\newblock \emph{The Principles of Mathematics},
\newblock Cambridge University Press, 1903.

\bibitem{Schreiber2013}
U.~Schreiber,
\newblock \emph{Differential cohomology in a cohesive $\infty$-topos},
\newblock arXiv:1310.7930, 2013.

\bibitem{Shulman2008}
M.A.~Shulman,
\newblock $2$-toposes,
\newblock \emph{Journal of Pure and Applied Algebra} 210 (2008), 273--287.

\bibitem{Streicher1991}
T.~Streicher,
\newblock \emph{Semantics of type theory},
\newblock Birkh\"auser Progress in Theoretical Computer Science, 1991.

\bibitem{Tarski1936}
A.~Tarski,
\newblock Der Wahrheitsbegriff in den formalisierten Sprachen,
\newblock \emph{Studia Philosophica} 1 (1936), 261--405.

\bibitem{HoTTBook2013}
The Univalent Foundations Program,
\newblock \emph{Homotopy Type Theory: Univalent Foundations of Mathematics},
\newblock Institute for Advanced Study, Princeton, 2013,
\newblock \url{https://homotopytypetheory.org/book/}.

\bibitem{VoevodskyFoundations2010}
V.~Voevodsky,
\newblock \emph{Foundations} Coq library,
\newblock software artefact, 2010--2015,
\newblock \url{https://github.com/vladimirias/Foundations}.

\bibitem{Werner1994}
B.~Werner,
\newblock \emph{Une Th\'eorie des Constructions Inductives},
\newblock PhD thesis, Universit\'e Paris 7, 1994.

\bibitem{Yanofsky2003}
N.S.~Yanofsky,
\newblock A universal approach to self-referential paradoxes, incompleteness and fixed points,
\newblock \emph{Bulletin of Symbolic Logic} 9 (2003), 362--386.

\end{thebibliography}

\end{document}
